% !TeX root = master.tex
% !BIB program = biber
\documentclass[%
english, titlepage=false % or `german`
]{./tum-templates/book/tumbook}

\pgfkeys{/TUM/book/.cd, thesis=student, coverpage=false}

% Keep mathematics in chapter headings and table-of-contents chapter entries
% consistent with their surrounding bold text.
\addtokomafont{chapter}{\boldmath}
\addtokomafont{chapterentry}{\boldmath}

\makeatletter %redefine some labels from the TUM template
\provideName{\@tum@examiner@}{Supervisor}{Themensteller} % or `Themenstellerin`
\provideName{\@tum@supervisor@}{Advisors}{Betreuer} % or `Advisor` / `Betreuerin`
\makeatother

\usepackage{booktabs}% for more beautiful tables
\usepackage{cleveref}% intelligent references
\usepackage{graphicx}% for including images
\usepackage{pgfplots}% for plotting graphs
\pgfplotsset{compat=1.17}% compatible with TeX Live 2021 and newer
\usepgfplotslibrary{groupplots}

\usepackage{diffcoeff}
\usepackage{amsmath}
\usepackage{amsthm}
\usepackage{framed}
\usepackage{amssymb}
\usepackage{mathtools}
\usepackage{tikz}
\usepackage{subcaption}
\usepackage{enumitem}
\usepackage[T1]{fontenc}
\usepackage{lmodern}
% Use the traditional LaTeX serif face for body text; headings retain the
% template's sans-serif styling.
\renewcommand{\familydefault}{\rmdefault}
\usepackage{microtype}
\setlength{\parindent}{0pt}

\newtheorem{corollary}{Corollary}
\newtheorem{theorem}{Theorem}
\newtheorem{lemma}{Lemma}
\newtheorem{definition}{Definition}
\newtheorem{remark}{Remark}
\newtheorem{assumption}{Assumption}
\newtheorem{property}{Property}
\newtheorem{result}{Result}
\newtheorem{construction}{Construction}
\newtheorem{conjecture}{Conjecture}


\DeclareMathOperator*{\argmax}{arg\,max}
\DeclareMathOperator*{\argmin}{arg\,min}

\pgfmathdeclarefunction{Uk}{2}{%
  \pgfmathparse{
    1
    + (#1-1)
      / (
        #2*#1
        * (
          1
          + sqrt(
            1
            + (
              4*(#1-1)
              * ((#1-1) + #2*(#1^2-2*#1+2))
            )
            / (#2*#1^2)
          )
        )
      )
  }%
}

%Literatur
\usepackage[%
    backend=biber,
    sortcites, % sort automatically
    sorting=nty, % sort order
    citestyle=alphabetic, %
    bibstyle=alphabetic, %
    hyperref=true, % provide clickable links
    maxbibnames=3, % shorten author list for more than 3 names
    maxcitenames=3, % use at most 3 names for key
    url=false, % do not print URLs
    doi=false, % do not print DOIs
    giveninits=true,
    ]%
{biblatex}
\addbibresource{literature.bib}

% automatische Anführungszeichen
\usepackage[autostyle=true]{csquotes}


%Kürzel
\definecolor{jobhighlight}{HTML}{9ABCE4}

\newcommand{\Cmax}{C_{\max}}
\newcommand{\Cmin}{C_{\min}}
\newcommand{\WSPT}{\mathrm{WSPT}}
\newcommand{\Cminmax}{L^{\min}}

\newcommand{\Problem}{P\!\parallel\!\sum_j p_j C^{\alpha}_j}
\newcommand{\LamI}{\lambda_{\alpha}(I)}
\newcommand{\Lam}{\lambda(\alpha)}
\newcommand{\Cost}{\gamma_{\alpha}(\sigma(I))}
\newcommand{\ComplAlpha}{C_{j}^{\alpha}(\sigma(I))}
\newcommand{\OptCost}{\gamma_{\alpha}^{*}(I)}
\newcommand{\WorstCost}{\gamma_{\alpha}^{W}(I)}
\newcommand{\Sstar}{\sigma^{*}(I)}
\newcommand{\SW}{\sigma^{W}(I)}
\newcommand{\machLoad}{L^{i}(\sigma(I))}
\newcommand{\machLoadMax}{L^{\max}(\sigma(I))}
\newcommand{\machLoadMin}{L^{\min}(\sigma(I))}

% Shared helpers for schedule diagrams.
\newcommand{\machinerow}[2]{%
    \node[machine, anchor=east] at (-0.3,-#1+0.5) {$M_{#2}$};
}

\newcommand{\dotsrow}[1]{%
    \foreach \dy in {-0.18,0,0.18} { % chktex 1
        \fill (-0.6,{-#1+0.5+\dy}) circle[radius=0.025];
        \fill (2,{-#1+0.5+\dy}) circle[radius=0.025];
    }
}

\newcommand{\jobblock}[5][white]{%
    \draw[job, fill=#1] (#3,-#2) rectangle ++(#4,1);
    \node at ({#3 + #4/2},{-#2 + 0.5}) {$#5$}; % chktex 1
}
\newcommand{\schedulebeforeunification}{
\begin{tikzpicture}[
    x=0.7cm,
    y=0.7cm,
    job/.style={
        draw,
        thick,
        fill=white,
        font=\small
    },
    machine/.style={font=\small}
]

% Displayed machine rows
\machinerow{1}{1}
\dotsrow{2}
\machinerow{3}{i^{w}}
\dotsrow{4}
\machinerow{5}{m}

% Edit the first schedule here.
\jobblock[jobhighlight]{1}{0}{1}{}
\jobblock{1}{1}{3}{}
\jobblock{1}{4}{2}{}


\jobblock{3}{0}{3}{}
\jobblock{3}{3}{2}{}
\jobblock{3}{5}{4}{k}

\jobblock{5}{0}{2}{}
\jobblock[jobhighlight]{5}{2}{2}{}
\jobblock{5}{4}{2}{}

\end{tikzpicture}
}
\newcommand{\schedulecaseone}{
\begin{tikzpicture}[
    x=0.7cm,
    y=0.7cm,
    job/.style={
        draw,
        thick,
        fill=white,
        font=\small
    },
    machine/.style={font=\small}
]

% Displayed machine rows
\machinerow{1}{1}
\dotsrow{2}
\machinerow{3}{u}
\machinerow{4}{u+1}
\dotsrow{5}
\machinerow{6}{m}

% Edit the second schedule here.

\jobblock{3}{0}{4}{p}
\jobblock[jobhighlight]{3}{4}{1}{s}

\jobblock{1}{0}{4}{p}
\jobblock[jobhighlight]{1}{4}{1}{s}

\jobblock[jobhighlight]{4}{0}{5}{s}

\jobblock[jobhighlight]{6}{0}{5}{s}


\end{tikzpicture}
}
\newcommand{\schedulecasetwo}{
\begin{tikzpicture}[
    x=0.7cm,
    y=0.7cm,
    job/.style={
        draw,
        thick,
        fill=white,
        font=\small
    },
    machine/.style={font=\small}
]

% Displayed machine rows
\machinerow{1}{1}
\dotsrow{2}
\machinerow{3}{u}
\machinerow{4}{u+1}
\dotsrow{5}
\machinerow{6}{m}

% Edit the second schedule here.

\jobblock{3}{0}{6}{p}


\jobblock{1}{0}{6}{p}

\jobblock[jobhighlight]{4}{0}{4}{s}

\jobblock[jobhighlight]{6}{0}{4}{s}


\end{tikzpicture}
}
\newcommand{\scheduleappendixm}{
\begin{tikzpicture}[
    x=0.7cm,
    y=0.7cm,
    job/.style={
        draw,
        thick,
        fill=white,
        font=\small
    },
    machine/.style={font=\small}
]

% Displayed machine rows
\machinerow{1}{1}
\machinerow{2}{2}

% Edit the first schedule here.
\jobblock{1}{0}{2}{1}
\jobblock{2}{0}{2}{1}
\jobblock{1}{2}{5}{x^{*}_{\alpha}}


\end{tikzpicture}
}
\newcommand{\scheduleappendixmm}{
\begin{tikzpicture}[
    x=0.7cm,
    y=0.7cm,
    job/.style={
        draw,
        thick,
        fill=white,
        font=\small
    },
    machine/.style={font=\small}
]

% Displayed machine rows
\machinerow{1}{1}
\machinerow{2}{2}

% Edit the first schedule here.
\jobblock{1}{0}{2}{1}
\jobblock{1}{2}{2}{1}
\jobblock{2}{0}{5}{x^*_{\alpha}}

\end{tikzpicture}
}

\newcommand{\scheduleworstcase}{
\begin{tikzpicture}[
    x=0.7cm,
    y=0.7cm,
    job/.style={
        draw,
        thick,
        fill=white,
        font=\small
    },
    machine/.style={font=\small}
]

% Displayed machine rows
\machinerow{1}{1}
\dotsrow{2}
\machinerow{3}{u}
\machinerow{4}{u+1}
\dotsrow{5}
\machinerow{6}{m}

% Edit the second schedule here.

\jobblock[jobhighlight]{1}{0}{3}{s}

\jobblock{1}{3}{5}{p}

\jobblock[jobhighlight]{3}{0}{3}{s}

\jobblock{3}{3}{5}{p}

\jobblock[jobhighlight]{4}{0}{3}{s}
\jobblock[jobhighlight]{6}{0}{3}{s}

\end{tikzpicture}
}
\newcommand{\scheduleafterunification}{
\begin{tikzpicture}[
    x=0.7cm,
    y=0.7cm,
    job/.style={
        draw,
        thick,
        fill=white,
        font=\small
    },
    machine/.style={font=\small}
]

% Displayed machine rows
\machinerow{1}{1}
\dotsrow{2}
\machinerow{3}{i^k}
\dotsrow{4}
\machinerow{5}{m}

% Edit the second schedule here.
\jobblock{1}{0}{2}{}
\jobblock{1}{2}{3}{}
\jobblock{1}{5}{1}{}

\jobblock[jobhighlight]{3}{0}{2}{}
\jobblock{3}{2}{4}{k}
\jobblock[jobhighlight]{3}{6}{1}{}

\jobblock{5}{0}{2}{}
\jobblock{5}{2}{5}{}


\end{tikzpicture}
}

\newcommand{\scheduleclosetimegaps}{
\begin{tikzpicture}[
    x=0.7cm,
    y=0.7cm,
    job/.style={
        draw,
        thick,
        fill=white,
        font=\small
    },
    machine/.style={font=\small}
]
% Displayed machine rows
\machinerow{1}{1}
\dotsrow{2}
\machinerow{3}{i^{w}}
\dotsrow{4}
\machinerow{5}{m}

% Edit the third schedule here.

\jobblock{1}{1}{3}{}
\jobblock{1}{4}{2}{}


\jobblock{3}{0}{3}{}
\jobblock{3}{3}{2}{}
\jobblock{3}{5}{4}{k}
\jobblock[jobhighlight]{3}{9}{1}{}
\jobblock[jobhighlight]{3}{10}{2}{}

\jobblock{5}{0}{2}{}

\jobblock{5}{4}{2}{}

\end{tikzpicture}
}

\newcommand{\schedulereschedulelastjobs}{
\begin{tikzpicture}[
    x=0.7cm,
    y=0.7cm,
    job/.style={
        draw,
        thick,
        fill=white,
        font=\small
    },
    machine/.style={font=\small}
]

% Displayed machine rows
\machinerow{1}{1}
\dotsrow{2}
\machinerow{3}{i^k}
\dotsrow{4}
\machinerow{5}{m}

% Edit the fourth schedule here.
\jobblock{1}{0}{2}{}
\jobblock{1}{2}{3}{}
\jobblock{1}{5}{1}{}


\jobblock[jobhighlight]{3}{0}{7}{j'}


\jobblock{5}{0}{2}{}
\jobblock{5}{2}{5}{}

\end{tikzpicture}
}

\newcommand{\schedulelistschedule}{
\begin{tikzpicture}[
    x=0.7cm,
    y=0.7cm,
    job/.style={
        draw,
        thick,
        fill=white,
        font=\small
    },
    machine/.style={font=\small}
]

% Displayed machine rows
\machinerow{1}{1}
\dotsrow{2}
\machinerow{3}{i^{w}}
\dotsrow{4}
\machinerow{5}{m}

% Edit the fifth schedule here.
\jobblock{1}{0}{3}{\text{moved}}
\jobblock{1}{3}{2}{\text{moved}}


\jobblock{3}{0}{3}{}
\jobblock{3}{3}{2}{}
\jobblock[jobhighlight]{3}{5}{7}{j'}


\jobblock{5}{0}{2}{}
\jobblock{5}{2}{2}{\text{moved}}

\end{tikzpicture}
}

\newcommand{\schedulefinalworst}{
\begin{tikzpicture}[
    x=0.7cm,
    y=0.7cm,
    job/.style={
        draw,
        thick,
        fill=white,
        font=\small
    },
    machine/.style={font=\small}
]

\machinerow{1}{1}
\dotsrow{2}
\machinerow{3}{i^k}
\dotsrow{4}
\machinerow{5}{m}

% Edit the sixth schedule here.
\jobblock{1}{0}{2}{}
\jobblock{1}{2}{3}{}
\jobblock{1}{5}{1}{}


\jobblock[jobhighlight]{3}{0}{7}{j'}


\jobblock{5}{0}{2}{}
\jobblock{5}{2}{5}{}

\end{tikzpicture}
}

\newcommand{\scheduleseventh}{
\begin{tikzpicture}[
    x=0.7cm,
    y=0.7cm,
    job/.style={
        draw,
        thick,
        fill=white,
        font=\small
    },
    machine/.style={font=\small}
]

\machinerow{1}{1}
\dotsrow{2}
\machinerow{3}{i^{w}}
\dotsrow{4}
\machinerow{5}{m}

% Edit the seventh schedule here.
\jobblock{1}{0}{3}{}
\jobblock{1}{3}{2}{}


\jobblock{3}{0}{3}{}
\jobblock{3}{3}{2}{}



\jobblock{5}{0}{2}{}
\jobblock{5}{2}{2}{}
\jobblock[jobhighlight]{5}{4}{7}{j'}
\end{tikzpicture}
}

\newcommand{\scheduleeighth}{
\begin{tikzpicture}[
    x=0.7cm,
    y=0.7cm,
    job/.style={
        draw,
        thick,
        fill=white,
        font=\small
    },
    machine/.style={font=\small}
]

\machinerow{1}{1}
\dotsrow{2}
\machinerow{3}{i^k}
\dotsrow{4}
\machinerow{5}{m}

% Edit the eighth schedule here.
\jobblock{1}{0}{2}{}
\jobblock{1}{2}{3}{}
\jobblock{1}{5}{1}{}

\jobblock[jobhighlight]{3}{0}{7}{j'}

\jobblock{5}{0}{2}{}
\jobblock{5}{2}{5}{}

\end{tikzpicture}
}
\newcommand{\schedulemtwoone}{
\begin{tikzpicture}[
    x=0.7cm,
    y=0.7cm,
    job/.style={
        draw,
        thick,
        fill=white,
        font=\small
    },
    machine/.style={font=\small}
]

% Displayed machine rows
\machinerow{1}{1}
\machinerow{2}{2}



% Edit the second schedule here.

\jobblock{1}{0}{1.5}{p_1}
\jobblock{1}{1.5}{2.5}{p_3}




\jobblock{2}{0}{3}{p_2}



\end{tikzpicture}
}

\newcommand{\schedulemtwotwo}{
\begin{tikzpicture}[
    x=0.7cm,
    y=0.7cm,
    job/.style={
        draw,
        thick,
        fill=white,
        font=\small
    },
    machine/.style={font=\small}
]

% Displayed machine rows
\machinerow{1}{1}
\machinerow{2}{2}



% Edit the second schedule here.

\jobblock{1}{0}{1.5}{p_1}
\jobblock{1}{1.5}{3}{p_2}



\jobblock{2}{0}{2.5}{p_3}



\end{tikzpicture}
}

\newcommand{\schedulethreemaxone}{
\begin{tikzpicture}[
    x=0.7cm,
    y=0.7cm,
    job/.style={
        draw,
        thick,
        fill=white,
        font=\small
    },
    machine/.style={font=\small}
]

% Displayed machine rows
\machinerow{1}{1}
\machinerow{2}{2}
\machinerow{3}{3}


% Edit the second schedule here.

\jobblock{1}{0}{4}{t_{\alpha}}
\jobblock{1}{4}{2}{\frac{t_{\alpha}}{2}}



\jobblock{2}{0}{4}{t_{\alpha}}
\jobblock{2}{4}{2}{\frac{t_{\alpha}}{2}}

\jobblock{3}{0}{10}{1}
\end{tikzpicture}
}

\newcommand{\schedulethreemaxtwo}{
\begin{tikzpicture}[
    x=0.7cm,
    y=0.7cm,
    job/.style={
        draw,
        thick,
        fill=white,
        font=\small
    },
    machine/.style={font=\small}
]

% Displayed machine rows
\machinerow{1}{1}
\machinerow{2}{2}
\machinerow{3}{3}


% Edit the second schedule here.
\jobblock{1}{0}{4}{t_{\alpha}}



\jobblock{2}{0}{4}{t_{\alpha}}

\jobblock{3}{0}{2}{\frac{t_{\alpha}}{2}}
\jobblock{3}{2}{2}{\frac{t_{\alpha}}{2}}

\jobblock{3}{4}{6}{1}

\end{tikzpicture}
}
\newcommand{\schedulethreeone}{
\begin{tikzpicture}[
    x=0.7cm,
    y=0.7cm,
    job/.style={
        draw,
        thick,
        fill=white,
        font=\small
    },
    machine/.style={font=\small}
]

% Displayed machine rows
\machinerow{1}{1}
\machinerow{2}{2}
\machinerow{3}{3}


% Edit the second schedule here.

\jobblock{1}{0}{2}{p_1}
\jobblock{1}{2}{2.5}{p_2}
\jobblock{1}{4.5}{2}{p_3}



\jobblock{2}{0}{2}{p_4}
\jobblock{2}{2}{2}{p_5}
\jobblock{2}{4}{1.5}{p_6}

\jobblock{3}{0}{6}{p_7}
\end{tikzpicture}
}

\newcommand{\schedulethreetwo}{
\begin{tikzpicture}[
    x=0.7cm,
    y=0.7cm,
    job/.style={
        draw,
        thick,
        fill=white,
        font=\small
    },
    machine/.style={font=\small}
]

% Displayed machine rows
\machinerow{1}{1}
\machinerow{2}{2}
\machinerow{3}{3}


% Edit the second schedule here.

\jobblock{1}{0}{2}{p_1}
\jobblock{1}{2}{2}{p_4}



\jobblock{2}{0}{2.5}{p_2}
\jobblock{2}{2.5}{2}{p_5}


\jobblock{3}{0}{2}{p_3}
\jobblock{3}{2}{1.5}{p_6}
\jobblock{3}{3.5}{6}{p_7}

\end{tikzpicture}
}

\title{Approximation Ratio of the WSPT Rule for Total Weighted \texorpdfstring{$\alpha$}{alpha}-Completion Time Scheduling}
\subtitle{}

\author{Jake Wefers}

\degree{Master of Science}% or `Bachelor of Science`
\dateSubmitted{23 September 2026}% preferably use some universally recognized date format


\examiner{Prof.\@ Dr.\@ Andreas S. Schulz}% `Themensteller`
\supervisor{M.\@ Sc.\@ Maximilian von Aspern}% `Betreuer`


\begin{document}
\raggedbottom

\frontmatter
\maketitle


\section*{Abstract}
\addcontentsline{toc}{chapter}{Abstract/Zusammenfassung}
The Weighted Shortest Processing Time ($\mathrm{WSPT}$) rule schedules jobs on 
identical parallel machines in non-increasing order of their weight-to-processing-time
ratios $w_j/p_j$. Its tight approximation ratio for minimizing the sum of weighted completion
times is $\frac{1}{2}(1+\sqrt{2})$, as established by Kawaguchi and Kyan \cite{KawaguchiKyan:1988}.
We investigate how this guarantee changes when completion times are replaced
by $\alpha$-completion times. We prove that
the approximation ratio is unbounded for $\alpha=0$ and equals
$1+\frac{1}{2(\alpha+\sqrt{2\alpha})}$ for
$\alpha\in\Big[\frac{1}{2},1\Big]$. For the unresolved range
$\alpha\in \Big(0,\frac{1}{2}\Big)$, we derive structural reductions for
worst-case instances, establish general and machine-dependent upper
bounds, and construct explicit lower bounds. We additionally determine
exact ratios for the specific case of two and three machines and formulate a mixed-integer
nonlinear program to investigate the remaining gap. Since deterministic
$\alpha$-completion-time bounds transfer to stochastic scheduling, these
results may also sharpen performance guarantees for the Weighted
Shortest Expected Processing Time ($\mathrm{WSEPT}$) rule.


\section*{Zusammenfassung}
Die Weighted-Shortest-Processing-Time-Regel ($\mathrm{WSPT}$) plant Aufträge auf
identischen parallelen Maschinen in nicht aufsteigender Reihenfolge ihrer
Gewicht-Prozesszeit-Verhältnisse $w_j/p_j$. Ihr scharfer Approximationsfaktor
für die Minimierung der Summe der gewichteten Fertigstellungszeiten beträgt
$\frac{1}{2}(1+\sqrt{2})$, wie Kawaguchi und Kyan \cite{KawaguchiKyan:1988} gezeigt haben.
Wir untersuchen, wie sich diese Garantie verändert, wenn Fertigstellungszeiten
durch $\alpha$-Fertigstellungszeiten ersetzt werden. Wir beweisen, dass
der Approximationsfaktor für $\alpha=0$ unbeschränkt ist und für
$\alpha\in[\frac{1}{2},1]$ den Wert
$1+\frac{1}{2(\alpha+\sqrt{2\alpha})}$ annimmt. Für den bislang ungeklärten Bereich
$\alpha\in(0,\frac{1}{2})$ leiten wir strukturelle Reduktionen für
Worst-Case-Instanzen her, beweisen allgemeine sowie von der Maschinenanzahl abhängige obere
Schranken und konstruieren explizite untere Schranken. Darüber hinaus bestimmen wir
exakte Approximationsfaktoren für die Spezialfälle mit zwei und drei Maschinen und formulieren ein
Optimierungsmodell, um die verbleibende Lücke zu untersuchen. Da sich deterministische
Schranken für $\alpha$-Fertigstellungszeiten auf die stochastische Ablaufplanung übertragen lassen, können diese
Ergebnisse auch die Gütegarantien für die Weighted-Shortest-Expected-Processing-Time-Regel
($\mathrm{WSEPT}$) verbessern.


\cleardoublepage
\tableofcontents

\chapter*{Notation}
\addcontentsline{toc}{chapter}{Notation}

\begin{tabular}{@{}p{0.22\textwidth}p{0.68\textwidth}@{}}
\toprule
Symbol & Meaning \\
\midrule
$I$ & Scheduling instance. \\[2mm]
$n$ & Number of jobs. \\[2mm]
$m$ & Number of identical parallel machines. \\[2mm]
$M_i$ & Machine $i$. \\[2mm]
$J$ & Set of jobs. Note that we will almost always write $j \in I$ instead of $j \in J$ as is common for scheduling problems. \\[2mm]
$\mathcal{M}$ & Set of machines $\{M_i:i\in[m]\}$. \\[2mm]
$p_j$ & Processing time of job $j$. \\[2mm]
$\sigma(I)$ & Schedule of instance $I$. \\[2mm]
$J_i(\sigma(I))$ & Set of jobs scheduled on machine $i$ in schedule $\sigma(I)$.\\[2mm]
$w_j$ & Weight of job $j$. \\[2mm]
$S_j(\sigma(I))$ & Start time of job $j$ in schedule $\sigma(I)$. \\[2mm]
$\alpha$ & Completion-time parameter. \\[2mm]
$P\!\parallel\!\sum_j w_j C^{\alpha}_j$ & Scheduling problem on identical parallel machines, minimizing total weighted $\alpha$-completion time. \\[2mm]
$\sigma^{\mathrm{WSPT}}(I)$ & A $\mathrm{WSPT}$ schedule of instance $I$. \\[2mm]
$C^{\alpha}_{j}(\sigma(I))$ & $\alpha$-completion time of job $j$ in schedule $\sigma(I)$.\\[2mm]
$\gamma_{\alpha}(\sigma(I))$ & Objective value of schedule $\sigma(I)$ for parameter value $\alpha$. \\[2mm]
$\sigma^{W}(I)$ & A cost-maximizing list schedule of instance $I$. \\[2mm]
$\sigma^{*}(I)$ & A cost-minimizing schedule of instance $I$. \\[2mm]
$\gamma^{*}_{\alpha}(I)$ & Minimum objective value out of all list schedules of instance $I$ for parameter value $\alpha$. \\[2mm]
$\gamma^{W}_{\alpha}(I)$ & Maximum objective value over all list schedules of instance $I$ for parameter value $\alpha$. \\[2mm]
$\lambda_{\alpha}(I)$ & Instance-specific performance ratio for parameter value $\alpha$.\\[2mm]
$\lambda_{\alpha}$ & Approximation ratio for parameter value $\alpha$. \\ [2mm]
$L^{i}(\sigma(I))$ & Load of machine $i$ in schedule $\sigma(I)$ of instance $I$. \\[2mm]
$L^{\max}(\sigma(I))$ & Maximal machine load of schedule $\sigma(I)$.\\[2mm]
$L^{\min}(\sigma(I))$ & Minimal machine load of schedule $\sigma(I)$.\\
\bottomrule
\end{tabular}

\mainmatter{}
\chapter{Introduction}\label{chap:introduction}


\section{Background and Motivation}

Scheduling problems arise whenever a collection of jobs must be
assigned to limited resources and arranged over time. In addition to
deciding which resource processes each job, one must determine the
order and start times of the jobs while respecting resource
availability. In practical applications, a suitable schedule is needed to
reduce waiting times, meet deadlines, balance workloads, and use
available resources efficiently. Examples include assigning production
orders to machines, patients to treatment rooms, computational tasks to
processors, and vehicles to charging stations.

The mathematical theory of scheduling provides a formal framework for
these allocation problems. A scheduling instance \(I\) has a job set
\[
    J\coloneqq\{j_{1},\ldots,j_{n}\}
\]
and a set of machines
\[
    \mathcal{M}\coloneqq\{M_i:i\in[m]\}.
\]
We refer to \(M_i\) simply as machine \(i\) when no ambiguity can arise.
Every job \(j\in J\) has a processing time \(p_j>0\) and must be assigned
to a machine. A schedule specifies both the assignment of jobs to
machines and the order in which the jobs assigned to the same machine
are processed. Depending on the application, further information such
as job weights, release dates, due dates, precedence constraints, or
machine-dependent processing times may be included.

The standard form, 
given by \cite{GrahamEtAl1979}, defines a scheduling problem in three parts:
\[
    \mathrm{A}\mid\mathrm{B}\mid\mathrm{C}.
\]
\(\mathrm{A}\) specifies the machine environment. The
second field \(\mathrm{B}\) specifies additional restrictions, and the
third field \(\mathrm{C}\) specifies the objective.


The notation \(1\) in the first field denotes a single machine, while
\(P\) denotes identical parallel machines and $Pm$ restricts this to $m$ identical 
parallel machines. On identical parallel
machines, every job has the same processing time regardless of the
machine it is assigned to. The symbols \(Q\) and \(R\) describe
uniform and unrelated parallel machines, respectively. In these
environments, a job's processing time may depend on the selected
machine.

An empty second field indicates that no additional restrictions are
imposed and is written as $\mathrm{A}\mid\mid\mathrm{C}$. In particular, all jobs are available at time zero and are
processed nonpreemptively. Possible restrictions in the second field
include release dates, precedence constraints between jobs, and preemption, 
which allows a job's processing to be interrupted and resumed later.
One may also distinguish between models that permit idle time and models in which machines 
must process jobs continuously whenever they are available.

\(\mathrm{C}\) determines what should be minimized. Common
objectives include the makespan
\[
    C_{\max}:=\max_{j\in J}C_j,
\]
which is the earliest point in time at which all jobs have been fully processed,
the total completion time
\[
    \sum_{j\in J}C_j,
\]
and the weighted total completion time
\[
    \sum_{j\in J}w_jC_j.
\]
In all three examples, $C_j$ denotes the completion time of job $j$ which is the 
point in time at which $j$ has been fully processed. 



\section{Problem Statement}

In the context of parallel machine scheduling, one particularly important problem is the minimization 
of total weighted completion time on identical parallel machines, denoted as $P \!\parallel\!\sum w_j C_j$. 
This problem involves scheduling $n$ jobs, each with a processing time $p_j$ and a weight $w_j$, on
identical parallel machines to minimize the sum of weighted completion times $\sum w_jC_j$, where $C_j$ 
is the completion time of job $j$.

A well-known heuristic for this problem is the \emph{Weighted-Shortest-Processing-Time} ($\mathrm{WSPT}$) rule, 
which orders jobs by 
non-increasing ratio $w_j / p_j$. Kawaguchi and Kyan~\cite{KawaguchiKyan:1988} proved that 
$\mathrm{WSPT}$ has a worst-case performance ratio of $\frac{1}{2}(1 + \sqrt{2}) \approx 1.207$, meaning the $\mathrm{WSPT}$ schedule 
is at most this factor worse than the optimal schedule. Schwiegelshohn~\cite{Schwiegelshohn:2011} later gave 
an alternative proof based on a sequence of structural reductions that progressively simplify the instances 
that must be considered to determine the approximation ratio.

This thesis extends that line of analysis to weighted $\alpha$-completion time.
To this end, consider the problem
\[
    P\mid\mid\sum_{j\in J}w_jC_j^\alpha \qquad \text{for} \qquad \alpha \in [0,1].
\]

For some instance and schedule, let \(S_j\) denote the starting time of job \(j\). For a parameter
\(\alpha\in [0,1]\), the \(\alpha\)-completion time of job \(j\) is defined by
\[
    C_j^\alpha:=S_j+\alpha p_j.
\]
It is the point in time at which an \(\alpha\)-fraction of job \(j\) has been completed. 
In particular,
\[
    C_j^1=S_j+p_j=C_j
\]
is the commonly known completion time.

The concept of \(\alpha\)-completion is particularly useful in the analysis
of stochastic scheduling problems. In such problems, the processing
time \(p_j\) of job \(j\) is unknown before scheduling and is therefore modeled as a random variable. 
Before the job finishes, only the distribution of \(p_j\) is known, while its actual
processing time becomes known only upon completion. For every possible value that $p_j$ might attain, the value
\[
    C_j^\alpha=S_j+\alpha p_j
\]
identifies the point at which an \(\alpha\)-fraction of the 
processing requirement of $j$ has been completed. Thus,
\(\alpha\)-completion times provide a connection between the stochastic
scheduling problem and its corresponding deterministic problem.
 
The main motivation for this thesis is the paper
\emph{Generalizing the Kawaguchi--Kyan Bound to Stochastic Parallel
Machine Scheduling} by Jäger and Skutella
\cite{JaegerSkutella2018}. The authors study the stochastic parallel
machine problem
\[
    P\mid p_j\sim\mathrm{stoch}\mid
    \mathbb{E}\!\left[\sum_{j\in J}w_jC_j\right]
\]
and analyze the Weighted-Shortest-Expected-Processing-Time ($\mathrm{WSEPT}$) rule.
In this stochastic model, processing times are random, and scheduling decisions must be made before their values are known.

Jäger and Skutella use deterministic weighted sums of
\(\alpha\)-completion times as an auxiliary objective for analyzing the stochastic
problem. Their results show that performance guarantees for the
deterministic problem
\[
    P\mid\mid\sum_{j\in J}w_jC_j^\alpha
\]
can be transferred to performance guarantees for $\mathrm{WSEPT}$ in the
stochastic setting. Thus, although the problem considered in this
thesis is deterministic and treats \(\alpha\) as a fixed parameter,
the resulting approximation bounds are motivated by their possible use
in the analysis of scheduling under uncertain processing times.


An improvement of the bound on the approximation ratio for the deterministic
\(\alpha\)-completion time problem has consequences beyond the deterministic
setting considered here. Suppose that the deterministic $\mathrm{WSPT}$ rule has
approximation ratio
\[
   1+\beta \qquad \text{for} \qquad P\mid\mid\sum_{j\in J}w_jC_j^\alpha.
\]
Jäger and Skutella show that this implies an approximation ratio of
\[
    1+\beta(1+\Delta)
\]
for $\mathrm{WSEPT}$ with respect to the expected weighted sum of
\(\alpha\)-completion times, where
\[
    \Delta
    \ge
    \frac{\operatorname{Var}(p_j)}
         {\mathbb{E}[p_j]^2} \qquad \forall j\in J.
\]

They further obtain the bound
\[
    1+\beta\max\{1,\alpha(1+\Delta)\}
\]
for the approximation ratio of $\mathrm{WSEPT}$ with respect to the expected weighted sum of 
completion times. Therefore, replacing a known deterministic bound
\(1+\beta_{\mathrm{old}}\) by a smaller bound
\(1+\beta_{\mathrm{new}}\) directly produces the improved stochastic
guarantees
\[
    1+\beta_{\mathrm{new}}(1+\Delta)
\]
and
\[
    1+\beta_{\mathrm{new}}
       \max\{1,\alpha(1+\Delta)\}.
\]


The range
\[
    0<\alpha<\frac12
\]
is of particular interest because some of the stronger 
\(\alpha\)-completion bounds used in the analysis of Jäger and Skutella apply
only when \(\alpha\ge\frac12\). 

This thesis investigates the entire parameter range $\alpha \in [0,1]$.

\section*{Objectives and Outline}

The main objective of this thesis is to determine or bound as sharply
as possible the approximation ratio $\lambda_{\alpha}$ for $\alpha \in [0,1]$. 
A central part of the analysis is the derivation of
structural properties of instances whose instance-specific performance ratio
can approach the approximation ratio.

\Cref{chap:introduction} introduces the notation and definitions used throughout the thesis.
\Cref{chap:alpha_all} presents reductions valid for all $\alpha$ in $[0,1]$.
\Cref{chap:alpha_g} analyzes the range $\alpha\in \Big[\frac12,1\Big]$, while \Cref{chap:alpha_l}
treats the endpoint $\alpha=0$ and addresses $\alpha\in\Big(0,\frac12\Big]$ through structural,
analytical, and computational bounds. Finally,
\Cref{chap:outlook} summarizes the results and gives insight into 
promising research directions.

The main contributions of this thesis are:
\begin{itemize}
    \item an exact analysis for $\alpha\in\Big[\frac12,1\Big]$;
    \item a family of instances whose instance-specific performance ratio is unbounded  for $\alpha=0$;
    \item structural reductions that restrict the instances relevant to the analysis of the approximation ratio for $\alpha\in\Big(0,\frac12\Big]$;
    \item upper- and lower-bound constructions for $\alpha\in\Big(0,\frac12\Big]$; and
    \item a mixed-integer nonlinear formulation used to explore the remaining gap computationally.
\end{itemize}

\section*{Use of External Tools and Artificial Intelligence}

Grammarly and WolframAlpha (including in-built AI tools) were used in the preparation of this thesis
exclusively for the following purposes:
\begin{itemize}
    \item translating text,
    \item improving language and grammar,
    \item computing maxima and derivatives, specifically: \eqref{Wolfram1}, \eqref{Wolfram2}, \eqref{Wolfram3}, \eqref{Wolfram4}, \eqref{Wolfram5}, \eqref{Wolfram6} and \eqref{Wolfram7}.
\end{itemize}

GitHub Copilot, integrated into Visual Studio Code, was used exclusively to assist
with improving and debugging the code implementing the MINLP and help in the formatting of
the figures: \ref{fig:g1-g2-comparison}, \ref{fig:construction}, \ref{fig:interval}, \ref{fig:objective-comparison} and \ref{fig:objective-zoom}.

\section{Definitions}
We begin by introducing several fundamental scheduling concepts that will be used throughout this thesis.

\begin{definition}[$P \!\parallel\! \sum w_j C_j^{\alpha}$]
    $P \!\parallel\! \sum w_j C_j^{\alpha}$ is defined to be the problem of
    finding a schedule for an instance $I$ that minimizes the total weighted completion time $\sum w_j C_j^{\alpha}$, where $w_j$ is the weight of job $j$ and
    $C_j^{\alpha}$ is its $\alpha$-completion time on an arbitrary number of identical parallel machines.
\end{definition}

\begin{definition}[Instance]
    An instance $I$ of $P \!\parallel\!\sum w_j C_j^{\alpha}$ is defined to be a set of $n$ jobs, each with a processing time $p_j > 0$ and a weight $w_j \ge 0$, as well as a set of $m$ identical parallel machines on which the jobs are to be scheduled.
\end{definition}

\begin{definition}[Schedule]
    A schedule $\sigma(I)$ of an instance $I$ of $P \!\parallel\!\sum w_j C_j^{\alpha}$ is defined to be an assignment of every job to a machine and a start time $S_j$ such that no two jobs that are
    assigned to the same machine are processed at the same time and all jobs are completed meaning $C_j(\sigma(I)) \ge S_j(\sigma(I)) + p_j$ for all jobs $j$.
\end{definition}

\begin{definition}[$\alpha$-completion time]
    Given an instance I of $P \!\parallel\! \sum w_j C_j^{\alpha}$, as well as a schedule $\sigma(I)$, the $\alpha$-completion 
    time $C_j^{\alpha}(\sigma(I))$ is defined as: 
    \begin{align*}
    C_j^{\alpha}(\sigma(I)) \coloneqq S_{j}(\sigma(I)) + \alpha p_j
    \end{align*}
\end{definition}

\begin{definition}[List schedule]
    A list schedule of an instance $I$ of $P \!\parallel\! \sum w_j C_j^{\alpha}$ is a schedule that results from considering the jobs in an ordered list
    $l$ and scheduling them in that order,
    assigning each job to a machine that becomes available earliest.
    We say that a list schedule $\sigma(I)$ is associated with a list $l$ if the schedule results out of $l$ and denote this by: 
    \[
    \sigma(I) \sim l
    \]
\end{definition}


\begin{definition}[Machine load]
The machine load $L^{i}(\sigma(I))$ of machine $i$ in a list schedule $\sigma(I)$ for an instance $I$  of  $P \!\parallel\! \sum w_j C_j^{\alpha}$,
is defined as
\[
    L^{i}(\sigma(I)) \coloneqq \sum_{j \in J_i(\sigma(I))} p_j,
\]
where $J_i(\sigma(I))$ is the set of jobs assigned to machine $i$ in schedule $\sigma(I)$.
For the purposes of this thesis, it suffices to define this only for list schedules. 
\end{definition}


\begin{definition}[Objective value]
Given an instance \(I\) of $P \!\parallel\! \sum w_j C_j^{\alpha}$ and a schedule \(\sigma(I)\), 
the objective value \(\gamma_{\alpha}(\sigma(I))\) is defined as
\[
    \gamma_{\alpha}(\sigma(I))
    \coloneqq
    \sum_{j \in I} w_j C_j^{\alpha}(\sigma(I)).
\]
Furthermore, 
\(\gamma_{\alpha}^{*}(I)\) is defined as the minimum objective value over all list schedules of instance \(I\), and \(\gamma_{\alpha}^{W}(I)\) is defined as the maximum objective value over all list schedules of instance \(I\):
\[
    \gamma_{\alpha}^{*}(I)
    \coloneqq   \min_{\sigma(I)} \gamma_{\alpha}(\sigma(I)), \qquad
    \gamma_{\alpha}^{W}(I)
    \coloneqq   \max_{\sigma(I)} \gamma_{\alpha}(\sigma(I)).
\]  
\end{definition}

\begin{definition}[$\mathrm{WSPT}$-schedule]
    A  weighted shortest processing time $(\mathrm{WSPT})$ schedule  $\sigma^{\mathrm{WSPT}}$ for an instance $I$ of $\Problem$ is a list schedule for which there exists
    an associated list where the jobs are in order of non-increasing ratio $w_j / p_j$.
\end{definition}

\begin{definition}[Worst list schedule]
    A worst list schedule $\sigma^{W}(I)$ for an instance $I$ of $P \!\parallel\!\sum w_j C_j^{\alpha}$ is defined to be a list schedule with maximal objective value for the given instance $I$.
    For the special case in which $w_j = p_j$ for all jobs $j$, a worst list schedule is a $\mathrm{WSPT}$ schedule.
\end{definition}

\begin{definition}[Optimal schedule]
    An optimal schedule $\sigma^{*}(I)$ for an instance $I$ of $P \!\parallel\! \sum w_j C_j^{\alpha}$ is defined to be a schedule with minimal objective value for the given instance $I$.
    Note that an optimal schedule for this problem is necessarily a list schedule. 
\end{definition}

\begin{definition}[Instance-specific performance ratio and Approximation ratio]
    Given an instance $I$ of $\Problem$, the instance-specific performance ratio of $I$ is defined as
    \[
        \lambda_{\alpha}(I) \coloneqq \frac{\gamma_{\alpha}^{W}(I)}{\gamma_{\alpha}^{*}(I)} 
    \]

    The approximation ratio is defined as
    \[
        \lambda_{\alpha} \coloneqq \sup_{I} \lambda_{\alpha}(I)
    \]

    We will also use 
    \[
    \lambda_{\alpha}(\sigma(I)) \coloneqq \frac{\gamma_{\alpha}(\sigma(I))}{\gamma_{\alpha}^*(I)}.
    \]
\end{definition}

\begin{definition}[$\lambda_\alpha$-dominance]
An instance \(I'\) is said to \(\lambda_\alpha\)-dominate an instance \(I\) if
\[
    \lambda_\alpha(I') \geq \lambda_\alpha(I).
\]
A subclass \(\mathcal D\) of instances of the problem $P \!\parallel\! \sum w_j C_j^{\alpha}$ is called
\(\lambda_\alpha\)-dominant if for every instance \(I\), there exists
an instance \(I'\in\mathcal D\) such that
\[
    \lambda_\alpha(I') \geq \lambda_\alpha(I).
\]
We say that an instance $I$ is $\lambda_{\alpha}$-dominant, if it is an element of a $\lambda_{\alpha}$-dominant subclass.
\end{definition}

\chapter{Structural Results for
\texorpdfstring{$\alpha \in [0,1]$}{alpha in [0,1]}}\label{chap:alpha_all}

We generalize Corollaries 1 and 2 from \cite{Schwiegelshohn:2011} so as to determine a $\lambda_{\alpha}$-dominant class of instances
with nontrivial structural properties. The results in \cite{Schwiegelshohn:2011} are stated implicitly for 
the case $\alpha=1$ only. The following two results extend directly and hold 
for the entire parameter range of $\alpha \in [0,1]$.

\begin{lemma}\label{lem: lemm1}
For every instance \(I\) of
$P\!\parallel\!\sum_j w_j C^{\alpha}_j$ with
\(\alpha \in [0,1]\), there exists an instance \(I'\) of
$\Problem$ such that
\[
    \lambda^{\mathrm{WSPT}}_{\alpha}(I') \geq \lambda^{\mathrm{WSPT}}_{\alpha}(I),
\]

where \[
\lambda^{\mathrm{WSPT}}_{\alpha}(I)
\coloneqq
\max_{\substack{\sigma(I)\text{ is a}\\ \mathrm{WSPT}\text{ schedule}}}
\lambda_{\alpha}(\sigma(I)).
\]
\end{lemma}

\begin{proof}
The proof follows the ideas of \cite{Schwiegelshohn:2011}.
Let
\[
        r_1>r_2>\cdots>r_d
\]
be the distinct ratios $w_j/p_j$ occurring in instance $I$, and set $r_{d+1}=0$.
Therefore $d$ is of course the number of distinct ratios. 
By setting $r_{d+1}=0$, we can write

\[
        r_i=\sum_{k=i}^d (r_k-r_{k+1}).
\]
Let $i(j)$ be the index  of the ratio for which
\[
      r_{i(j)}=\frac{w_j}{p_j}.
\]

Then it holds that
\[
        w_j
        =
        r_{i(j)}p_j
        =
        \left(\sum_{k=i(j)}^d (r_k-r_{k+1})\right)p_j.
\]
And thus for any list schedule $\sigma(I)$,
\begin{equation}
    \label{eq: Ik identity}
\begin{aligned}
        \sum_{j\in I} w_j C_j^{\alpha}(\sigma(I))
        &=
        \sum_{j\in I}
        \left(
            \sum_{k=i(j)}^d (r_k-r_{k+1})
        \right)p_j C_j^{\alpha}(\sigma(I))                                      \\
        &=
        \sum_{k=1}^d \left((r_k-r_{k+1})
        \sum_{\substack{j\in I\\ i(j)\leq k}} p_j C_j^{\alpha}(\sigma(I))\right).
\end{aligned}
\end{equation}
Fix a cost-maximizing $\mathrm{WSPT}$-schedule $\sigma^{\mathrm{WSPT}}(I)$ for $I$.
Next we define
\[
        I_k \coloneqq \{j\in I: i(j)\leq k\} \qquad \text{for } k=1,\ldots,d.
\]

$I_k$ consists of exactly those jobs whose ratios are among the $k$
largest $w_j/p_j$ ratios in $I$. Now regard each $I_k$ as an instance of
$\Problem$ by setting $w_j=p_j$ for every job $j\in I_k$.
Let $\gamma^{*}_{\alpha}(I_k)$ denote the cost of an optimal schedule.
For every instance $I_k$, we consider the $\mathrm{WSPT}$-schedule $\sigma^{\mathrm{WSPT}}(I_k)$ that is obtained by restricting $\sigma^{\mathrm{WSPT}}(I)$ to the jobs of $I_k$ and denote it by $\sigma^{\mathrm{WSPT}}(I_k)$.
The $\alpha$-completion time of every job $j\in I_k$ is the same in $\sigma^{\mathrm{WSPT}}(I_k)$ as in
$\sigma^{\mathrm{WSPT}}(I)$, because every job excluded from $I_k$
appears later in the list corresponding to $I$.
We can thus write the cost of our original instance $I$ as a sum over the instances $I_k$ for both the $\mathrm{WSPT}$ and optimal
cost schedule. To this end, we will make use of \eqref{eq: Ik identity}.
We can write:
\[
\begin{aligned}
        \lambda^{WSPT}_{\alpha}(I)
        =
        \lambda_{\alpha}(\sigma^{\mathrm{WSPT}}(I))
        &=
        \frac{
        \sum_{k=1}^d (r_k-r_{k+1}) \gamma_{\alpha}(\sigma^{\mathrm{WSPT}}(I_k))
        }{
        \sum_{k=1}^d \left(
        (r_k-r_{k+1})
        \sum_{j\in I_k} p_j C_j^{\alpha}(\sigma^{*}(I))\right)
        }                                                    \\
        &\le
        \frac{
        \sum_{k=1}^d (r_k-r_{k+1}) \gamma_{\alpha}(\sigma^{\mathrm{WSPT}}(I_k))
        }{
        \sum_{k=1}^d
        (r_k-r_{k+1}) \gamma^{*}_{\alpha}(I_k)
        }                                                    \\
        &\leq
        \max_{1\leq k\leq d}
        \frac{\gamma_{\alpha}(\sigma^{\mathrm{WSPT}}(I_k))}{\gamma^{*}_{\alpha}(I_k)} \\
        &=
        \max_{1\leq k\leq d}\lambda_{\alpha}(\sigma^{\mathrm{WSPT}}(I_k))\le \max_{1\leq k\leq d}\lambda^{\mathrm{WSPT}}_{\alpha}(I_k).
\end{aligned}
\]
The first equality holds since $\sigma^{\mathrm{WSPT}}(I)$ is a cost-maximizing $\mathrm{WSPT}$-schedule for $I$.
Because of 
\[
\sum_{k=1}^d (r_k-r_{k+1}) \gamma_{\alpha}(\sigma^{\mathrm{WSPT}}(I_k)) = \gamma_{\alpha}(\sigma^{\mathrm{WSPT}}(I)),
\]
which is a compacted form of \eqref{eq: Ik identity}, the second equality holds.
For the first inequality, observe that $\sum_{j\in I_k} p_j C_j^{\alpha}(\sigma^{*}(I))$ is the cost of a schedule for the instance $I_k$. It might not be the cost of a list schedule, but since every optimal schedule 
is a list schedule, we have:
\[
\gamma^{*}_{\alpha}(I_k)
\leq
\sum_{j\in I_k}p_j C_j^{\alpha}(\sigma^{*}(I)).
\]

Therefore, choosing an index $k$ maximizing $\lambda^{\mathrm{WSPT}}_{\alpha}(I_k)$ and setting
$I'=I_k$ gives an instance of $\Problem$ with
$\lambda^{\mathrm{WSPT}}_{\alpha}(I')\geq \lambda^{\mathrm{WSPT}}_{\alpha}(I)$.



\end{proof}
It is important to note that since the instance $I'$ constructed in the proof is 
an instance for which $w_j=p_j$ holds for all jobs $j$, any list schedule is a $\mathrm{WSPT}$-schedule.
In particular, there exists a $\mathrm{WSPT}$-schedule $\sigma^{\mathrm{WSPT}}(I')$ such that
\[
    \gamma_{\alpha}^{\mathrm{WSPT}}(I') = \gamma^{W}_{\alpha}(I')
\]
and consequently
\[
    \lambda^{\mathrm{WSPT}}_{\alpha}(I') = \lambda_{\alpha}(I').
\]

We can therefore restrict our attention to instances of $\Problem$ in the remainder of this thesis.
We will work only with cost-maximizing list schedules $\sigma^{W}(I)$ and optimal schedules $\sigma^{*}(I)$.

Lemma \ref{lem: lemm1} implies that the original problem admits a $\lambda_{\alpha}$-dominant subclass of instances satisfying $w_j=p_j$ for all $j$. 
We continue by showing a useful identity.

\begin{lemma}\label{corr:queyranne}
For any schedule $\sigma(I)$, without intermediate idle times, of an instance $I$ of $1\!\parallel\!\sum_{j} p_j C_j^{\alpha}$ and $\alpha \in [0,1]$,
\[
        \sum_{j\in I}p_j C_j^{\alpha}(\sigma(I))
        =
        \frac12\left((2\alpha-1)\left(\sum_{j\in I}p_j^2\right)+(\machLoadMax)^2\right),
\]
where $L^{max}(\sigma(I)) \coloneqq \max_{i \in [m]} \{ L^{i}(\sigma(I))\}$.
\end{lemma}

\noindent
The proof is provided in
\hyperref[app:proof-queyranne]
{Appendix~\ref*{app:detailed-calculations}
(Proof of Lemma~\ref*{corr:queyranne})}.
The result is a generalization of a result by Queyranne~\cite{queyranne1993structure} for the case $\alpha=1$.

We now extend this to instances with an arbitrary number of machines.
This extension is done equivalently to the extension of Corollary 2 in \cite{Schwiegelshohn:2011}.
We can apply Lemma \ref{corr:queyranne} to every machine separately.
Formally, define
\[
J_i(\sigma(I))
\coloneqq
\{j\in I\mid j\text{ is scheduled on }M_i\text{ in }\sigma(I)\}.
\]
Then for any schedule $\sigma(I)$ of instance $I$ of $\Problem$ with $\alpha \in [0,1]$, we have

\[
        \sum_{j\in I}p_j C^{\alpha}_j(\sigma(I))
        =
        \sum_{i=1}^m \sum_{j\in J_i(\sigma(I))}p_j C^{\alpha}_j(\sigma(I))
        \overset{(\ref{corr:queyranne})}{=}
        \frac12\left((2\alpha-1)\left(\sum_{j\in I}p_j^2\right)+\sum_{i=1}^m (\machLoad)^2\right).
\]

Furthermore, define
\[
        \machLoadMin \coloneqq \min_{i \in [m]} \{\machLoad\}
\]
as the minimum machine load in the parallel identical machine scenario.
In any list schedule $\sigma(I)$, no job can start after $\machLoadMin$, and at most
$m-1$ jobs can ($\alpha$)-complete after $\machLoadMin$.
We will make use of this property in \Cref{chap:alpha_g}.


\chapter{Structural Reductions and Exact Approximation Ratio for
\texorpdfstring{$\alpha \in \Big[\frac12,1\Big]$}{alpha in [1/2,1]}}\label{chap:alpha_g}

As in \Cref{chap:alpha_all}, we follow the ideas of \cite{Schwiegelshohn:2011}
to simplify the structure of instances that may attain or approach the
approximation ratio $\lambda_{\alpha}$, focusing on their free parameters.
\section{Structural Reductions}
We next extend Corollaries 3–5 of \cite{Schwiegelshohn:2011}. 
To this end, we restrict attention to \(\alpha \ge \frac12\), 
which ensures that \(2\alpha-1\ge 0\). This nonnegativity condition 
is essential for the following result. We will use the results from this section
as well as the results from \Cref{chap:alpha_all} to show the following result:


\begin{theorem}\label{th: firstframed}
For every \(\alpha\in\Big[\frac12,1\Big]\), the
subclass of instances \(I\) with the
following properties:
\begin{enumerate}[label=(\roman*)]
    \item \(w_j=p_j\) for every job \(j\)

    \item There are \(u<m\) long jobs, all of processing time \(p\);
          all remaining jobs have arbitrarily small processing times

    \item The small jobs can be normalized to have total workload \(m\)

    \item In any cost-maximizing schedule \(\sigma^{W}(I)\), the small jobs are scheduled in
          \([0,1]\), and all long jobs start at time \(1\)

    \item In any optimal schedule \(\sigma^{*}(I)\), every long job starts at time \(0\).
          Every machine $i$ processing small jobs has machine load
          \[
              L^{i}(\sigma^{*}(I))
              =
              L^{\min}(\sigma^{*}(I))
          \]
            and every machine with larger load contains exactly one long job.
\end{enumerate}
is \(\lambda_\alpha\)-dominant.
\end{theorem}

We give the proof at the end of this section.
We will also show that the properties in Theorem~\ref{th: firstframed}
reduce the scheduling problem to a nonlinear optimization problem in one variable.
We begin by generalizing Corollary 3 in \cite{Schwiegelshohn:2011}.

\begin{lemma}\label{corr:splitting-jobs}
For every instance $I$ of $\Problem$ with $\alpha \in \Big[\frac12,1\Big]$, there is an
instance $I'$ of $\Problem$ such that
$\lambda_{\alpha}(I')\geq\LamI$, and every job in $I'$ that starts before
$L^{\min}(\sigma^{W}(I))$ has arbitrarily small processing time where $\sigma^W(I)$ is an arbitrary worst-case list schedule.
\end{lemma}

\begin{proof}
For an arbitrary list schedule $\sigma(I)$, we generate an instance $I'$ by taking a job $j_0$ from
$I$ and replacing it in schedule $\sigma(I)$ with two jobs $j_1$ and $j_2$ such
that
\[
        p_{j_1}+p_{j_2}=p_{j_0},
        \qquad
        C_{j_2}(\sigma(I'))=C_{j_0}(\sigma(I)) \qquad \text{and}
        \qquad
        C_{j_1}(\sigma(I'))=C_{j_0}(\sigma(I))-p_{j_2}.
\]


This replacement does not alter the completion time of any other job in
$\sigma(I)$. Therefore, 
\[
       \machLoad=L^i(\sigma(I')) \qquad \forall i \in [m].
\]

And thus, using the identity \eqref{eq:queyranne-identity}
\begin{align*}
        \sum_{j\in I}p_j C^{\alpha}_j(\sigma(I))-\sum_{j\in I'}p_j C^{\alpha}_j(\sigma(I'))
        &=
        \frac12((2\alpha-1)(p_{j_0}^2-p_{j_1}^2-p_{j_2}^2))\\
        &=
        (2\alpha-1)p_{j_1}p_{j_2}\\
        &\ge 0.
\end{align*}
If $C_{j_1}(\sigma(I'))\leq L^{\min}(\sigma(I))$, then 
$\sigma(I')$ is again a list schedule. If however $C_{j_1}(\sigma(I'))\ > L^{\min}(\sigma(I))$, 
then $C_{j_2}$ would not be scheduled at the earliest available time which implies that it does 
not correspond to a list schedule.
It is because of this, that
for any instance $I$, we perform the splitting operation only on $j_0$ if it satisfies
$S_{j_0}(\sigma(I))\leq L^{\min}(\SW)$ and pick the processing times of the split jobs such that $C_{j_1}(\sigma(I'))\leq L^{\min}(\sigma(I))$. 
Splitting the same job into two consecutive parts in an optimal schedule of $I$ might not directly result 
in a list schedule, but rescheduling of the jobs that we split, will only decrease the costs further by "evening" the 
machine loads. Therefore, the minimum cost reduction for the optimal schedule is the gain
from the two shortened job lengths ($p_{j_1}$ and $p_{j_2}$ instead of $p_{j_0}$).
With
\[
\delta=(2\alpha-1)p_{j_1}p_{j_2}\geq0,
\]
we therefore have
\[
\gamma^{W}_{\alpha}(I') \geq\WorstCost-\delta
\qquad\text{and}\qquad
\gamma^{*}_{\alpha}(I')\leq\OptCost-\delta.
\]
Consequently,
\[
\lambda_{\alpha}(I')= = \frac{\gamma_{\alpha}(I')}{\gamma_{\alpha}^W(I')}
\geq
\frac{\WorstCost-\delta}{\OptCost-\delta}
\geq
\frac{\WorstCost}{\OptCost}
=\LamI,
\]
where the second inequality uses $\WorstCost\geq\OptCost$.
We can repeat this splitting operation arbitrarily often, yielding an instance in which all jobs that
start before $L^{\min}(\sigma^{W}(I))$ have arbitrarily small processing time where $\sigma^W(I)$ is an arbitrary cost-maximizing list schedule.
Since at most $m-1$ jobs can start at $L^{\min}(\sigma^{W}(I))$, all but at most
$m-1$ jobs have arbitrarily small processing time.
\end{proof}

By the preceding reductions, we from now on only consider instances $I$
with $w_j=p_j$ for all jobs and in which all jobs have arbitrarily small
processing time, except for at most $m-1$ \emph{long jobs} that start at
$L^{\min}(\sigma^{W}(I))$ in the cost-maximizing schedule. Let $L(I)\subseteq I$ denote the set of these long jobs.

Furthermore, let $K(I)\subseteq L(I)$ be the largest set of jobs such that
\[
        p_k\geq \frac1{m-|K(I)|}\sum_{j\in I\setminus K(I)}p_j
\]
holds for all $k\in K(I)$. We obtain $K(I)$ by adding long jobs in decreasing
order of their processing times and testing this property.
We will see in the construction of an optimal schedule, that $K(I)$ is defined in a way, 
such that any $j\in K(I)$ has a processing time so large, that it is "worth" reserving 
a machine for $j$ alone. 
All jobs in
$I\setminus L(I)$ are called \emph{small jobs}.

Since every small job has arbitrarily small processing time, we can ignore
the term $(2\alpha-1)\sum_{j\in I\setminus L(I)}p_j^2$ in
Lemma~\ref{corr:queyranne}. Thus, for any list schedule $\sigma(I)$,
\begin{equation}
\label{eq:queyranne-identity}
        \sum_{j\in I}p_j C^{\alpha}_j
        =
        \frac12
        \left(
            (2\alpha-1)\sum_{j\in L(I)}p_j^2
            +
            \sum_{i=1}^m(L^{i}(\sigma(I)))^2
        \right).
\end{equation}
Applying \eqref{eq:queyranne-identity} to $\sigma^{W}(I)$ gives
\begin{equation}
\label{eq:wspt-value}
\begin{aligned}
        \WorstCost
        &=
        \underbrace{\frac12(2\alpha-1)\left(\sum_{j\in L(I)}p_j^2\right)}_{\text{Fixed Cost}}
        +
        \underbrace{\frac12\sum_{j\in L(I)}
            (p_j+L^{\min}(\sigma^{W}(I)))^2}_{\text{Machine load cost with long job}}                    
        +
        \underbrace{\frac12(m-|L(I)|)(L^{\min}(\sigma^{W}(I)))^2}_{\text{Machine load with only short jobs}}.
\end{aligned}
\end{equation}

Next, we address the construction of an optimal schedule $\sigma^{*}(I)$ for such instances.
Consider a list for $I$ in which the long jobs precede all small jobs.
In $\sigma^{*}(I)$, all long jobs start at time
$0$, which is possible because $|L(I)|<m$. We have
\[
        L^{\min}(\sigma^{*}(I))
        \leq
        \frac1{m-|K(I)|}\sum_{j\in I\setminus K(I)}p_j.
\]
If a small job $j_0$ finishes after
\[
        \frac1{m-|K(I)|}\sum_{j\in I\setminus K(I)}p_j,
\]
we split $j_0$ at time
\[
        \frac1{m-|K(I)|}\sum_{j\in I\setminus K(I)}p_j
\]
as in the proof of Lemma~\ref{corr:splitting-jobs}. In a list corresponding to  $\sigma^{*}(I)$, the first new job
$j_1$ replaces $j_0$, while the second new job $j_2$ is placed at the end of
the list. After at most $m-|K(I)|-1$ such splitting steps,
\[
        L^{i}(\sigma^{*}(I))
        =
        L^{\min}(\sigma^{*}(I))
        =
        \frac1{m-|K(I)|}\sum_{j\in I\setminus K(I)}p_j
\]
holds for every machine $i$ that executes a small job. Therefore, every
machine $i$ with
\[
        L^{i}(\sigma^{*}(I))>L^{\min}(\sigma^{*}(I))
\]
executes only a single long job. By \eqref{eq:queyranne-identity} and the
convexity of the square function, it is impossible to further reduce
\[
        \sum_{i=1}^m(L^{i}(\sigma^{*}(I)))^2
\]
by transferring jobs. Hence, $\sigma^{*}(I)$ is optimal.

We next establish a structural reduction that directly generalizes Corollary 5 of \cite{Schwiegelshohn:2011}.

We introduce a transformation that replaces a set $M$ of positive
numbers with a number $p$ that is used $t$ times such that the following
equations hold:
\begin{align}
        \sum_{i\in M} i &= t\cdot p, \\
        \sum_{i\in M} i^2 &= t\cdot p^2.
\end{align}
These equations admit a unique solution:

\begin{equation}
\label{eq:t-and-p}
t=\frac{\left(\sum_{i\in M} i\right)^2}{\sum_{i\in M} i^2}
\qquad\text{and}\qquad
p=\frac{\sum_{i\in M} i^2}{\sum_{i\in M} i}.
\end{equation}

This transformation will drastically reduce the number of distinct processing times 
in the instance which will make it easier for us to express the cost of any schedule. 

In \eqref{eq:queyranne-identity}, this transformation is either applied to the set $L(I)$ of long
jobs or to some makespans. In~\eqref{eq:t-and-p}, $t$ might not be an integer but may be a
rational number with denominator $q$ resulting in an instance with a
fractional number of jobs. To obtain an instance with an integral number of
jobs, we simply replace each machine and each job of the original instance
$I$ with $q$ identical copies. In any cost-maximizing order, all replicas of a job occur
one after another. Therefore, the assignment and order of jobs on each machine is replicated for
the cost maximizing schedule $\sigma^{W}(I)$.
Replicating the schedule of each machine for schedule
$\sigma^{*}(I)$ also yields an optimal schedule for the same reasons. 
Hence, $\gamma^{W}_{\alpha}(I)$ and $\gamma^{*}_{\alpha}(I)$ are simply multiplied by $q$, and
$\LamI$ remains unchanged.

The application of the transformation to the set $L(I)$ or to machine
makespans of a list schedule $\sigma(I)$ does not change
$\sum_{i=1}^m L^{i}(\sigma(I))$ due to Lemma~\ref{corr:splitting-jobs}. 
Furthermore, we have
\begin{equation}
\begin{aligned}
        t
        &=
        \frac{\left(\sum_{j\in L(I)}p_j\right)^2}
        {\sum_{j\in L(I)}p_j^2}                                      \\
        &\leq
        \frac{\left(\sum_{j\in L(I)}p_j\right)^2}
        {|L(I)|\left(\frac{\sum_{j\in L(I)}p_j}{|L(I)|}\right)^2}       \\
        &=
        \frac{|L(I)|\left(\sum_{j\in L(I)}p_j\right)^2}
        {\left(\sum_{j\in L(I)}p_j\right)^2}
        =
        |L(I)|.
\end{aligned}
\end{equation}

This property implies $t\le m-1$ which is important to ensure that we are able to
schedule all $t$ jobs of length $p$ at time zero in the optimal schedule. 

We apply this transformation in the following lemma which is a generalization of Corollary 5 in \cite{Schwiegelshohn:2011}.

\begin{lemma}\label{corr:same-length-jobs}
For every instance $I$ of $\Problem$,
there is an instance $I'$ of the same problem that
contains at most $m-1$ jobs with equal processing time,
while all its other jobs have an arbitrarily small processing time. Furthermore, 
\[
    \lambda_{\alpha}(I') \geq \lambda_{\alpha}(I)
    \qquad \forall \alpha\in\left[\frac12,1\right]
\]
holds.
\end{lemma}

\begin{proof}
Due to Lemma~\ref{corr:splitting-jobs}, we consider a schedule $\sigma^{W}(I)$ in which
only jobs with an arbitrarily small processing time start before
$L^{\min}(\sigma^{W}(I))$.

We apply our transformation to the set $L(I)$ to generate an instance $I'$ in
which all $t$ long jobs have the same processing time. In a cost maximizing order of
both instances, long jobs appear after all small jobs. Therefore, we have
\[
        L^{\min}(\sigma^{W}(I))=L^{\min}(\sigma^{W}(I')).
\]
We discuss the consequences of the transformation to the cost maximizing schedule
and an optimal schedule separately. Due to \eqref{eq:wspt-value} and \eqref{eq:t-and-p}
we have
\[
\begin{aligned}
        2\WorstCost
        &=
        (2\alpha-1)\sum_{j\in L(I)}p_j^2
        +
        \sum_{j\in L(I)}(p_j+L^{\min}(\sigma^{W}(I)))^2                  
        +(m-|L(I)|)(L^{\min}(\sigma^{W}(I)))^2                           \\
        &=
        (2\alpha-1)\sum_{j\in L(I)}p_j^2
        +
        \sum_{j\in L(I)}p_j^2
        +
        2\sum_{j\in L(I)}(p_jL^{\min}(\sigma^{W}(I)))               
        +m(L^{\min}(\sigma^{W}(I)))^2                                    \\
        &=
        (2\alpha-1)tp^2+tp^2
        +2tpL^{\min}(\sigma^{W}(I^{\prime})))
        +m(L^{\min}(\sigma^{W}(I^{\prime})))^2                                  \\[1mm]
        &=
        (2\alpha -1)tp^2
        +t(p+L^{\min}(\sigma^{W}(I^{\prime})))^2
        +(m-t)(L^{\min}(\sigma^{W}(I^{\prime})))^2                              \\
        &=
        2\gamma^{W}_{\alpha}(I').
\end{aligned}
\]

Next we address the optimal schedules for instances $I$ and $I'$. To show
that $\gamma^{*}_{\alpha}(I')\leq \OptCost$ holds, we consider two cases:

\begin{enumerate}[label=(\arabic*)]
\item
$L^{i}(\sigma^{*}(I'))=L^{\min}(\sigma^{*}(I'))$ holds for all machines $i$.

Note that any list schedule $\sigma(I')$ is optimal if
$L^{i}(\sigma(I'))=L^{\min}(\sigma(I'))$ holds for all machines $i$ due to the
convexity of the square function. Furthermore, the transformation does not
change $\sum_{i=1}^m L^{i}(\sigma(I))$ for any list schedule $\sigma(I)$.
Therefore, because of \eqref{eq:wspt-value} and \eqref{eq:t-and-p}:
\[
\begin{aligned}
        2\gamma^{*}_{\alpha}(I)
        &=
        (2\alpha-1)\sum_{j\in L(I)}p_j^2
        +
        \sum_{i=1}^m(L^{i}(\sigma^{*}(I)))^2                 \\
        &=
        (2\alpha-1)t\cdot p^2
        +
        \sum_{i=1}^m(L^{i}(\sigma^{*}(I)))^2                 \\
        &\geq
        (2\alpha-1)t\cdot p^2
        +
        m(L^{\min}(\sigma^{*}(I')))^2
        =
        2\gamma^{*}_{\alpha}(I').
\end{aligned}
\]

\item Suppose that
$L^{i}(\sigma^{*}(I'))>L^{\min}(\sigma^{*}(I'))$ holds for some machine $i$.
In the following sequence of equations, we show that $\gamma^{*}_{\alpha}(I')\leq \gamma^{*}_{\alpha}(I)$
also holds for this case. We later reason for all of these equations.
\begin{align}
    2\gamma^{*}_{\alpha}(I)
    &=
    (2\alpha-1)\sum_{j\in L(I)}p_j^2
    +\sum_{j\in K(I)}p_j^2
    +(m-|K(I)|)\bigl(L^{\min}(\sigma^{*}(I))\bigr)^2
    \label{eq:cost-decomposition}
    \\
    &=
    (2\alpha-1)tp^2+t'(p')^2
    +(m-|K(I)|-k)\bigl(L^{\min}(\sigma^{*}(I))\bigr)^2
    \label{eq:cost-substitution}
    \\
    &\geq
    (2\alpha-1)tp^2+t'(p')^2
    +(m-t')
    \left(
        \frac{m-|K(I)|-k}{m-t'}L^{\min}(\sigma^{*}(I))
    \right)^2
    \label{eq:cost-first-bound}
    \\
    &\geq
    (2\alpha-1)tp^2+tp^2
    +(m-t)\bigl(L^{\min}(\sigma^{*}(I'))\bigr)^2
    \label{eq:cost-second-bound}
    \\
    &=
    2\alpha tp^2
    +(m-t)\bigl(L^{\min}(\sigma^{*}(I'))\bigr)^2
    =
    2\gamma^{*}_{\alpha}(I').
    \label{eq:cost-final}
\end{align}

First, we introduce the size
\[
        k\coloneqq
        \frac{\sum_{j\in L(I)\setminus K(I)}p_j}
        {L^{\min}(\sigma^{*}(I))}
        \leq
        |L(I)|-|K(I)|.
\]
If necessary, we duplicate our instance, as in \eqref{eq:t-and-p}, to guarantee an integral $k$.
\eqref{eq:cost-decomposition} is an application of \eqref{eq:queyranne-identity}. \eqref{eq:cost-substitution} is
the application of the exchange argument \eqref{eq:t-and-p} on schedule $\sigma^{*}(I)$.\medskip

We apply our transformation twice to \eqref{eq:cost-substitution}:
\begin{enumerate}[label=(\alph*)]
\item on the set $L(I)$ to transform $\sum_{j\in L(I)}p_j^2$;
\item on the makespan of all machines that execute jobs in $K(I)$ and of
$k$ additional machines with
$L^{i}(\sigma^{*}(I))=L^{\min}(\sigma^{*}(I))$.
\end{enumerate}
The second transformation produces the values $t'$ and $p'$. Remember that
we use replication to guarantee an integral $t'$. Due to the convexity of the
square function and $0\leq p_j\leq L^{\min}(\sigma^{*}(I))$ for all jobs
$j\in L(I)\setminus K(I)$, we have

\begin{align}\label{eq:first-moment-prime}
    t'\cdot p'
    &=
    \sum_{j\in K(I)}p_j
    +k\cdot L^{\min}(\sigma^{*}(I))\\
    &=
    \sum_{j\in K(I)}p_j +\sum_{j\in L(I)\setminus K(I)}p_j
    =
    \sum_{j\in L(I)}p_j
    =
    t\cdot p,
\end{align}

\begin{equation}\label{eq:second-moment-prime}
\begin{aligned}
    t'\cdot(p')^2
    &=
    \sum_{j\in K(I)}p_j^2
    +k\bigl(L^{\min}(\sigma^{*}(I))\bigr)^2
    \\
    &\geq
    \sum_{j\in K(I)}p_j^2
    +\sum_{j\in L(I)\setminus K(I)}p_j^2
    =
    \sum_{j\in L(I)}p_j^2
    =
    t\cdot p^2,
    \\
    &\Rightarrow
    p'\geq p
    \quad\text{and}\quad
    t'\leq t.
\end{aligned}
\end{equation}

In \eqref{eq:cost-first-bound}, we evenly distribute the total load of all small jobs, which is
\[
        (m-|K(I)|-k)L^{\min}(\sigma^{*}(I))
\]
onto $m-t'$ machines. Indeed, since 
$t' \le t \le |L(I)| \le m-1$, $m-t'$ is a feasible number of machines. 
This distribution reduces the total
weighted completion time due to the convexity of the square function.

Due to \eqref{eq:first-moment-prime}, only
$m-t\leq m-t'$ machines are available to execute the small jobs. Therefore, we have
\[
\begin{aligned}
\frac{m-|K(I)|-k}{m-t'}L^{\min}(\sigma^*(I))
&=
\frac{\sum_{j\in I\setminus K(I)}p_j
      -\sum_{j\in L(I)\setminus K(I)}p_j}{m-t'} \\
&=
\frac{\sum_{j\in I\setminus L(I)}p_j}{m-t'} \\
&\leq
\frac{\sum_{j\in I\setminus L(I)}p_j}{m-t} \\
&=
L^{\min}(\sigma^*(I'))
\leq
L^i(\sigma^*(I')).
\end{aligned}
\]
Because of this and \eqref{eq:second-moment-prime}, \eqref{eq:cost-second-bound} holds. 
\end{enumerate}
\end{proof}

\begin{proof}[Proof of Theorem \ref{th: firstframed}]
 We arrive at this result by applying the construction steps from Lemma \ref{lem: lemm1}, Lemma \ref{corr:queyranne},
 Lemma \ref{corr:splitting-jobs} and \ref{corr:same-length-jobs} in succession.
\end{proof}

\section{Determining the Approximation Ratio}
We are now ready to determine the approximation ratio.
We consider the $\lambda_{\alpha}$-dominant class determined in Theorem $\ref{th: firstframed}$ whose instances have the properties:

\begin{framed}\label{prop:instance-structure}
\begin{enumerate}[label=(\roman*)]
    \item \(w_j=p_j\) for every job \(j\).

    \item There are \(u<m\) long jobs, all of processing time \(p\);
          all remaining jobs have arbitrarily small processing times.

    \item The small jobs have total workload \(m\).

    \item In any \(\sigma^{W}(I)\), the small jobs are scheduled in
          \([0,1]\), and all long jobs start at time \(1\).

    \item In any \(\sigma^{*}(I)\), every long job starts at time \(0\).
          Every machine $i$ processing small jobs has load
          \[
              L^{i}(\sigma^{*}(I))
              =
              L^{\min}(\sigma^{*}(I)).
          \]
          Every machine with larger load contains exactly one long job.
\end{enumerate}
\end{framed}

We will determine the approximation ratio by determining the cost of the worst-case schedule and the optimal schedule 
for both cases of Lemma~\ref{corr:same-length-jobs} separately. 
We will then determine the maximum of the two ratios to find the approximation ratio.

\subsection*{Case (1) of Lemma~\ref{corr:same-length-jobs}}

For case (1) of Lemma~\ref{corr:same-length-jobs},
\[
    L^{i}(\sigma^{*}(I)) = L^{\min}(\sigma^{*}(I))
    \qquad \forall i \in [m].
\]
This implies that all machines have equal load in $\sigma^{*}(I)$ which is depicted in Figure \ref{figure-scheduleawde}.

\begin{figure}[htbp]
\centering
\setlength{\fboxsep}{6pt}
\fbox{%
\begin{minipage}{0.94\textwidth}
    \begin{minipage}[t]{0.47\linewidth}
        \centering
        \scalebox{0.8}{\schedulecaseone}
        \par\smallskip
        \small(a) \(\sigma^{*}(I_1)\)
    \end{minipage}
    \hfill
    \begin{minipage}[t]{0.47\linewidth}
        \centering
        \scalebox{0.8}{\scheduleworstcase}
        \par\smallskip
        \small(b) \(\sigma^{W}(I_1)\)
    \end{minipage}
\end{minipage}%
}
\caption{Schedules for case (1). Shaded regions represent the load of an unspecified number of short jobs.}\label{figure-scheduleawde}
\end{figure}

In case (1) we are able to distribute the entire processing time $up + m$ of all jobs
in the instance evenly onto $m$ machines. This is only possible, if
\[
    m + up \ge mp
    \quad\Longleftrightarrow\quad
    p \le \frac{1}{1-\frac{u}{m}}.
\]
Let $I_1$ be an element of the $\lambda_{\alpha}$-dominant class from Theorem \ref{th: firstframed} for which additionally $p\le \frac{1}{1-\frac{u}{m}}$ holds.
We can then explicitly write down the cost of $\sigma^{*}(I_1)$ and $\sigma^{W}(I_1)$:
\begin{align*}
    \gamma^{*}_{\alpha}(I_1)
    &=
    \frac12 \left(
        \bigl(2\alpha-1\bigr)
        \sum_{j \in L(I_1)} p_j^2
        +
        \sum_{i=1}^m
        \bigl(L^{i}(\sigma^{*}(I_1))\bigr)^2
    \right) \\
    &=
    \frac12 \left(
        \bigl(2\alpha-1\bigr)up^2
        +
        m\left(\frac{m+up}{m}\right)^2
    \right)
\end{align*}


and
\begin{align*}
    \gamma^{W}_{\alpha}(I_1)
    &=
    \frac12 \left(
        \bigl(2\alpha-1\bigr)
        \sum_{j \in L(I_1)} p_j^2
        +
        \sum_{i=1}^m
        \bigl(L^{i}(\sigma^{W}(I_1))\bigr)^2
    \right) \\
    &=
    \frac12 \left(
        \bigl(2\alpha-1\bigr)up^2
        +
        m-u
        +
        u(p+1)^2
    \right).
\end{align*}


Hence
\begin{align}
    \label{Wolfram1}
    \lambda_{\alpha}(I_1)
    &=
    \frac{
        \frac12\left(
            \bigl(2\alpha-1\bigr)up^2
            + m-u + u(p+1)^2
        \right)
    }{
        \frac12\left(
            \bigl(2\alpha-1\bigr)up^2
            + m\left(\frac{m+up}{m}\right)^2
        \right)
    } \notag \\
    &=
    \frac{
        \bigl(2\alpha-1\bigr)up^2
        + m-u + up^2 + 2up + u
    }{
        \bigl(2\alpha-1\bigr)up^2
        + \frac{m^2+2mup+u^2p^2}{m}
    } \notag \\
    &=
    \frac{
        2\alpha\frac{u}{m}p^2 + 1 + 2\frac{u}{m}p
    }{
        \bigl(2\alpha-1\bigr)\frac{u}{m}p^2
        + 1 + 2\frac{u}{m}p
        + \left(\frac{u}{m}\right)^2p^2
    }.
\end{align}
And we can finally maximize this function over the parameters $u$ and $p$ to find the approximation ratio for case (1).
Case (1) is defined by the following constraints:
\[
    0 \le \frac{u}{m} \le 1,
    \qquad
    0 \le p \le \frac{1}{1-\frac{u}{m}}.
\]
The approximation ratio restricted to instances aligning with case (1), which we denote by $\lambda_{\alpha}^{(1)}$, therefore is

\[
\lambda_{\alpha}^{(1)} = 1+\frac{1}{(1+\sqrt{2\alpha})^2}.
\]

\subsection*{Case (2) of Lemma~\ref{corr:same-length-jobs}}
Now consider case (2) of Lemma~\ref{corr:same-length-jobs}. In this case, we have
\[
    L^{i}(\sigma^{*}(I)) > L^{\min}(\sigma^{*}(I))
    \qquad \text{for some machine } i.
\]
This can be expressed as,
\[
mp > m + up \quad\Longleftrightarrow\quad p > \frac{1}{1-\frac{u}{m}}.
\]

\begin{figure}[htbp]
\centering
\setlength{\fboxsep}{6pt}
\fbox{%
\begin{minipage}{0.94\textwidth}
    \begin{minipage}[t]{0.47\linewidth}
        \centering
        \scalebox{0.8}{\schedulecasetwo}
        \par\smallskip
        \small(a) \(\sigma^{*}(I_2)\)
    \end{minipage}
    \hfill
    \begin{minipage}[t]{0.47\linewidth}
        \centering
        \scalebox{0.8}{\scheduleworstcase}
        \par\smallskip
        \small(b) \(\sigma^{W}(I_2)\)
    \end{minipage}
\end{minipage}%
}
\caption{Schedules for case (2). Shaded regions represent the load of an unspecified number of short jobs.}\label{fig:schedule-comparison}
\end{figure}

Let $I_2$ be an element of the dominant class from Theorem \ref{th: firstframed} for which $p > \frac{1}{1-\frac{u}{m}}$ holds.
Then we can again find a closed form for the cost of $\sigma^{*}(I_2)$ and $\sigma^{W}(I_2)$.
In this case, a job of processing time $p$ does not share its machine with any short jobs in the 
optimal schedule, as is depicted in Figure \ref{fig:schedule-comparison}.

For $\sigma^{*}(I_2)$,
\begin{align*}
    \gamma^{*}_{\alpha}(I_2)
    &=
    \frac12 \left(
        \bigl(2\alpha-1\bigr)
        \sum_{j \in L(I_2)} p_j^2
        +
        \sum_{i=1}^m
        \bigl(L^{i}(\sigma(I_2))\bigr)^2
    \right) \\
    &=
    \frac12 \left(
        \bigl(2\alpha-1\bigr)up^2
        +
        (m-u)\left(\frac{m}{m-u}\right)^2
        +
        up^2
    \right) \\
    &=
    \frac12 \left(
        2\alpha up^2
        +
        (m-u)\left(\frac{m}{m-u}\right)^2
    \right).
\end{align*}



$\sigma^{W}(I_2)$ remains structurally unchanged compared to case (1), so
\begin{align}
    \label{Wolfram2}
    \lambda_{\alpha}(I_2)
    &=
    \frac{
        \frac12\left(2\alpha up^2 + m + 2up\right)
    }{
        \frac12\left(2\alpha up^2 + (m-u)\left(\frac{m}{m-u}\right)^2\right)
    } \notag\\
    &=
    \frac{
        2\alpha \frac{u}{m}p^2 + 1 + 2\frac{u}{m}p
    }{
        2\alpha \frac{u}{m}p^2 + \frac{m}{m-u}
    } \notag\\
    &=
    \frac{
        2\alpha \frac{u}{m}p^2 + 1 + 2\frac{u}{m}p
    }{
        2\alpha \frac{u}{m}p^2 + \frac{1}{1-\frac{u}{m}}
    }.
\end{align}

Case (2) is restricted by the following constraints:
\[
    0 \le \frac{u}{m} \le 1,
    \qquad
    \frac12 \le \alpha \le 1,
    \qquad
    p > \frac{1}{1-\frac{u}{m}}.
\]
We maximize to obtain
\[
    \lambda_{\alpha}^{(2)} =\frac{1}{1-\frac{2}{(2\sqrt{\alpha}+\sqrt{2})^2}},
    \qquad \text{for } \alpha \ge \frac12,
\]
Where we used the same convention for the naming of $\lambda_{\alpha}^{(2)}$ as we did for case (1).
\subsection*{Dominant Case}
It remains to compare the two cases and find the 
maximum for every value of $\alpha$.
This yields the following result.
\begin{result}\label{Res:Case-Comp}
\[
\boxed{
\lambda_{\alpha}^{(2)}>\lambda_{\alpha}^{(1)} \qquad \forall \alpha \in \Bigg[\frac12,1\Bigg]
}
\]
\end{result}
Details on how we arrive at this result are provided in \hyperref[app:proof-cases]
{Appendix~\ref*{app:detailed-calculations}
(Proof of Result~\ref*{Res:Case-Comp})}.
Consequently, for \(\alpha \in \Big[\frac12,1\Big]\), the approximation ratio of the \(\mathrm{WSPT}\)-rule for $\Problem$ is
\[
\lambda_{\alpha} = \frac{1}{1-\frac{2}{(2\sqrt{\alpha}+\sqrt{2})^2}} = 1+\frac{1}{2(\alpha+\sqrt{2\alpha})}.
\]
This result was independently obtained by Jäger and Skutella in \cite{JaegerSkutella2018}, Theorem 7.
It is easy to see that $1+\frac{1}{2(\alpha+\sqrt{2\alpha})}$ is decreasing as a function of $\alpha$. 
Therefore, its maximum is attained for $\alpha =\frac12$ and evaluates to $\lambda_{\frac12}=\frac43$. 
This will be used in subsequent proofs.



\chapter{Structural Reductions and Approximation Bounds for
\texorpdfstring{$\alpha \in \Big[0,\frac12\Big)$}{alpha in [0,1/2)}}\label{chap:alpha_l}

For the previously unresolved case $\alpha \in [0,\frac12)$, we seek
bounds on the approximation ratio. To this end, we follow the approach
developed in \cite{Schwiegelshohn:2011}, aiming to establish structural
properties of $\lambda_{\alpha}$-dominant instances $I$ that simplify
the determination of $\gamma_{\alpha}^W(I)$ and $\gamma_{\alpha}^*(I)$.


\section{Unboundedness for \texorpdfstring{$\alpha = 0$}{alpha equals 0}}\label{sec: alpha zero}
We show that for $\alpha=0$, the $\mathrm{WSPT}$-rule is arbitrarily bad for $P \!\parallel\! \sum_{j} w_j C_j^{\alpha}$.\medskip

Consider the instance $I$ with three jobs and processing times $p_1 =1$, $p_2=1$, $p_3=p$ and weights
$w_1 =p_1$, $w_2=p_2$, $w_3=p_3$ on $m=2$ machines.
Define the schedules through list correspondance:  $\sigma^{*}(I) \sim [p_1,p_3,p_2]$ and $\sigma^{W}(I) \sim [p_1,p_2,p_3]$.

\begin{figure}[ht]
\centering
\begin{subfigure}[t]{0.47\textwidth}
\centering
\begin{tikzpicture}[
    x=0.7cm,
    y=0.7cm,
    job/.style={
        draw,
        thick,
        fill=white,
        font=\small
    },
    machine/.style={font=\small}
]

% ------------------------------------------------
% Template idea:
% row      = displayed vertical row in the picture
% machine  = actual machine label shown on the left
%
% Use \machinerow{row}{machine label}
% Use \dotsrow{row}
% Use \jobblock{row}{start}{length}{job label}
% Use \jobblock[jobhighlight]{row}{start}{length}{job label} for shading
% ------------------------------------------------

% Displayed machine rows
\machinerow{1}{1}
\machinerow{2}{2}



% Example jobs
\jobblock{1}{0}{1}{1}
\jobblock{1}{1}{1}{1}
\jobblock{2}{0}{3}{p}



\end{tikzpicture}

\caption{Schedule $\sigma^{*}(I)$.}
\end{subfigure}
\hfill
\begin{subfigure}[t]{0.47\textwidth}
\centering
\begin{tikzpicture}[
    x=0.7cm,
    y=0.7cm,
    job/.style={
        draw,
        thick,
        fill=white,
        font=\small
    },
    machine/.style={font=\small}
]

% ------------------------------------------------
% Template idea:
% row      = displayed vertical row in the picture
% machine  = actual machine label shown on the left
%
% Use \machinerow{row}{machine label}
% Use \dotsrow{row}
% Use \jobblock{row}{start}{length}{job label}
% Use \jobblock[jobhighlight]{row}{start}{length}{job label} for shading
% ------------------------------------------------

% Displayed machine rows
\machinerow{1}{1}
\machinerow{2}{2}



% Example jobs
\jobblock{1}{0}{1}{1}
\jobblock{2}{0}{1}{1}
\jobblock{2}{1}{3}{p}



\end{tikzpicture}

\caption{Schedule $\sigma^{W}(I)$.}
\end{subfigure}
\caption{Schedules for $\alpha=0$.}
\label{fig:instance-43-ratio}
\end{figure}
Therefore,
\begin{align*}
    \lambda_{0} \ge \frac{\gamma_{0}(\sigma^{W}(I))}{\gamma_{0}(\sigma^{*}(I))} =\frac{(p+1)^2+1-p^2-2}{2}=p
\end{align*}

which is unbounded in $p$. Thus the approximation ratio for $\alpha=0$ is unbounded.

\section{Structural Results}\label{sec:Struct-Results}
For $\alpha < \frac12$, Lemma \ref{corr:splitting-jobs} does not hold. 
It is because of this that in this section, we present some alternative results that allow us to restrict the analysis to instances
with stronger structural properties.
We will formally see why Lemma \ref{corr:splitting-jobs} does not hold after we introduce the concept of pairwise unification (Definition \ref{def:pairwise-unification}) and the corresponding
result (Lemma \ref{corr:unification}).

In this section we will work towards the following result:
\begin{theorem}\label{thm: secondframe}
For \(\alpha\in\Big(0,\frac12\Big]\), the collection of instances with the following properties:
\begin{enumerate}[label=(\roman*)]
    \item \(w_j=p_j\) for every job \(j\).

    \item There exist a cost-maximizing list schedule
      \(\sigma^{W}(I)\) and an optimal schedule
      \(\sigma^{*}(I)\) such that:
      \begin{enumerate}[label=(\alph*)]
          \item No two jobs are assigned to the same machine
                in both schedules

          \item Each machine contains at most \(m\) jobs
                in either schedule, and
                \[
                    n \leq 3m-2
                \]

          \item There exists a job \(j^*\) that is a singleton
                in \(\sigma^{*}(I)\) and finishes last
                in \(\sigma^{W}(I)\), possibly tied with
                other jobs

          \item The job \(j^*\) is not a singleton
                in \(\sigma^{W}(I)\)

          \item Every job larger than \(j^*\) is scheduled alone on a machine
                in \(\sigma^{*}(I)\). In particular, at least
                one largest job is a singleton in
                \(\sigma^{*}(I)\)

          \item With \(p_{\max}\coloneqq\max_j p_j\),
                \[
                    L^{\max}\bigl(\sigma^{W}(I)\bigr)
                    \leq \frac43 p_{\max}
                \]
      \end{enumerate}
\end{enumerate}

is a \(\lambda_\alpha\)-dominant class.
\end{theorem}



\subsection*{Unification}

\begin{definition}[Pairwise Unification]\label{def:pairwise-unification}
    The pairwise unification of two jobs $j,k$ of instance $I$ for a schedule $\sigma(I)$ is defined by replacing jobs $j,k$
    with a single job $l$ of
    processing time $p_l=p_j+p_k$. We refer to this instance as $I'$ throughout this section and denote the
    schedule that results from $\sigma(I)$ by replacing $j$ and $k$ with $l$ by $\sigma(I')$.
    \begin{quote}
        \small\normalfont
        \textbf{Note.}
        We may assume without loss of generality that either
        \(S_j(\sigma(I))=C_k(\sigma(I))\) or
        \(S_k(\sigma(I))=C_j(\sigma(I))\) as swapping the order of two jobs of which neither finishes last on the same machine,
        does not alter the load of any machine and the resulting schedule remains a list schedule.
    \end{quote}
\end{definition}

\begin{lemma}[Unification]\label{corr:unification}
    Consider an instance $I$ of  $P\!\parallel\!\sum p_jC_j^{\alpha}$ where $\alpha \in \Big(0,\frac12\Big]$.
    Then there exists an instance $I'$ of $P\!\parallel\!\sum p_jC_j^{\alpha}$ such that no two jobs
    $j,k$ that are scheduled on the same machine in any cost-maximizing schedule $\sigma^{W}(I')$ are also
    scheduled on the same machine in any optimal schedule $\sigma^{*}(I')$ and vice versa. Furthermore
    \[
    \lambda_{\alpha}(I)\le\lambda_{\alpha}(I')
    \]
    holds.
\end{lemma}

This is essentially the reverse operation compared to the splitting of jobs performed in Lemma \ref{corr:splitting-jobs}.

\begin{proof}[Proof]
    Suppose there exist two jobs $j,k$ that are scheduled on the same machine $i$ in both $\sigma^{W}(I)$ and $\sigma^{*}(I)$.
    We now perform the pairwise unification of $j,k$ in both schedules. Let $I'$ be the instance that emerges from this unification.
    Let $\sigma^{W}(I')$ and $\sigma^{*}(I')$ be the schedules that emerge from $\sigma^{W}(I)$ and $\sigma^{*}(I)$, respectively, after the pairwise unification.
    Then, due to \eqref{eq:queyranne-identity}
    \begin{align*}
        \gamma_{\alpha}(\sigma^{W}(I)) - \gamma_{\alpha}(\sigma^{W}(I'))
        &= (1-2\alpha)p_jp_k \ge 0.
    \end{align*}

    Analogously for $\sigma^{*}(I)$,

    \begin{align*}
        \gamma_{\alpha}(\sigma^{*}(I)) - \gamma_{\alpha}(\sigma^{*}(I'))
        &= (1-2\alpha)p_jp_k \ge 0.
    \end{align*}

    Combining these two results, we have
    \begin{align*}
        \lambda_{\alpha}(I) = \frac{\gamma_{\alpha}(\sigma^{W}(I))}{\gamma_{\alpha}(\sigma^{*}(I))}
        &= \frac{\gamma_{\alpha}(\sigma^{W}(I')) + (1-2\alpha)p_jp_k}{\gamma_{\alpha}(\sigma^{*}(I'))+ (1-2\alpha)p_jp_k } \\
        &\le \frac{\gamma_{\alpha}(\sigma^{W}(I'))}{\gamma_{\alpha}(\sigma^{*}(I'))}
        \le \lambda_{\alpha}(I').
    \end{align*}

    We can repeat this process until no two jobs are scheduled on the same machine in both $\sigma^{W}(I')$ and $\sigma^{*}(I')$.
    We retain dominance in every pairwise unification which is why the resulting instance is $\lambda_{\alpha}$-dominant.
\end{proof}
Since Lemmas \ref{corr:splitting-jobs} and \ref{corr:unification} describe opposing operations, 
each of which strictly increases the competitive ratio on its respective \(\alpha\)-interval (excluding $\alpha=\frac12$), 
they cannot both apply whenever the inequality  
\[
\lambda_{\alpha}(I') \geq \lambda_{\alpha}(I)
\]is strict.
The only value for which both results are correct is \(\alpha=\frac12\).

\begin{definition}[Maximally Unified Instance]\label{def: maximally unified}
    An instance \(I\) is called \emph{maximally unified} 
    if no pairwise unification can be applied to any best–worst schedule pair $\sigma^*(I)/\sigma^W(I)$ without changing 
    the load on any machine in either schedule and the resulting schedules remain list schedules.
\end{definition}

The following two results yield another characterization by bounding the 
number of jobs one needs to consider in the analysis of the approximation ratio. 
The first one follows immediately from Lemma \ref{corr:unification}.

\begin{corollary}[Job bound 1]\label{corr: job bound1}
For every instance $I$ of $\Problem$ with
$\alpha\in\Big(0,\frac12\Big]$, there exists an instance $I'$ of the same
problem such that
\[
    \lambda_\alpha(I')\geq\lambda_\alpha(I),
\]
and every machine contains at most $m$ jobs in every optimal schedule
and every cost-maximizing schedule of $I'$.
\end{corollary}

\begin{proof}
    Suppose the statement is false.
    Then there exists a machine with more than \(m\) jobs assigned to it
    in one of the schedules, say \(\sigma^{W}(I)\). Since \(\sigma^{*}(I)\)
    is a schedule on \(m\) machines, the pigeonhole principle implies that two of
    these jobs must be assigned to the same machine in \(\sigma^{*}(I)\). The
    argument is symmetric if the overloaded machine occurs in \(\sigma^{*}(I)\).    
    Thus, we can perform pairwise unification of these two jobs and obtain an instance $I'$ with $\lambda_{\alpha}(I)\le\lambda_{\alpha}(I')$.
    We repeat this process until no machine has more than \(m\) jobs assigned to it in either schedule.
\end{proof}

It follows from this result, that the subclass of instances with at most \(m\) jobs per machine in any 
optimal or cost-maximizing schedule is \(\lambda_{\alpha}\)-dominant for \(\alpha \in \Big(0, \frac{1}{2}\Big]\).
This of course also gives  a bound on the total number of jobs:
\[
n \le m^2.
\]

We want to further restrict the number of jobs. To this end, it is convenient to consider the problem on a fixed but arbitrary 
number of machines as we can avoid the technicalities of the dominance term and instead work with maximizing instances.
We will later relate this to $\Problem$.

\begin{definition}[$\lambda^{m}_{\alpha}$]\label{def:lambda-m}
Fix any $m \ge 2$ and define
\[
 \lambda^{m}_{\alpha} \coloneqq \max_{\left\{ I \mid I \text{ has } m \text{ machines} \right\}}\lambda_{\alpha}(I).
\]
\newline
$\lambda^{m}_{\alpha}$ is well defined since: 
\begin{itemize}
     \item $m$ is fixed and thus bounded
     \item $n \le m^2$ for the relevant set of instances we maximize over. 
     \item We may assume, without loss of generality, that the job lengths are within a compact set $p_j \in [0,1]$
\end{itemize}
Thus, the problem amounts to maximizing a continuous function over a compact set. For technical convenience, we allow jobs with \(p_j=0\), 
ensuring that the underlying set is closed. Such jobs can be discarded, since removing them does not change the instance.
A instance $I$ of $Pm \!\parallel\! \sum_{j} p_jC^{\alpha}_j$ is called $\lambda^m_{\alpha}$-maximizing, 
if 
\[
\lambda_{\alpha}(I) = \lambda^m_{\alpha}.
\]
\end{definition}

\begin{theorem}[Job bound 2]\label{thm:job-bound}
Let $m \ge 2$ and $I$ be a $\lambda_\alpha^m$-maximizing instance of $Pm\!\parallel\!\sum_j p_j C_j^\alpha$ 
where $\alpha \in \Big(0,\frac12\Big)$. Then
\[
    n \le 3m-2.
\]
\end{theorem}



\begin{proof}
Consider an arbitrary assignment of jobs from an instance $I$ onto its machines. We denote this as a partition of the jobs into $m$
sets $P(I)=(P_1,\ldots,P_m)$.
Write
\[
    L^{i}(P(I)):=\sum_{j\in P_i}p_j,
 \qquad L^{\min}(P(I)):=\min_i L^i(P(I)),
 \qquad p^{\max}_i(P(I)):=\max_{j\in P_i}p_j,
\]
where $p^{\max}_i(P(I)) \coloneqq 0$ if $P_i$ is empty. 
We say that a schedule $\sigma(I)$ agrees with or corresponds to a partition $P(I)$ and vice versa if the assignment of jobs to machines is the same for $P(I)$ and $\sigma(I)$.
We now put a condition on when $P(I)$ corresponds to a list schedule.

$P(I)$ corresponds to a list schedule, if and only if, 

\begin{equation}
    L^{i}(P(I))-p^{\max}_i(P(I)) \le L^{\min}(P(I)).
    \label{eq:list-characterization}
\end{equation}


Indeed, if a list produces $P(I)$, consider the last job assigned to
machine $i$.  Immediately before that assignment, machine $i$ has load
at least $L^{i}(P(I))-p^{\max}_i(P(I))$ and is a least-loaded machine.  A machine of final
load $L^{\min}(P(I))$ has current load at most $L^{\min}(P(I))$, so necessarily $L^{i}(P(I))-p^{\max}_i(P(I)) \le L^{\min}(P(I))$ for
a list schedule.

Conversely, given a partition $P(I)$ such that \eqref{eq:list-characterization} holds,
choose a largest job on every machine and place it last in that
machine's processing order. Order every other job arbitrarily on its
machine. 
The start time of every job is at most
$L^i(P(I))-p^{\max}_i(P(I))\leq L^{\min}(P(I))$. Create a list by ordering the jobs in order
of non-decreasing start time. Break ties arbitrarily. We now schedule according to this list, 
breaking ties by scheduling according to $P(I)$.\medskip 

Whenever a job of prescribed start time $s$ is selected from the list,
every machine with an unassigned job has current load at least $s$, and
every machine that already has all of its jobs assigned has final load at least $ L^{\min}(P(I))\geq s$.  Hence the
designated machine is currently least loaded, and the order is a
list agreeing with the partition $P(I)$.
Let $I$ be an instance of $Pm\!\parallel\!\sum_{j}p_j C^{\alpha}_j$ that attains $\lambda_{\alpha}^m$.

\noindent
Choose a cost-maximizing list partition $P^{W}(I)$ and a minimizing
partition $P^*(I)$. A cost maximizing/minimizing list partition is a partition such that 
there exists a cost maximizing/minimizing list schedule $\sigma^*(I)$/$\sigma^{W}(I)$ that agrees with $P^*(I)$/$P^{W}(I)$.\medskip
 
\noindent
Form a bipartite undirected graph $G$: its right vertices are associated to
the $m$ machines of $P^{*}(I)$, its left vertices are associated with the $m$ machines of
$P^W(I)$, and every job $j$ corresponds to an edge, of weight $p_j$, incident to the two machine 
nodes to which it is assigned in $P^{W}(I)$/$P^{*}(I)$.


\noindent
For a left vertex $v$, write
\[
 w_v:=L^v(P^W(I)),
 \qquad L^{\min}(P^W(I)):=\min_v w_v,
 \qquad p^{\max}_v(P^W(I)):=\max_{e\ni v} p_e.
\]
$p^{\max}_v(P^W(I))$ is of course just the maximal processing time of a job assigned to machine $v$ in $P^{W}(I)$.
We call an edge $e_v$ \emph{tight} if
\begin{enumerate}
     \item $ w_v-p^{\max}_v(P^W(I))=L^{\min}(P^{W}(I))$ \label{eq:critical}, and 
     \item $e_v$ is the unique edge of weight $p^{\max}_v$ incident with $v$
\end{enumerate}


Let $T$ be the set of tight edges.
We show that every cycle of $G$ contains an edge of $T$.  Suppose to the
contrary that a cycle $C$ avoids $T$. Starting arbitrarily, give the successive edges of $C$
alternating signs $d_e\in\{+1,-1\}$ and set $d_e=0$ for any edge not in the cycle.
We perturb the processing times:
\[
 p_e(\varepsilon):=p_e+\varepsilon d_e.
\]
For all sufficiently small $\varepsilon$, every processing time remains
positive. We will denote the instance that results from these perturbations by
$I(\varepsilon)$. At every vertex of the cycle, one incident edge increases
and one decreases.  Hence all loads in both partitions $P^{*}(I(\varepsilon))$ and
$P^W(I(\varepsilon))$ remain unchanged.

The partition $P^W(I(\varepsilon))$ remains list-realizable for both sufficiently
small signs of $\varepsilon$.  This follows from
\eqref{eq:list-characterization}.
If the inequality is tight and the largest incident edge is unique,
that edge is tight, hence it is not on $C$ and remains
unchanged by perturbation. If the maximum $p^{\max}_v(P^W(I))$ is attained by several incident edges, then
either an unchanged maximum remains (if it is not part of the cycle) and thus the maximum does not
fall below $p^{\max}_v(P^W(I))$ after perturbation, which implies that $w_v-p^{\max}_v(P^W(I))\leq L^{\min}(P^{W}(I(\varepsilon)))$ remains true, or the two cycle edges are all the
maximizers. In this case, one of them increases for either non-zero
sign of $t$.  In all cases
$w_v-p^{\max}_v(P^W(I))\leq L^{\min}(P^{W}(I(\varepsilon)))$ is true. 
Thus our perturbed instance and the partitions we chose for $I$ still correspond to list schedules.

Define
\[
Q(I)\coloneqq \sum_{j \in I} p^2_j
\]

This object satisfies
\begin{equation}\label{eq:Q-perturbation}
 Q(I(\varepsilon)) \coloneqq \sum_{j=1}^{n} (p_j(\varepsilon))^2 = \sum_{j=1}^{n} (p_j+\varepsilon d_j)^2 = \sum_{j=1}^{n}p_j^2+ \sum_{j=1}^{n}2\varepsilon p_jd_j + \varepsilon^2 \sum_{j=1}^{n}d_j^2.
\end{equation}
For at least one of the two signs of every sufficiently small non-zero $\varepsilon$, this
is strictly larger than $Q(I)$.

Let
\[
 A^W:=\sum_{i=1}^{m} \left(L^i(P^{W}(I))\right)^2,
 \qquad A^{*}:=\sum_{i=1}^{m} \left(L^i(P^{*}(I))\right)^2.
\]
Since we consider only instances that attain the non-trivial ratio
$\lambda^{m}_{\alpha}>1$, we must have
$A^{W}>A^{*}$.


Consider now a cost maximizing schedule $\sigma^W(I)$ that agrees with the partition $P^W(I)$ 
as well as an optimal schedule $\sigma^{*}(I)$ that agrees with the partition $P^{*}(I)$.
It holds for these that

\[
\lambda_{\alpha}^{m}=
\frac{\gamma_{\alpha}(\sigma^W(I))}{\gamma_{\alpha}(\sigma^{*}(I))}
<
\frac{\gamma_{\alpha}(\sigma^W(I(\varepsilon)))}{\gamma_{\alpha}(\sigma^{*}(I(\varepsilon)))} \le \lambda_{\alpha}^m
\]

Where we choose in \eqref{eq:Q-perturbation} the sign for $d_{e}$ that increases $Q(I(\varepsilon))$.
The inequality follows from an expansion which is carried out in \hyperlink{proof:JobNumber}{Appendix~\ref{app:proof-JobNumber}
(Proof of Theorem~\ref*{thm:job-bound})}.


This contradicts the definition of
$\lambda_{\alpha}^m$.  Therefore every cycle contains a tight
edge, or equivalently,
\begin{equation}\label{eq:forest}
 G-T \qquad \text{is a forest}.
\end{equation}

There is at most one tight edge at each left vertex.  Moreover, a left
vertex of minimum load cannot have a positive tight edge.
\begin{equation}\label{eq:C-bound}
 |T|\leq m-1.
\end{equation}

The bipartite graph has $2m$ vertices.  By
\eqref{eq:forest}, its non-tight part has at most $2m-1$ edges (it is a forest).
Every job is one edge of $G$, so
\[
 n=|E(G-T)|+|T|
 \leq(2m-1)+(m-1)=3m-2.
\]
\end{proof}

For the special case of $\alpha=\frac12$ this statement is incorrect as 
we can always split jobs according to Lemma \ref{corr:splitting-jobs} and remain 
cost equal. 
We thus show that there at least \emph{exists} a $\lambda_\alpha^m$-maximizing 
instance of $Pm\!\parallel\!\sum_j p_j C_j^\alpha$.\medskip

Let $(\alpha_r)_{r\in\mathbb N}$ satisfy
$0<\alpha_r<\frac12$ and $\alpha_r\uparrow\frac12$.
For each $r$, choose a $\lambda_{\alpha_r}^m$-maximizing
instance $I_r$ (by Theorem \ref{thm:job-bound} it has at most $3m-2$ jobs). Normalize
$\sum_j p_j^{(r)}=1$ and allow formally $p_j=0$. This is only done, so that 
a vector of job lengths $I_r=\bigl(p_1^{(r)},\ldots,p_{3m-2}^{(r)}\bigr)$ is an element of a compact set.
If the processing time of any job is in fact $0$, we do not consider it to be part of the instance in the end. 
Because this set is compact, there exists a subsequence
$(I_{r_k})_{k\in\mathbb N}$, with $r_1<r_2<\cdots$, such that
\[
\bigl(p_1^{(r_k)},\ldots,p_{3m-2}^{(r_k)}\bigr)
\xrightarrow[k\to\infty]{}
\bigl(p_1^*,\ldots,p_{3m-2}^*\bigr)\in K.
\]
The limit defines an instance $I^*$ after
removing zero-length jobs. Moreover,
\[
\alpha_{r_k}\xrightarrow[k\to\infty]{}\frac12.
\]
Continuity of $\lambda^m_{\alpha}(I)$ implies that indeed $\lambda^m_{\alpha}(I^*) = \lambda^m_{\alpha}$

Thus, for any $m\ge2$ and $\alpha=\frac12$ there exists a $\lambda^m_{\alpha}$-maximizing
instance $I$ with at most $3m-2$ jobs.  

Our proof was carried out only for instances of $Pm\!\parallel\!\sum_{j}p_j C^{\alpha}_j$.
We now use it to obtain a corresponding
$\lambda_\alpha$-dominance result for $\Problem$.

\begin{corollary}
Let $\alpha\in\Big(0,\frac12\Big]$, and let $I$ be an instance of
$\Problem$ on $m$ machines. Then there exists an instance $I'$ of the
same problem such that
\[
    \lambda_\alpha(I')\geq\lambda_\alpha(I)
\]
and
\[
    |J(I')|\leq3m'-2.
\]
\end{corollary}

\begin{proof}
We can show the existence of this instance by setting $m'=m$.
Now let $I'$ be a $\lambda_\alpha^{m'}$-maximizing instance with $|J(I')| \le 3m-2$. By the
definition of $\lambda_\alpha^{m'}$,
\[
    \lambda_\alpha(I')
    =\lambda_\alpha^{m'}
    \geq\lambda_\alpha(I).
\]
Since $I'$ is $\lambda_\alpha^{m'}$-maximizing,
Theorem~\ref{thm:job-bound} implies that
\[
    |I'|\leq3m'-2,
\]
where $m'$ is the number of machines of instance $I'$.
\end{proof}


\subsection*{Singleton}

We know from Lemma~\ref{corr:unification} that both the total number of jobs as well as the number of jobs assigned
to any individual machine are bounded. Next, we show that in order to determine $\lambda_{\alpha}$, it is sufficient to
consider optimal schedules with a singleton.

\begin{definition}[Singleton]
Let \(I\) be an instance of \(P\!\parallel\!\sum_{j=1}^{n}p_j C_j^{\alpha}\).
We call a job $j$ of $I$ a singleton in a schedule $\sigma(I)$, if it is the only job scheduled on 
some machine in $\sigma(I)$.
\end{definition}
We will need two auxiliary lemmas to show the singleton result.
Both are purely algebraic.

\begin{lemma}\label{lem:peak-increase}
Let $x_1,\dots,x_n \in \mathbb{R}_{+}$ 
and suppose that $x_1 \ge x_i$ for all $i$.
Furthermore, let $x_1', \dots, x_n' \in \mathbb{R}_{+}$ such that
$\delta\geq0$ and 
$\delta_2,\dots,\delta_n \ge 0$ exist that satisfy
\[
\sum_{i=2}^n \delta_i = \delta.
\]
and
\[
x_1' = x_1+\delta,
\qquad
x_i' = x_i-\delta_i \quad \text{for } i=2,\dots,n.
\]

Then
\[
\sum_{i=1}^n (x_i')^2 \geq \sum_{i=1}^n x_i^2.
\]
If \(\delta>0\), then the inequality is strict.
\end{lemma}

\begin{proof}
We compute
\begin{align*}
\sum_{i=1}^n (x_i')^2-\sum_{i=1}^n x_i^2
&=
(x_1+\delta)^2-x_1^2
+
\sum_{i=2}^n \left((x_i-\delta_i)^2-x_i^2\right)\\
&=
2\delta x_1+\delta^2
-
2\sum_{i=2}^n \delta_i x_i
+
\sum_{i=2}^n \delta_i^2.\\
&=
2\sum_{i=2}^n \delta_i(x_1-x_i)
+
\delta^2
+
\sum_{i=2}^n \delta_i^2.
\end{align*}

All terms in the last expression are nonnegative. Hence,
\[
\sum_{i=1}^n (x_i')^2-\sum_{i=1}^n x_i^2 \geq 0.
\]
If \(\delta>0\), then \(\delta^2>0\), so the inequality is strict.

\end{proof}

\begin{lemma}\label{lem:ratio-monotonicity}
Let
\[
A\geq C>B\ge 0,\qquad r\geq0,\qquad d>0,\qquad B+d\leq C,
\qquad \alpha \in \Big(0, \frac12\Big].
\]
Then
\[
\frac{A+(2\alpha-1)B}{C+(2\alpha-1)B}
\le
\frac{A+r+(2\alpha-1)(B+d)}{C+(2\alpha-1)(B+d)}.
\]
The inequality is strict if either $r>0$, or
\(\alpha<\frac12\) and $A>C$.
\end{lemma}

\begin{proof}

Let
\[
x=1-2\alpha.
\]
Since \(\alpha\in\Big(0,\frac12\Big]\), we have \(x\in[0,1)\). The inequality becomes
\[
\frac{A-xB}{C-xB}
\le
\frac{A+r-x(B+d)}{C-x(B+d)}.
\]

Both denominators are positive. Equivalently:
\[
(A-xB)\bigl(C-x(B+d)\bigr)
\le
\bigl(A+r-x(B+d)\bigr)(C-xB).
\]
After cancellation, this is equivalent to
\[
r(C-xB)+x(A-C)d\ge 0.
\]

Since \(x\in[0,1)\) and \(C>B\geq0\), we have \(C-xB>0\). Also,
\[
    C-x(B+d)\geq C-xC=(1-x)C>0.
\]
Moreover, \(r\geq0\), \(x\geq0\), \(A-C\geq0\), and \(d>0\).
Thus both terms in the last expression are nonnegative, which proves the
inequality.
If \(r>0\), or \(\alpha<\frac12\) and \(A>C\), the inequality is strict.
\end{proof}

\begin{theorem}[Singleton]\label{corr: singletonjob}
Let $I$ be an instance of $\Problem$ and $0 < \alpha \leq \frac{1}{2}$.
Then there exists an instance \(I'\) of $\Problem$, a job $j'$ of $I'$, an
optimal schedule \(\sigma^*(I')\), and a cost-maximizing schedule \(\sigma^W(I')\)
such that:
\begin{enumerate}
    \item \(j'\) is a singleton in \(\sigma^*(I')\)
    \item no job finishes after \(j'\) in \(\sigma^W(I')\)
    \item $\lambda_\alpha(I)\leq\lambda_\alpha(I')$.
\end{enumerate}
\end{theorem}

\begin{proof}
We first describe one step of the reduction and show that it does not decrease
$\lambda_{\alpha}(I)$. Repeating this step will produce an instance with all
the claimed properties. By Lemma~\ref{corr:unification}, we may assume before
each step that the current instance is maximally unified
(Definition \ref{def: maximally unified}).\medskip

Consider any maximally-unified (Definition \ref{def: maximally unified}) instance $I$ of $P\!\parallel\!\sum_{j=1}^{n}p_j C_j^{\alpha}$ and fix any pair $\sigma^{W}(I),\sigma^{*}(I)$. Fix any job $k$ that finishes last in $\sigma^{W}(I)$.
If $k$ is already a singleton in $\sigma^*(I)$, the statement is true
with $I'=I$ and $j'=k$. Hence, assume that $k$ is not a singleton in
$\sigma^*(I)$.
Let $i^k$ be the machine on which $k$ is scheduled in $\sigma^{*}(I)$.
Apply pairwise unification (Definition \ref{def:pairwise-unification}) repeatedly to unify all jobs scheduled on $i^k$ in $\sigma^{*}(I)$ into a single job $j'$.
The first observation is that:
\[
L^{i}(\sigma^{*}(I)) = L^{i}(\sigma(I')) \quad \forall i \in [m].
\]
where $I'$ is the instance that differs from $I$ only through unification to $j'$.
And $\sigma(I')$ is the schedule that differs from $\sigma^{*}(I)$ only through unification to $j'$.
It is immediately clear that $\sigma(I')$ is indeed optimal.
It is therefore justified to denote $\sigma(I')$ by $\sigma^*(I')$. \medskip

\noindent
Next we consider what happens to the schedule $\sigma^{W}(I)$ when we unify the jobs on $i^k$ to become the singleton $j'$.
To this end, let $i^{w}$ be the machine on which $k$ is scheduled in $\sigma^{W}(I)$.

Starting with $\sigma^{W}(I)$ we perform the following steps of construction:
\begin{construction}\label{building SW}
\leavevmode
\begin{enumerate}
    \item Move all jobs that were affected by unification except $k$ to be scheduled after $k$ on $i^{w}$ in $\sigma^{W}(I)$. The resulting schedule is shown in Figure \ref{fig:schedule-structure} (c).
     These can now be unified with $k$ to become $j'$.
    \item On each machine, remove all idle-time gaps by shifting jobs to start earlier without changing their order. The resulting schedule is shown in Figure \ref{fig:schedule-structure} (e).
    \item Reschedule the subset of jobs that are not scheduled according to list scheduling. Do this by going through the jobs in order of
     non-decreasing start time, rescheduling a job $j$ if and only if $S_j(\sigma(I))>L^{\min}(\sigma^{j}(I))$. $\sigma^{j}(I)$ denotes the intermediate schedule 
     right before rescheduling $j$. The resulting schedule is shown in Figure \ref{fig:schedule-structure} (e).
    \item The resulting schedule is named $\sigma^{C}(I')$. The resulting schedule is shown in Figure \ref{fig:schedule-structure} (g).
\end{enumerate}
\end{construction}


We do not need to show that
$\gamma_{\alpha}(\sigma^C(I'))=\gamma^W_{\alpha}(I')$. We only show that
$\sigma^C(I')$ is a list schedule whose cost is sufficient to establish
\[
\lambda_{\alpha}(I') \ge \frac{\gamma_{\alpha}(\sigma^C(I'))}{\gamma_{\alpha}^*(I')}\geq \frac{\gamma^W_{\alpha}(I)}{\gamma^*_{\alpha}(I)}= \lambda_{\alpha}(I).
\]
The last-finishing-job property will follow by iterating the reduction.
\begin{figure}[p]
\centering
\setlength{\fboxsep}{6pt}
\fbox{%
\begin{minipage}{0.94\textwidth}
    \begin{minipage}[t]{0.47\linewidth}
        \centering
        \resizebox{\linewidth}{!}{\schedulebeforeunification}
        \par\smallskip
        \small(a) \(\sigma^{W}(I)\)
    \end{minipage}
    \hfill
    \begin{minipage}[t]{0.47\linewidth}
        \centering
        \resizebox{\linewidth}{!}{\scheduleafterunification}
        \par\smallskip
        \small(b) \(\sigma^{*}(I)\)
    \end{minipage}
    \par\smallskip
    \centering\small Initial schedules of $I$.
\end{minipage}%
}

\medskip

\fbox{%
\begin{minipage}{0.94\textwidth}
    \begin{minipage}[t]{0.47\linewidth}
        \centering
        \resizebox{\linewidth}{!}{\scheduleclosetimegaps}
        \par\smallskip
        \small(c) Moving jobs from $M_{i^k}$ to $M_{i^{w}}$
    \end{minipage}
    \hfill
    \begin{minipage}[t]{0.47\linewidth}
        \centering
        \resizebox{\linewidth}{!}{\schedulereschedulelastjobs}
        \par\smallskip
        \small(d) Unification
    \end{minipage}
    \par\smallskip
    \centering\small  Step 1 in Construction \ref{building SW}.
\end{minipage}%
}

\medskip

\fbox{%
\begin{minipage}{0.94\textwidth}
    \begin{minipage}[t]{0.47\linewidth}
        \centering
        \resizebox{\linewidth}{!}{\schedulelistschedule}
        \par\smallskip
        \small(e) Removing idle time
    \end{minipage}
    \hfill
    \begin{minipage}[t]{0.47\linewidth}
        \centering
        \resizebox{\linewidth}{!}{\schedulefinalworst}
        \par\smallskip
        \small(f) \(\sigma^{*}(I')\)
    \end{minipage}
    \par\smallskip
    \centering\small  Step 2 in Construction \ref{building SW}.
\end{minipage}%
}

\medskip

\fbox{%
\begin{minipage}{0.94\textwidth}
    \begin{minipage}[t]{0.47\linewidth}
        \centering
        \resizebox{\linewidth}{!}{\scheduleseventh}
        \par\smallskip
        \small(g) \(\sigma^{C}(I')\)
    \end{minipage}
    \hfill
    \begin{minipage}[t]{0.47\linewidth}
        \centering
        \resizebox{\linewidth}{!}{\scheduleeighth}
        \par\smallskip
        \small(h) \(\sigma^{*}(I')\)
    \end{minipage}
    \par\smallskip
    \centering\small Step 3 in Construction \ref{building SW}.
\end{minipage}%
}
\caption{Singleton schedule construction.}
\label{fig:schedule-structure}
\end{figure}


$\sigma^{C}(I')$ is indeed a list schedule, since we reschedule all jobs that might not be scheduled
according to a list. It is therefore true by construction.\medskip

Observe that each machine "contributes" at most one job to the unification
in $\sigma^{W}(I)$. If this were not the case, another pairwise unification (Definition \ref{def:pairwise-unification}) would have been possible contradicting the fact that $I$ is maximally unified.
Next, we reason for:
\[
L^{\max}(\sigma^{C}(I')) \geq L^{\max}(\sigma^{W}(I))
\]

We consider two cases.
\begin{enumerate}
    \item   $j'$ is \textbf{not rescheduled} in Construction \ref{building SW}. As $k$ was a job that completes last in $\sigma^{W}(I)$ and no other machine's load
            is increased above $L^{\max}(\sigma^{W}(I))$ in the construction process, the job $j'$ will have the same starting time as $k$ and 
            has processing time $p_{j'} > p_k$ which implies $C_{j'}(\sigma^{C}(I')) > C_{k}(\sigma^{W}(I))$. Therefore there is no job that completes after $j'$ in $\sigma^{C}(I')$.
            Formally, $L^{\max}(\sigma^{C}(I')) = C_{j'}(\sigma^{C}(I'))$.
            Taking all observations into consideration, we get 
            \begin{align*}
            L^{\max}(\sigma^{C}(I'))=C_{j'}(\sigma^{C}(I')) > C_{k}(\sigma^{W}(I))=L^{\max}(\sigma^{W}(I)).
            \end{align*}
    \item   Without loss of generality, $j'$ \textbf{is rescheduled} to machine $m$ in Construction \ref{building SW}
            (as depicted in (e) to (g) of Figure \ref{fig:schedule-structure}). Since $j'$ was not scheduled on $m$ in $\sigma^{W}(I)$, it must hold that $S_{j'}(\sigma^{W}(I))\le L^{m}(\sigma^{W}(I))$. 
            The rescheduling of $j'$ to $m$ implies that a job $v$
            that was scheduled on machine $m$ in $\sigma^{W}(I)$ was unified with $k$ to $j'$. This means that the new load on machine $m$ in $\sigma^C(I')$ is at least 
            as large as the load on machine $i^{w}$ in $\sigma^{W}(I)$.Thus, 
            \begin{align*}
            L^{\max}(\sigma^{C}(I')) = L^{m}(\sigma^{W}(I))-p_{v} +p_{j'} \ge L^{m}(\sigma^{W}(I)) + p_{k} \ge L^{i_w}(\sigma^{W}(I)) = L^{\max}(\sigma^{W}(I)).
            \end{align*}
\end{enumerate}


For any given instance, $\sum_{j=1}^{n}p_j$ is fixed. Therefore, $\sum_{i=1}^{m}L^{i}(\sigma(I)) = \sum_{i=1}^{m}L^{i}(\sigma'(I))$ for any two list schedules $\sigma(I)$ and $\sigma'(I)$.
This implies that the sum of machine loads is also fixed for any instance. Assume, without loss of generality, 
that machine $1$ has maximal load of all machines in both $\sigma^{C}(I')$ and $\sigma^{W}(I)$. 

\noindent
Define

\[x_1 \coloneqq L^{\max}(\sigma^{W}(I)) \qquad x'_1 \coloneqq L^{\max}(\sigma^{C}(I'))\]

\noindent
From the preceding argument, we have $0 \le x_1'- x_1 \eqqcolon \delta$.

\noindent
Define
\[
x_i \coloneqq L^{i}(\sigma^{W}(I)) \qquad \forall i \in \{2,\dots, m\} 
\]
Define $l_i$ to be the job that is scheduled last on machine $i$ in $\sigma^{W}(I)$. 
\[
x_i' \coloneqq L^{\overline{i}}(\sigma^{C}(I')) \qquad \forall i \in \{2,\dots, m\}
\]
where $\overline{i}$ is the index of the machine in $\sigma^{C}(I')$ on which $l_i$ is scheduled. 

\noindent
Then, 
\[
\delta_i \coloneqq  x_i-x_i' \ge 0 \qquad \forall i \in \{2,\dots, m\}
\]
and 
\begin{align*}
\sum_{i=2}^{m} \delta_i &= \sum_{i=2}^{m} (x_i-x_i') + x_1-x_1' - (x_1-x_1') \\
&= \sum_{i=1}^{m}L^{i}(\sigma^{W}(I)) - \sum_{i=1}^{m}L^{i}(\sigma^{C}(I')) -(x_1-x_1')\\
&= x_1'-x_1 = \delta.
\end{align*}
The third equality holds because the load is equal for both instances $I'$ and $I$.
We can now apply Lemma~\ref{lem:peak-increase} to conclude that 
\[
r:= \sum_{i=1}^m (L^{i}(\sigma^{C}(I')))^2 - \sum_{i=1}^m (L^{i}(\sigma^{W}(I)))^2 \ge 0
\]

\noindent
If we additionally define the values

\[
A:= \sum_{i=1}^m (L^{i}(\sigma^{W}(I)))^2 \qquad B:=\sum_{i=1}^{n}p_i^2 \qquad C:=\sum_{i=1}^m (L^{i}(\sigma^{*}(I)))^2
\]
as well as
\[
d
:=
\left(p_k+\sum_{i\in U}p_i\right)^2
-p_k^2-\sum_{i\in U}p_i^2
=2\sum_{\{i,j\}\subseteq U\cup\{k\}}p_ip_j > 0 ,
\]
where $U$ is the set of jobs that are unified with $k$ to form $j'$ in
Construction~\ref{building SW}. The inequality is strict since we assume that $k$ is not a singleton in $\sigma^*(I)$ and thus $U$ is nonempty. Thus $B+d$ is the sum of squared processing
times in $I'$.
We can apply Lemma~\ref{lem:ratio-monotonicity} to get

\[
\lambda_{\alpha}(I') \ge \frac{A+r+(2\alpha-1)(B+d)}{C+(2\alpha-1)(B+d)} \ge \frac{A+(2\alpha-1)B}{C+(2\alpha-1)B} = \lambda_{\alpha}(I)
\]
Thus this procedure, which marks one step in our reduction, 
does not decrease the approximation ratio.\medskip

Because $k$ is not a singleton in $\sigma^*(I)$, at least two jobs are
unified, and hence $|I'|<|I|$.\medskip

We can now repeat the reduction step on the newly created instance $I'$.
$I'$ is now in the role of $I$ in the previous step.\medskip

Choose a cost-maximizing schedule $\sigma^W(I')$ and the 
optimal schedule $\sigma^*(I')$ we just constructed. Choose a job that finishes 
last in the cost-maximizing schedule.
If this job is a singleton in the optimal schedule, we are done. Otherwise,
apply the same reduction described in Construction~\ref{building SW}.\medskip

Every step reduces the number of jobs and since we know that $n\le m^2$ the procedure must end
after finitely many steps.

Let $\widehat I$ be the final instance and let $\widehat j$ be a job that
finishes last in the cost-maximizing schedule selected at the last step.
$\widehat j$ is a singleton in the selected optimal
schedule of $\widehat I$. By construction, no job finishes after $\widehat j$ in
the selected cost-maximizing schedule. Moreover, chaining the inequalities
from all reduction steps gives
\[
    \lambda_\alpha(I)\leq\lambda_\alpha(\widehat I).
\]
Finally, taking $I':=\widehat I$ and $j':=\widehat j$ proves the theorem.
\end{proof}


\noindent
\subsection*{Perturbing job length}
We can further characterize singleton jobs in $\lambda_{\alpha}$-dominant instances.
The results in this section rely on perturbing the processing time of jobs.
\begin{definition}[Single-job processing-time perturbation]\label{def:singleojobperturbation}
Let $I$ be an instance of
$P \!\parallel\! \sum_j p_jC_j^{\alpha}$, and let $j^*$ be a job with
processing time $p_{j^*}=p $. For $\varepsilon \in \mathbb{R}$, define
\(I_{j^*}(\varepsilon)\) to be the instance obtained from \(I\) by perturbing
only the processing time of \(j^*\) by \(\varepsilon\). Formally,
\[
p_j(\varepsilon)
:=
\begin{cases}
p+\varepsilon, & \text{if } j=j^*,\\
p_j,           & \text{if } j\neq j^*.
\end{cases}
\]
We call $I_{j^*}(\varepsilon)$ the $\varepsilon$-perturbation of $I$ at $j^*$.
\end{definition}

\begin{corollary}[The critical singleton is not a singleton in the worst schedule]
\label{corr:singleton-not-worst}
Let $m\geq2$, let $\alpha\in\Big(0,\frac12\Big]$, and let $I$ be a
$\lambda_\alpha^m$-maximizing instance. Let $j$ be a job with the
properties of the singleton constructed in
Theorem~\ref{corr: singletonjob}. Then $j$ is not a singleton in some cost maximizing schedule
$\sigma^W(I)$.
\end{corollary}

\begin{proof}
By Theorem~\ref{corr: singletonjob}, we know that there must exist a job \(j\) that is a singleton in
some optimal schedule \(\sigma^{*}(I)\), and no job finishes after \(j\) in the worst schedule constructed in Theorem \ref{corr: singletonjob} which we call
\(\sigma^{W}(I)\).

Suppose, for contradiction, that \(j\) is also a singleton in
\(\sigma^W(I)\), and write \(p:=p_{j}\). Since no job finishes
after \(j\), every job appearing after \(j\) in a list
producing \(\sigma^W(I)\) starts strictly before \(p\). Hence, for
sufficiently small \(\varepsilon>0\), decreasing \(p_{j}\) from
\(p\) to \(p-\varepsilon\) leaves the assignment produced by this list
unchanged. We will thus work with the $(-\varepsilon)$-perturbation of $I$ at $j$: \(I_{j}(-\varepsilon)\).

Since \(j\) is a singleton in both schedules, using the expression determined in 
Lemma~\ref{corr:queyranne} shows that both schedule costs decrease
by
\[
x
:=
\alpha\bigl(p^2-(p-\varepsilon)^2\bigr)>0.
\]
Therefore,
\[
\lambda_\alpha(I_{j}(-\varepsilon))
\ge
\frac{\gamma_\alpha^W(I)-x}
     {\gamma_\alpha^*(I)-x}.
\]
Using the same subtraction argument as in the proof of
Lemma~\ref{corr:splitting-jobs} and
\(\gamma_\alpha^W(I)>\gamma_\alpha^*(I)\), we obtain
\[
\frac{\gamma_\alpha^W(I)-x}
     {\gamma_\alpha^*(I)-x}
>
\frac{\gamma_\alpha^W(I)}
     {\gamma_\alpha^*(I)}
=
\lambda_\alpha(I).
\]
Thus,
\[
    \lambda_\alpha\bigl(I_j(-\varepsilon)\bigr)
    >
    \lambda_\alpha(I)
    =
    \lambda_\alpha^m,
\]
contradicting the definition of $\lambda_\alpha^m$. Hence $j$ is not
a singleton in $\sigma^W(I)$.
\end{proof}

Since Corollary~\ref{corr:singleton-not-worst} is not stated directly
in terms of \(\lambda_\alpha\)-dominance, we now make this connection
explicit. Together with the preceding structural results, the corollary
implies that any instance \(I\) that does not admit a pair consisting of
an optimal schedule and a cost-maximizing schedule satisfying all the
structural properties established so far in this chapter is not
\(\lambda_\alpha^m\)-maximizing.

Indeed, there exists a \(\lambda_\alpha^m\)-maximizing instance \(I'\)
that admits such a pair of schedules. This follows from the fact that a $\lambda^m_{\alpha}$-maximizing instance
exists and that the properties we impose on it imply dominance. Therefore,
\[
    \lambda_\alpha(I')
    =
    \lambda_\alpha^m
    \geq
    \lambda_\alpha(I).
\]
Consequently, the subclass of instances admitting such a pair of schedules
is \(\lambda_\alpha\)-dominant. 

\begin{corollary}[Jobs larger than a singleton]\label{cor: larger than singleton}
Suppose $j$ is a singleton in $\sigma^{*}(I)$ of $\Problem$.
Then every job $k$ with
\[
    p_k>p_j,
\]
is also a singleton in $\sigma^{*}(I)$.
Consequently, if $j^*$ is the only singleton in $\sigma^{*}(I)$,
then $j^*$ is a largest job.
\end{corollary}

\begin{proof}
Let $M_1$ be the machine on which the singleton $j$ is scheduled in $\sigma^{*}(I)$. Define $p \coloneqq p_j$ so that 
the load of $M_1$ in $\sigma^{*}(I)$ is $p$.  Consider any job $k$ with processing time
\[
    p_k>p,
\]
and let $M_i$ be the machine on which job $k$ is scheduled.
If such a job does not exist we are done. 
Else, exchange $j$ and $k$ and set
\[
    \delta:=p_k-p>0.
\]

Before the exchange, the two machine loads are $(p,L^{i}(\sigma^{*}(I)))$.  After the
exchange, they are $(p_k,L^{i}(\sigma^{*}(I))-\delta)$.
All other machine loads remain unchanged.  Hence, using the term introduced in Lemma \ref{corr:queyranne}
it is easy to see that the change in cost is
\begin{align*}
    \frac12\left(p_k^2+(L^{i}(\sigma^{*}(I))-\delta)^2-p^2-L^{i}(\sigma^{*}(I))^2\right)
      &=\delta(p_k-L^{i}(\sigma^{*}(I))).
\end{align*}
Because $L^{i}(\sigma^{*}(I))\geq p_k$, this quantity is nonpositive.  If $M_i$
contains any job besides $k$, then all processing times are positive
and the exchange would strictly decrease the cost of the schedule, which
contradicts its optimality.
Note that the exchange might not result in a list schedule.
But since the optimal schedule is a list schedule, we can infer that the 
cost would only decrease further.
Equivalently, at least one of the longest jobs must be scheduled alone.
Finally, every job larger than $p$ must also be a singleton in $\sigma^{*}(I)$.
\end{proof}


\begin{lemma}[Singleton perturbation]
\label{lem:critical-singleton-perturbation}
Let $I$ be an instance of $\Problem$ with $\alpha\in \Big(0,\tfrac12\Big]$. Suppose there exists a 
job $j^*$ such that
\begin{enumerate}
    \item $j^*$ is a singleton in some optimal schedule $\sigma^{*}(I)$ 
    \item $j^*$ is not a singleton in some cost-maximizing list schedule $\sigma^{W}(I)$
    \item $j^*$ is scheduled last on the machine which it is assigned to in $\sigma^{W}(I)$.
\end{enumerate}

Then for every $\varepsilon\geq0$
\begin{equation}
    \lambda_{\alpha}(I(\varepsilon))
    \geq
    \frac{
        \WorstCost+(S_{j^*}(\sigma(I))+2\alpha p)\varepsilon
        +\alpha\varepsilon^2
    }{
        \OptCost+2\alpha p\varepsilon
        +\alpha\varepsilon^2
    },
    \label{eq:ratio-perturbation-bound}
\end{equation}
where $p\coloneqq p_{j^*}$
\end{lemma}

\begin{proof}
For an arbitrary list schedule $\sigma(I)$ in which a job $j$ is scheduled last on
a machine, increasing its processing time $p$ to
$p+\varepsilon$ where $\varepsilon \ge 0$, changes the assignment's cost by
\begin{align}
    \Delta(\sigma(I),\varepsilon)
      &=
      \frac12\bigl((p+\varepsilon+S_{j}(\sigma(I)))^2-(p+S_{j}(\sigma(I)))^2\bigr)
      +
      \left(\alpha-\frac12\right)
      \bigl((p+\varepsilon)^2-p^2\bigr)
      \notag\\
      &=(S_{j}(\sigma(I))+2\alpha p)\varepsilon+\alpha\varepsilon^2.
    \label{eq:fixed-assignment-cost-change}
\end{align}

Consider the optimal schedule $\sigma^*(I)$ in which $j^*$ is
a singleton. This assignment is feasible for $I_{j^*}(\varepsilon)$.
As $j^*$ is scheduled alone in an optimal schedule, $S_{j^*}(\sigma^{*}(I))=0$. An application of \eqref{eq:fixed-assignment-cost-change} shows that
the cost of the perturbed instance is
\[
    \OptCost+2\alpha p\varepsilon+\alpha\varepsilon^2.
\]
Therefore,
\begin{equation}
    \gamma^{*}_{\alpha}(I(\varepsilon))
    \leq
    \OptCost+2\alpha p\varepsilon+\alpha\varepsilon^2.
    \label{eq:opt-candidate-bound}
\end{equation}

Next, fix a list corresponding to the cost-maximizing schedule $\sigma^W(I)$ from Theorem \ref{corr: singletonjob}.
Increasing the processing time of $j^*$ leaves all assignment
decisions before $j^*$ unchanged.  After $j^*$ is assigned,
only the availability time of its machine increases.  Since
$j^*$ is last on that machine, every later job in the fixed list
was assigned to another machine. Making the machine of $j^*$
available later cannot change any of those choices.  Thus the same
list produces the same assignment for $I_{j^*}(\varepsilon)$.

Taking $S_j(\sigma(I))=S_{j^*}(\sigma^{W}(I))$ in \eqref{eq:fixed-assignment-cost-change} gives the
cost
\[
    \WorstCost+(S_{j^*}(\sigma^{W}(I))+2\alpha p)\varepsilon
    +\alpha\varepsilon^2.
\]
It remains a list schedule for $I_{j^*}(\varepsilon)$.
Consequently,
\begin{equation}
    \gamma^{W}_{\alpha}(I_{j^*}(\varepsilon))
    \geq
    \WorstCost+(S_{j^*}(\sigma^{W}(I))+2\alpha p)\varepsilon
    +\alpha\varepsilon^2.
    \label{eq:worst-candidate-bound}
\end{equation}

Dividing \eqref{eq:worst-candidate-bound} by
\eqref{eq:opt-candidate-bound} proves
\eqref{eq:ratio-perturbation-bound}.
\end{proof}


\begin{corollary}[Strict improvement when the accompanying load is large]
\label{thm:strict-perturbation-improvement}
Let $I$ be an instance of $\Problem$ and $\alpha\in\Big(0,\tfrac12\Big]$. Suppose that
\begin{equation}
    \LamI\leq 1+\frac{1}{6\alpha}
    \label{eq:rho-upper-bound}
\end{equation}
Let $I$ and $j^*$ satisfy the
assumptions of
Lemma~\ref{lem:critical-singleton-perturbation}. 
If
\[
    S_{j^*}(\sigma^{W}(I))>\frac{p_{j^*}}{3},
\]
then
\begin{equation}
    \boxed{\lambda_{\alpha}(I_{j^*}(\varepsilon))>\LamI}
    \label{eq:strict-ratio-improvement}
\end{equation}
for every
\begin{equation}
    0<\varepsilon<
    6\left(S_{j^*}(\sigma^{W}(I))-\frac p3\right).
    \label{eq:epsilon-improvement-range}
\end{equation}
\end{corollary}

\begin{proof}
Fix $\alpha$ and write $\lambda:=\LamI$.  By
Lemma~\ref{lem:critical-singleton-perturbation}, it is enough to show
\[
    \frac{
        \WorstCost+(S_{j^*}(\sigma^{W}(I))+2\alpha p)\varepsilon
        +\alpha\varepsilon^2
    }{
        \OptCost+2\alpha p\varepsilon
        +\alpha\varepsilon^2
    }
    >\lambda.
\]
Substituting $\WorstCost=\lambda\OptCost$ and canceling,
shows that this strict inequality is equivalent to
\begin{equation}
    S_{j^*}(\sigma^{W}(I))>
    (\lambda-1)(2\alpha p+\alpha\varepsilon).
    \label{eq:exact-improvement-condition}
\end{equation}
\eqref{eq:rho-upper-bound} gives
\[
    (\lambda-1)(2\alpha p+\alpha\varepsilon)
    \leq
    \frac{1}{6\alpha}
    (2\alpha p+\alpha\varepsilon)
    =
    \frac p3+\frac{\varepsilon}{6}.
\]
If \eqref{eq:epsilon-improvement-range} holds, then
\[
    \frac p3+\frac{\varepsilon}{6}<S_{j^*}(\sigma^{W}(I)).
\]
Thus \eqref{eq:exact-improvement-condition} holds, and
\eqref{eq:ratio-perturbation-bound} yields
\eqref{eq:strict-ratio-improvement}.
\end{proof}

In \Cref{sec:Upper Bound}, we prove that the bound assumed in
\Cref{thm:strict-perturbation-improvement} holds for every instance
\(I\) of \(\Problem\) with \(\alpha\in\Big(0,\frac12\Big]\). This implies 
that there exists a non-empty $\lambda_{\alpha}$-dominant subclass for which this property holds.\medskip

\begin{corollary}\label{corr: makespandbound}
Let $\alpha\in\Big(0,\frac12\Big)$, and let $I$ be an instance of $\Problem$
on $m$ machines. Then there exists an instance $I'$ of the same
problem on $m$ machines and a cost-maximizing schedule $\sigma^W(I')$ satisfying
\[
    L^{\max}\bigl(\sigma^W(I')\bigr)
    \leq\frac43p_{\max},
    \qquad
    p_{\max}:=\max_{j\in I'}p_j,
\]
such that
\[
    \lambda_\alpha(I')\geq\lambda_\alpha(I)
\]

\end{corollary}

\begin{proof}
Apply Theorem~\ref{corr: singletonjob} to a
$\lambda_\alpha^m$-maximizing instance and denote the resulting
instance by $I'$. Then there exists a job $j^*$ that is a singleton
in an optimal schedule of $I'$ and after which no job finishes in a
cost-maximizing schedule $\sigma^W(I')$. Moreover,
\[
    \lambda_\alpha(I')
    \geq\lambda_\alpha^m
    \geq\lambda_\alpha(I).
\]
Since $I'$ has $m$ machines, it follows that
$\lambda_\alpha(I')=\lambda_\alpha^m$.

If $j^*$ is also a singleton in $\sigma^W(I')$, the claimed bound is
immediate. Otherwise,
Corollary~\ref{thm:strict-perturbation-improvement} and the maximality
of $I'$ imply
\[
    S_{j^*}\bigl(\sigma^W(I')\bigr)
    \leq\frac{p_{j^*}}3.
\]
Therefore,
\[
    L^{\max}\bigl(\sigma^W(I')\bigr)
    \leq\frac43p_{j^*}
    \leq\frac43p_{\max}.
\]
\end{proof}
\newpage
\begin{framed}
Combining the results from Section \ref{sec:Struct-Results}, we conclude that for
$\alpha\in\Big(0,\frac12\Big]$, the class of instances $I$ for which there
exist an optimal schedule $\sigma^*(I)$, a cost-maximizing schedule
$\sigma^W(I)$, and a job $j^*$ satisfying the following properties:

\begin{enumerate}[label=(\roman*)]
    \item \(w_j=p_j\) for every job \(j\)

    \item There exist a cost-maximizing list schedule
      \(\sigma^{W}(I)\) and an optimal schedule
      \(\sigma^{*}(I)\) such that:
      \begin{enumerate}[label=(\alph*)]
          \item No two jobs are assigned to the same machine
                in both schedules

          \item Each machine contains at most \(m\) jobs
                in either schedule, and
                \[
                    n \leq 3m-2
                \]

          \item There exists a job \(j^*\) that is a singleton
                in \(\sigma^{*}(I)\) and finishes last
                in \(\sigma^{W}(I)\), possibly tied with
                other jobs

          \item The job \(j^*\) is not a singleton
                in \(\sigma^{W}(I)\)

          \item Every job larger than \(j^*\) is a singleton
                in \(\sigma^{*}(I)\). In particular, at least
                one largest job is a singleton in
                \(\sigma^{*}(I)\)

          \item With \(p_{\max}=\max_j p_j\),
                \[
                    L^{\max}\bigl(\sigma^{W}(I)\bigr)
                    \leq \frac43 p_{\max}
                \]
      \end{enumerate}
\end{enumerate}
is $\lambda_\alpha$-dominant.
\end{framed}

\begin{proof}[Proof of Theorem \ref{thm: secondframe}]
    We need to show that this class is nonempty and $\lambda_{\alpha}$-dominant.
    It is nonempty, because as in the proof of Theorem \ref{th: firstframed}, we can apply the results obtained in \Cref{chap:alpha_all}
    and \Cref{chap:alpha_l} in succession. 
    In particular, we apply the construction steps described in the respective proofs of Lemma \ref{lem: lemm1}, Corollary \ref{corr:unification}, Corollary \ref{corr: job bound1} and 
    Theorem \ref{thm:job-bound} which can hold simultaneously since we do not require the existence of anything that is not needed for the other results respectively. 
    After this, we apply the construction steps from the proof of Theorem \ref{corr: singletonjob}. The structural result is now only 
    provably true for a possibly unique pair of schedules $\sigma^W(I)$ and $\sigma^*(I)$. 
    But since all other relevant results, namely Corollary \ref{corr:singleton-not-worst}, Corollary \ref{cor: larger than singleton}, Corollary \ref{thm:strict-perturbation-improvement}, Lemma \ref{lem:critical-singleton-perturbation},
    and Corollary \ref{corr: makespandbound} extend on exactly this pair of 
    schedules, the set of instances for which all results are true is nonempty. 
    These instances are $\lambda_{\alpha}$-dominant, since the instance specific performance ratio $\lambda_{\alpha}(I)$ is shown
    to not decrease for every result. 
\end{proof}

We will use some of these results in the determination of upper and lower bounds in the following sections.

\section{Fixed Machine Number}\label{sec:FixedmachineNumber}
For \(m=2\) and \(m=3\), the approximation ratio admits an exact characterization.

We present our results in two Lemmas.

\begin{lemma}[Approximation ratio for $m=2$]\label{lem:mtwo}
    Let $\lambda^{2}_{\alpha}$ denote the approximation-ratio of $2\!\parallel\!\sum_{j} w_j C^{\alpha}_j$. Then for 
    $\alpha \in \Big[0,\frac12\Big]$,

   \[
\lambda^2_{\alpha}
=
\begin{cases}
1+\dfrac{1}{2\alpha\left(1+\sqrt{3+\frac{1}{\alpha}}\right)},
& \alpha\in \Big(0,\frac12\Big],\\[1.2ex]
\infty,
& \alpha=0.
\end{cases}
\]
\end{lemma}
The proof can be found in
\hyperlink{proof:mtwo}{Appendix~\ref*{app:proofmtwo}
(Proof of Lemma~\ref*{lem:mtwo})}.


\begin{lemma}[Approximation ratio for $m=3$]\label{lem: mthree}

    Let $\lambda^{3}_{\alpha}$ denote the approximation-ratio of $3\!\parallel\!\sum_{j} w_j C^{\alpha}_j$. Then for 
    $\alpha \in \Big[0,\frac12\Big]$,

   \[
\lambda^3_{\alpha}
=
\begin{cases}
1+
\frac{2}{\sqrt{\alpha(16+49\alpha)}+3\alpha}, & \alpha\in \Big(0,\frac12\Big],\\[1.2ex]
\infty,
& \alpha=0.
\end{cases}
\]
\end{lemma}

The proof can be found in \hyperlink{proof:mthree}{Appendix~\ref*{app:proof-mthree}
(Proof of Lemma~\ref*{lem: mthree})}.\medskip

\noindent
It is notable to state that $\lambda^3_{\frac12} = \lambda_{\frac12}$, which follows from \Cref{chap:alpha_g}, and that the expression 
found for $m=3$ is equal to the lower bound we find in Lemma \ref{lem: LB} for $\alpha \in \Big[\frac{1}{15}, \frac12\Big]$.


\newpage

\section{Upper Bound}\label{sec:Upper Bound}

In this section, we present two upper bounds, one of which is dependent on the number of machines.
\begin{framed}
For every \(\alpha\in \Big(0,\frac12\Big]\), the following bounds hold:

\begin{enumerate}[label=(\roman*)]
    \item
    \[
        \lambda_\alpha
        \leq
        1+\frac{1}{6\alpha}.
    \]

    \item For \(m\geq3\),
    \[
        \lambda_\alpha^m\leq U_m(\alpha),
    \]
    where
    \[
        U_m(\alpha)=
        \begin{cases}
            \displaystyle
            1+\frac{C_m}{2\sqrt{\alpha}},
            & 0<\alpha\leq z_m,\\[2.5mm]
            \displaystyle
            1+\frac{1+2z_m}{6(\alpha+z_m)},
            & z_m\leq\alpha\leq\frac12,
        \end{cases}
    \]
    with
    \[
        C_m
        =
        m\left(
            1+(2m-3)\sqrt{\frac{m-1}{2}}
        \right),
        \qquad
        z_m
        =
        \left(
            \frac{1}
            {3C_m+\sqrt{9C_m^2-2}}
        \right)^2.
    \]
\end{enumerate}

Moreover,
\[
    U_m(\alpha)
    \leq
    1+\frac{1}{6\alpha},
\]
with equality only at \(\alpha=\frac12\).
\end{framed}


We begin with the first upper bound, which is independent of the number of machines.
\newline
Consider the general expression of the instance specific approximation ratio using the identity from Lemma \ref{corr:queyranne}:
\[
    \lambda_{\alpha}(I)=\frac{\sum_{i=1}^{m} \left( L^{i}(\sigma^{W}(I)) \right)^2+(2\alpha-1)\sum_{j=1}^{n}p_{j}^{2}}
    {\sum_{i=1}^{m} \Big( L^{i}(\sigma^{*}(I)) \Big)^2+(2\alpha-1)\sum_{j=1}^{n}p_{j}^{2}}
\]
Defining
\[
    C:=\sum_{i=1}^{m} \left(L^{i}(\sigma^{*}(I))\right)^2 \quad \text{and} \quad K:=\sum_{j=1}^{n}p_{j}^{2}.
\]

Thus we can write $\lambda_{\alpha}(I)$ as
\[
    \lambda_{\alpha}(I)=\frac{\sum_{i=1}^{m} \left(L^{i}(\sigma^{W}(I))\right)^2+(2\alpha-1)K}{C+(2\alpha-1)K}.
\]
We know from \Cref{chap:alpha_g}, that for \(\alpha = \frac12\),
\[
\lambda_{\alpha}(I)
=
\frac{\sum_{i=1}^{m} \left(L^{i}(\sigma^{W}(I))\right)^2}{C}
\leq \frac43.
\]
This implies for any feasible instance

\begin{equation}
\label{eq: with C}
    \sum_{i=1}^{m} \left(L^{i}(\sigma^{W}(I))\right)^2
    \leq
    \frac43 C.
\end{equation}
This is true for any $\alpha$ (and not just for the case $\alpha=\frac12$) since if this were not the case 
for some $\alpha < \frac12$, the same instance  would achieve $\lambda_{\frac12} > \frac43$.

The function
 \[
    R\left(\sum_{i=1}^{m} L^{i}(\sigma^{W}(I))\right)
    :=
    \frac{\sum_{i=1}^{m}\left( L^{i}(\sigma^{W}(I))\right)^2 + \bigl(2\alpha-1\bigr)K}
         {C + \bigl(2\alpha-1\bigr)K},
\]
is monotonically increasing in $L^{i}(\sigma^{W}(I))$. Therefore, to consider the maximum, we may consider from now on only the case:
\[
    \sum_{i=1}^{m} \left(L^{i}(\sigma^{W}(I))\right)^2 = \frac43 C.
\]
Furthermore, we know that for any feasible instance, we have $C>0$ and $0 \le K \le C$. The
boundary case $K=C$ implies that there are exactly $m$ jobs, which gives $\lambda_{\alpha}(I) = 1$. The case
$K=0$ implies that all jobs have zero processing time, which is trivial.
Therefore, we can restrict our attention to $0<K<C$.

\noindent
Inserting the identity \eqref{eq: with C}

\[
    \lambda_{\alpha}(I) \le 1+\frac{\frac13 C}{C+(2\alpha-1)K}.
\]
The denominator is positive, so the expression is well-defined.


Because $C>0$, define
\[
    x:=\frac{K}{C}.
\]
The constraint $0<K<C$ is equivalent to
\[
    0<x<1.
\]
Using $K=xC$, we obtain
\[ 
    \lambda_{\alpha}(I) \le 1+\frac{1}{3(1+(2\alpha-1)x)}.
\]
The function
\[
    x\mapsto 1+\frac{1}{3(1+(2\alpha-1)x)}
\]
is strictly increasing on $(0,1)$.

\noindent
Therefore the supremum is reached only as $x \to 1$. This gives
\[
    \lambda_{\alpha}(I) \le 1+\frac{1}{3(1+(2\alpha-1))} = 1+\frac{1}{6\alpha}.
\]
As $I$ is arbitrary, it follows that
\[
\boxed{\lambda_{\alpha} \le 1+\frac{1}{6\alpha} \qquad \forall \alpha \in \Big(0,\frac12\Big]}.
\]

\subsection*{Machine-Dependent Upper Bounds}
We derive a strictly better upper bound of the approximation ratio at the cost of
introducing a dependency on the number of machines. 
We only consider instances with \(m\ge 3\), as we have already shown the exact approximation ratio for
$m=2$ in Section \ref{sec:FixedmachineNumber}.

We will make use of the following properties.

There exists a $\lambda_{\alpha}$-dominant instance $I$ such that:
\begin{enumerate}
 \item there exist $\sigma^*(I)$ and $\sigma^{W}(I)$ such that at most $m$ jobs are scheduled on any machine in both schedules.
 \item $n \le 3m-2$.
 \item some job $j^*$ is a singleton in an optimal schedule $\sigma^{*}(I)$ and there exists a cost maximizing
 schedule $\sigma^W(I)$ such that $j^*$ is the
 last job scheduled in $\sigma^{W}(I)$. Additionally, no job in $I$ finishes strictly after $j^*$ in $\sigma^{W}$.
 \item $\lambda_{\frac12} = \frac43$.
 \item $n > m$, (else $\lambda_{\alpha}(I) =1$).
\end{enumerate}

Next, we derive a form of the objective very similar to the one from Lemma \ref{corr:queyranne}, which was based on a result from \cite{queyranne1993structure}.
Let $J_i(\sigma(I))$ denote the set of jobs scheduled on machine $i$ in schedule $\sigma(I)$.
Fix a machine $i$ and write its jobs in order of processing in schedule $\sigma(I)$. 
\[
 j_{i,1},j_{i,2},\ldots,j_{i,r_i},
\]
where $r_i \coloneqq |J_i(\sigma(I))|$.
The cost they contribute to the total cost of the schedule is:
\[
\begin{aligned}
\sum_{s=1}^{r_i}p_{j_{i,s}}C_{j_{i,s}}^{\alpha}(\sigma(I))
&=
\sum_{s=1}^{r_i}
p_{j_{i,s}}
\left(
\sum_{t=1}^{s-1}p_{j_{i,t}}
+\alpha p_{j_{i,s}}
\right)\\
&=
\alpha\sum_{j \in J_i(\sigma(I))}p_j^2
+
\sum_{\{j,k\}\subseteq J_i(\sigma(I))}p_jp_k.
\end{aligned}
\]

\noindent
Summing over all machines gives
\begin{equation}
     \gamma_{\alpha}(\sigma(I))=\alpha Q+P(\sigma(I)).
     \label{eq: F,P,Q}
\end{equation}

\noindent
where:
\[
 Q:=\sum_jp_j^2
\]
and,
\[
 P(\sigma(I))\coloneqq
 \sum_{i=1}^m\ \sum_{\{j,k\}\subseteq J_i(\sigma(I))}p_jp_k.
\]

We therefore split the cost up into fixed costs $Q$ independant of the schedule $\sigma(I)$
and "assignment cost" which we denote by $P(\sigma(I))$.\medskip 
\noindent
To avoid cumbersome notation, we write
\[
 P_{W}:=P(\sigma^{W}(I)),\qquad P^*:=P(\sigma^{*}(I)).
\]

\noindent
We will make use of Theorem \ref{thm:job-bound}, which bounded the number of jobs of a $\lambda^{m}_{\alpha}$-maximizing instance,
\begin{equation}
 n\le 3m-2.
 \label{eq:job-number}
\end{equation}

\noindent
Define
\[
 \mu_m:=\sqrt{\frac{m-1}{2}},
 \qquad
 K_m:=2(m-1),
 \qquad
 C_m:=m\left(1+(K_m-1)\mu_m\right).
\]

$K_m$ is chosen such that $n-m \le K_m$. We will see why this is useful in the proof of Lemma $\ref{lem:pair-cost-gap}$.

\begin{lemma}
\label{lem:pair-cost-gap}
Under the assumptions above,
\begin{equation}
    P_{W}-P^*\le C_m\sqrt{QP^*}. 
    \label{eq:pair-cost-gap}
\end{equation} 
\end{lemma}

\begin{proof}
Consider the jobs in order of any list $l \sim \sigma^{W}(I)$.  When job \(j\) is scheduled, let \(\ell_j\)
be the current load of the machine on which it is to be scheduled. The assignment incurs 
cost \(p_j\ell_j\), and therefore
\[
 P_{W}=\sum_jp_j\ell_j.
\]

\noindent
Fix any job \(j\) with \(\ell_j>0\) (this exists as we schedule at least $m+1$ jobs).
We now consider only the partial schedule consisting of all jobs scheduled before $j$.
Since $j$ is assigned to a least-loaded machine,
every machine has current load at least \(\ell_j\).
Choose a largest currently assigned job from each machine (excluding j).  At that moment
each machine contains at most \(m\) jobs, so every selected job has size at
least
\[
 \frac{\ell_j}{m}.
\]
The \(m\) selected jobs and \(j\) are \(m+1\) distinct jobs in total. Thus, by the pigeonhole principle,
two of them share a machine in $\sigma^*(I)$.

\noindent
If the colliding pair consists of two selected jobs, then
\[
 P^*\ge\frac{\ell_j^2}{m^2}.
\]
Since \(p_j\le\sqrt Q\), it follows that
\begin{equation}
 p_j\ell_j\le m\sqrt{QP^*}.
 \label{eq: pjlj-bound}
\end{equation}

\noindent
If the colliding pair consists of \(j\) and a selected job, then
\[
 P^*\ge\frac{p_j\ell_j}{m},
\]
and hence
\begin{equation}
 p_j\ell_j\le mP^*.
 \label{eq: pjlj-bound-2}
\end{equation}

\noindent
We make use of the following inequality. 
For any \(k\) positive numbers \(x_1,\ldots,x_k\),

\begin{equation}
 \sum_{r<s}x_rx_s
 \le \frac{k-1}{2}\sum_{r=1}^kx_r^2.
 \label{eq: Cauchy-Schwarz}
\end{equation}

This follows from the Cauchy-Schwarz inequality and was shown in \cite{rozenberg2020inequality}.
\noindent
Applying \eqref{eq: Cauchy-Schwarz} to every  set $J_{i}(\sigma^{*}(I))$ and making use of the fact that $|J_{i}(\sigma^{*}(I))|\le m$, then summing over all machines gives:
\begin{equation}
 P^*\le\frac{m-1}{2}Q.
\label{eq: P* bound}
\end{equation}

\noindent
Multiplying \eqref{eq: P* bound} by \(P^*\ge 0\), we obtain
\[
(P^*)^2
\le
\frac{m-1}{2}QP^*.
\]

\noindent
Taking square roots on both sides yields
\[
P^*
\le
\sqrt{\frac{m-1}{2}}\sqrt{QP^*}.
\]

\noindent
Substituting this estimate into \eqref{eq: pjlj-bound-2}, we conclude that

\begin{equation}
p_j\ell_j
\le mP^*
\le
m\sqrt{\frac{m-1}{2}}\sqrt{QP^*}.
\label{eq: pjlj bound 3}
\end{equation}


\noindent
Consequently, every placement with \(\ell_j>0\) satisfies
\begin{equation}
 p_j\ell_j\le m\mu_m\sqrt{QP^*}.
\label{eq: pjlj bound 4}
\end{equation}

The first \(m\) jobs of an ordinary list schedule are placed on empty
machines and have \(\ell_j=0\).  By Theorem \ref{thm:job-bound}, the number of placements with strictly
positive $\ell_j$ is at most
\[
 n-m\le (3m-2)-m =2(m-1) =K_m.
\]

We now use both special properties of the singleton $j^*$. Because of Theorem \ref{corr: singletonjob}, $j^*$ may be assumed to be scheduled last on some machine
and is not a singleton in $\sigma^{W}(I)$ (Corollary \ref{corr:singleton-not-worst}). Thus
at least $m$ jobs precede it in some corresponding list $l \sim \sigma^{W}(I)$. Therefore, its preceding load \(\ell_{j^*}\) is positive. Immediately before placing $j^*$,
choose a largest job from each of the \(m\) machines in $\sigma^{W}(I)$. Since the instance is assumed to be maximally unified, none of these
selected jobs can be assigned to the same machine as $j^*$ in $\sigma^*(I)$. Consequently the \(m\) selected jobs occupy only
\(m-1\) optimal machines, so by the pigeonhole principle two selected jobs share a machine in $\sigma^*(I)$.
An application of \eqref{eq: pjlj-bound} to $j^*$ gives
\begin{equation}
 p_{j^*}\ell_{j^*}\le m\sqrt{QP^*}.
 \label{eq: j^*lj bound}
\end{equation}

There are at most \(K_m-1\) other positive-load assignments ($\ell_j > 0$). Summing \eqref{eq: pjlj bound 4}
over those placements and adding \eqref{eq: j^*lj bound}, we obtain
\[
\begin{aligned}
 P_{W}
 \le
 m\left(1+(K_m-1)\mu_m\right)\sqrt{QP^*}
 =C_m\sqrt{QP^*}.
\end{aligned}
\]
Finally \(P_W-P^*\le P_W \le C_m \sqrt{QP^*}\) and thus \eqref{eq:pair-cost-gap} holds.
\end{proof}


Expanding the established result for $\alpha = \frac12$ by means of \eqref{eq: F,P,Q} gives
\[
\gamma_{\alpha}(\sigma^W(I))=
 \frac12Q+P_{W}
 \le
 \frac43\left(\frac12Q+P^*\right),
\]
and hence
\begin{equation}
 \gamma_{\alpha}^W(I)- \gamma_{\alpha}^*(I)P_W-P^*\le\frac16Q+\frac13P^*.
\label{eq: PW, P* bound}
\end{equation}
We now have two bounds on the cost difference between maximal cost 
schedule and the optimal schedule. We will use these to bound
the approximation ratio. By \eqref{eq: F,P,Q},
\[
\begin{aligned}
\frac{\gamma_{\alpha}^W(I)}{\gamma_{\alpha}^*(I)}
=
\frac{\alpha Q+P_W}{\alpha Q+P^*}
=
1+\frac{P_W-P^*}{\alpha Q+P^*}.
\end{aligned}
\]
Using \eqref{eq: PW, P* bound} to bound the numerator gives
\[
\frac{\gamma_{\alpha}^W(I)}{\gamma_{\alpha}^*(I)}
\le
1+
\frac{\frac16Q+\frac13P^*}{\alpha Q+P^*}.
\]
Using Lemma~\ref{lem:pair-cost-gap} instead gives
\[
\frac{\gamma_{\alpha}^W(I)}{\gamma_{\alpha}^*(I)}
\le
1+
\frac{C_m\sqrt{QP^*}}{\alpha Q+P^*}.
\]
Both inequalities hold simultaneously. We can therefore take the smaller
of the two right-hand sides:
\[
\frac{\gamma_{\alpha}^W(I)}{\gamma_{\alpha}^*(I)}
\le
1+
\min\left\{
\frac{\frac16Q+\frac13P^*}{\alpha Q+P^*},
\frac{C_m\sqrt{QP^*}}{\alpha Q+P^*}
\right\}.
\]
Finally, since \(Q>0\), define
\[
 z:=\frac{P^*}{Q}.
\]
Dividing the numerator and denominator of both fractions by \(Q\) and
using
\[
\frac{\sqrt{QP^*}}{Q}
=
\sqrt{\frac{P^*}{Q}}
=\sqrt z,
\]
gives
\begin{equation}
 \frac{\gamma_{\alpha}^W(I)}{\gamma_{\alpha}^*(I)}
 \le
 1+
 \min\left\{
 \frac{\frac16+\frac13z}{\alpha+z},
 \frac{C_m\sqrt z}{\alpha+z}
 \right\}.
\label{eq: ratio bound by min}
\end{equation}
Next we determine which of the two expressions is smaller for each value of $z$.
To this end, define \(z_m\) to be the smaller positive solution of
\begin{equation}
 \frac16+\frac13z=C_m\sqrt z.
\label{eq: equal bounds}
\end{equation}
We are only interested in \emph{positive} solutions as negative solutions do not 
correspond to a feasible instance. We are only interested in the \emph{smallest} positive solution
since this will give us the better bound. 

Equivalently, define
\begin{equation}
     z_m:=\left(\frac{1}{3C_m+\sqrt{9C_m^2-2}}\right)^2.
     \label{eq: z_m}
\end{equation}

Note that, since $C_m > 1$, $z_m$ is indeed real, positive and $z_m < \frac12$.

\begin{corollary}
\label{thm:main}
Let $I$ be an instance of $Pm\!\parallel\!\sum_{j} p_jC^{\alpha}_j$
Then for every $\alpha\in \Big(0,\frac12\Big]$,
\begin{equation}
 \lambda_{\alpha}(I)
 \le U_m(\alpha),
 \label{eq: lambdam Um}
\end{equation}
where
\begin{equation}
 U_m(\alpha)\coloneqq
 \begin{cases}
 \displaystyle
 1+\frac{C_m}{2\sqrt\alpha},
 &0<\alpha\le z_m,\\[3mm]
 \displaystyle
 1+\frac{C_m\sqrt{z_m}}{\alpha+z_m}
 =
 1+\frac{1+2z_m}{6(\alpha+z_m)},
 &z_m\le\alpha\le\frac12.
 \end{cases}
 \label{eq: Um}
\end{equation}
Moreover,
\[
 U_m(\alpha)<1+\frac{1}{6\alpha}
 \qquad\text{for} \qquad \alpha \in \Big(0,\frac12\Big),
\]
and
\[
 U_m\left(\frac12\right)=\frac43 \qquad \forall m\ge 2.
\]
\end{corollary}

\begin{proof}
For fixed \(\alpha\), define
\[
 g_1(z):=\frac{\frac16+\frac13z}{\alpha+z},
 \qquad
 g_2(z):=\frac{C_m\sqrt z}{\alpha+z}.
\]
For \(0<\alpha\le \frac12\),
\[
 g_1'(z)=\frac{\alpha/3-1/6}{(\alpha+z)^2}\le0.
\]
By \eqref{eq: equal bounds}, \(g_1(z_m)=g_2(z_m)\).
We determine which of the two functions gives the tighter bound for each
value of $z$, as illustrated in Figure~\ref{fig:g1-g2-comparison}.
Suppose first that \(\alpha\ge z_m\). Since $C_m > 0$, the function \(g_2\) is increasing
on \([0,\alpha]\) and hence on \([0,z_m]\).  Thus, for \(z\le z_m\),
\[
 \min\{g_1(z),g_2(z)\}\le g_2(z_m),
\]
while for \(z\ge z_m\), monotonicity of \(g_1\) gives
\[
 \min\{g_1(z),g_2(z)\}\le g_1(z)\le g_1(z_m).
\]
This proves the second branch of \eqref{eq: Um}.

Suppose next that \(0<\alpha\le z_m\).  The global maximum of \(g_2\) occurs
at \(z=\alpha\), where
\[
 g_2(\alpha)=\frac{C_m}{2\sqrt\alpha}.
\]
For \(z\le z_m\), the minimum in \eqref{eq: ratio bound by min} is at most this global maximum.  For
\(z\ge z_m\), it is at most
\[
 g_1(z_m)=g_2(z_m)\le g_2(\alpha).
\]
This proves the first branch of \eqref{eq: Um} and therefore \eqref{eq: lambdam Um}.

\begin{figure}[ht]
\centering
\begin{minipage}{0.48\textwidth}
\centering
\begin{tikzpicture}
% Numerical values used only to render the qualitative plot.
% The visible labels remain symbolic.
\def\alphaval{0.08}
\def\Cmplot{1.0}
\def\zone{0.03136}
\begin{axis}[
    xmode=log,
    width=\textwidth,
    height=0.72\textwidth,
    axis lines=left,
    xlabel={},
    ylabel={value},
    xmin=0.008,
    xmax=20,
    ymin=0,
    ymax=3.0,
    samples=400,
    domain=0.008:20,
    legend style={draw=none, fill=none, at={(0.98,0.98)}, anchor=north east, font=\scriptsize},
    xtick={\zone,\alphaval},
    xticklabels={$z_m$,$\alpha$},
    ytick=\empty,
    tick style={draw=none},
    clip=false,
]
  \addplot[black, thick]
    {(1/6 + x/3)/(\alphaval + x)};
  \addlegendentry{$g_1(z)=\frac{1/6+z/3}{\alpha+z}$}

  \addplot[black, thick, dashed]
    {\Cmplot*sqrt(x)/(\alphaval + x)};
  \addlegendentry{$g_2(z)=\frac{C_m\sqrt z}{\alpha+z}$}

  \draw[dotted] (axis cs:\zone,0) --
    (axis cs:\zone,{(1/6+\zone/3)/(\alphaval+\zone)});
  \draw[dotted] (axis cs:\alphaval,0) --
    (axis cs:\alphaval,{\Cmplot*sqrt(\alphaval)/(\alphaval+\alphaval)});
\end{axis}
\end{tikzpicture}
\[
 z_m<\alpha
\]
\end{minipage}
\hfill
\begin{minipage}{0.48\textwidth}
\centering
\begin{tikzpicture}
% Same qualitative plot, but now the maximum of g_2 occurs before z_m.
% The numerical values are rescaled only to keep both panels visually comparable.
\def\alphaval{0.035}
\def\Cmplot{0.8}
\def\zone{0.05311}
\begin{axis}[
    xmode=log,
    width=\textwidth,
    height=0.72\textwidth,
    axis lines=left,
    xlabel={},
    ylabel={value},
    xmin=0.03,
    xmax=20,
    ymin=0,
    ymax=3.0,
    samples=400,
    domain=0.03:20,
    legend style={draw=none, fill=none, at={(0.98,0.98)}, anchor=north east, font=\scriptsize},
    xtick={\alphaval,\zone},
    xticklabels={$\alpha$,$z_m$},
    ytick=\empty,
    tick style={draw=none},
    clip=false,
]
  \addplot[black, thick]
    {(1/6 + x/3)/(\alphaval + x)};
  \addlegendentry{$g_1(z)=\frac{1/6+z/3}{\alpha+z}$}

  \addplot[black, thick, dashed]
    {\Cmplot*sqrt(x)/(\alphaval + x)};
  \addlegendentry{$g_2(z)=\frac{C_m\sqrt z}{\alpha+z}$}

  \draw[dotted] (axis cs:\zone,0) --
    (axis cs:\zone,{(1/6+\zone/3)/(\alphaval+\zone)});
  \draw[dotted] (axis cs:\alphaval,0) --
    (axis cs:\alphaval,{\Cmplot*sqrt(\alphaval)/(\alphaval+\alphaval)});
\end{axis}
\end{tikzpicture}
\[
 \alpha<z_m
\]
\end{minipage}
\caption{Comparison of \(g_1\) and \(g_2\).}
\label{fig:g1-g2-comparison}
\end{figure}

It remains to compare \eqref{eq: Um} with the known bound. 

If
\(z_m\le\alpha<1/2\), then \(g_1\) is strictly decreasing, so
\[
 \frac{1+2z_m}{6(\alpha+z_m)}
 =g_1(z_m)<g_1(0)=\frac{1}{6\alpha}.
\]
If \(0<\alpha\le z_m\), \eqref{eq: equal bounds} gives
\[
 3C_m\sqrt{z_m}=\frac12+z_m<1.
\]
Here \(z_m<1/2\) follows immediately from \(C_m\ge2\) and \eqref{eq: z_m}.
Therefore
\[
 3C_m\sqrt\alpha\le3C_m\sqrt{z_m}<1,
\]
which implies
\[
 1+\frac{C_m}{2\sqrt\alpha}<1+\frac{1}{6\alpha}.
\]


Finally, substituting \(\alpha=1/2\) into the second branch of \eqref{eq: Um} gives
\[
 \frac{1+2z_m}{6(\frac12+z_m)}=\frac13,
\]
and hence this bound is exact for $\alpha =\frac12$.
\end{proof}

\begin{quote}
    \small\normalfont
    \textbf{Note.}
    For all $\alpha \in (0,\frac12]$ and $m\ge2$, our bound $U_m(\alpha)$
    strictly improves upon the upper bound $1+\frac{m-1}{2m\alpha}$
    given in \cite{JaegerSkutella2018}.
\end{quote}


\newpage
\section{Lower Bound}\label{sec:lower-bound}

\subsection*{Instance family}

We aim to find a lower bound on $\lambda_{\alpha}$ by considering a family of instances and corresponding schedules. 
\begin{lemma}\label{lem: LB}
Given any $\alpha \in \Big(0,\frac12\Big]$,
define 
\[
B_m(\alpha)
\coloneqq
1+
\frac{m-1}{
\alpha m\left(
1+\sqrt{
1+
\frac{
4(m-1)\left((m-1)+\alpha(m^2-2m+2)\right)
}{\alpha m^2}
}
\right)}, \qquad m \ge 3
\]

Then it holds for the approximation ratio of $\Problem$ that
\[
\boxed{
\lambda_{\alpha}\ge 
\begin{cases}
B_3(\alpha),
& \displaystyle \alpha_3\leq\alpha\leq\frac12,\\[6pt]
B_{m+1}(\alpha),
& \alpha_{m+1}\leq\alpha<\alpha_{m},\quad m\geq 3,
\end{cases}
}
\]
where 
\[
\alpha_m \coloneqq \frac{1}{m^4-2m^3-2m^2+m+3}, \qquad m \geq 3
\]
and every bound \(B_m(\alpha)\) is attained by an instance with \(m\) machines.
\end{lemma}


\begin{proof}
We arrive at this bound by constructing a family of instances that make
use of some of the structural results of $\lambda_{\alpha}$-dominant instances determined in
\Cref{sec:Struct-Results}. We do impose, however, some additional restrictions, which simplify the determination 
of cost for both a cost maximizing and optimal schedule. 
Our generalized instance $I$ has $u$ many jobs of processing time $p$ which are singleton in the optimal schedule $\sigma^{*}(I)$. This aligns with the findings of Theorem \ref{corr: singletonjob}.
The remaining machines all have equal makespan, which we normalize to $1$ and have at least two jobs scheduled. We only allow two different machine loads in $\sigma^{*}(I)$, as this greatly simplifies the
determination of cost. The schedule $\sigma^*(I)$ is depicted in Figure \ref{fig:instance-structure}.\medskip

The schedule $\sigma^{W}(I)$ results from $\sigma^{*}(I)$ by spreading the load of the $m-u$ machines with load $1$ to all $m$ machines resulting in an equal load of
$\frac{m-u}{m}$ on all machines. 
We then schedule the remaining $u$ jobs of size $p$ on top of this load. This again follows the idea of the finding that there must be a singleton in some optimal schedule after which no other job finishes 
in some cost-maximizing schedule $\sigma^{W}(I)$. 

We assume for this family of jobs that $u \geq 1$.
All other steps of construction are such that we determine a smallest set of jobs that meet the criteria named above. 
We are interested in a smallest set, since we know from Lemma \ref{corr:unification} that any unification and therefore reduction of the job 
number is beneficial in the search for a $\lambda_{\alpha}$-dominant instance. \medskip

We now restrict attention to instances of the above form for which $\frac{1}{k+1} \leq \frac{u}{m} \leq \frac{1}{k}$ for some $k \geq 2$.
We define the lengths of the jobs that are not singleton in $\sigma^*$ as follows. We have $m-u$ jobs of size $1-\frac{u}{m}$ and $(k-1)u$ jobs of size $\frac{u}{m}$.
There remains a block of jobs of total load $(m-ku)\frac{u}{m}$. We split this block into jobs similar to the construction in Lemma \ref{corr:splitting-jobs}, using as few cuts as possible so that its load can be distributed equally over $u$ machines.
The load that we need for each machine in $\sigma^{W}(I)$ is $\frac{m-ku}{m}$.
\begin{figure}[ht]
\centering

\begin{tikzpicture}[
    x=0.7cm,
    y=0.7cm,
    job/.style={
        draw,
        thick,
        fill=white,
        font=\small
    },
    machine/.style={font=\small}
]

% ------------------------------------------------
% Template idea:
% row      = displayed vertical row in the picture
% machine  = actual machine label shown on the left
%
% Use \machinerow{row}{machine label}
% Use \dotsrow{row}
% Use \jobblock{row}{start}{length}{job label}
% Use \jobblock[jobhighlight]{row}{start}{length}{job label} for shading
% ------------------------------------------------

% Displayed machine rows
\machinerow{1}{1}
\dotsrow{2}
\machinerow{3}{u(k-1)}
\machinerow{4}{u(k-1)+1}
\dotsrow{5}
\machinerow{6}{m-u}
\machinerow{7}{m-u+1}
\dotsrow{8}
\machinerow{9}{m}

% Example jobs
\jobblock{1}{0}{2}{\frac{u}{m}}
\jobblock{1}{2}{3}{1-\frac{u}{m}}

\jobblock{3}{0}{2}{\frac{u}{m}}
\jobblock{3}{2}{3}{1-\frac{u}{m}}

\jobblock[jobhighlight]{4}{0}{2}{split}
\jobblock{4}{2}{3}{1-\frac{u}{m}}

\jobblock[jobhighlight]{6}{0}{2}{split}
\jobblock{6}{2}{3}{1-\frac{u}{m}}
\jobblock{7}{0}{10}{p}
\jobblock{9}{0}{10}{p}

\end{tikzpicture}

\caption{Optimal schedule with $u$ singleton jobs of size $p$ and $m-u$ machines with load $1$.}
\label{fig:instance-structure}

\end{figure}
\subsection*{Construction}

For simplicity, we assume throughout that
\[
\gcd(m,u)=1
\]

The problem we are faced with is having \(N = m-ku\) unit lengths which we want to split using as few cuts as possible, so that they can be regrouped into \(u\) parts of equal total length.

Since \(0<N<u\), each of the \(u\) final parts must have total length
\[
\frac{N}{u}
=\frac{m-ku}{u}.
\]
Since \(\gcd(m,u)=1\), we also have
\[
\gcd(N,u)=1.
\]

\begin{figure}[ht]
\centering
\begin{tikzpicture}[x=1.1cm,y=1cm]
  \draw[thick] (0,0) -- (8,0);
  \foreach \x/\lab in {0/0,1.6/1,3.2/2,4.8/{\cdots},6.4/{N-1},8/N} {
    \draw[thick] (\x,-0.16) -- (\x,0.16);
    \node[below,font=\small,fill=white,inner sep=1pt] at (\x,-0.18) {\(\lab\)};
  }

  \foreach \x/\lab in {1.05/{N/u},2.1/{2N/u},3.15/{3N/u},5.25/{\cdots},7.15/{(u-1)N/u}} {
    \draw[dashed] (\x,-0.5) -- (\x,0.5);
    \node[above,font=\small,fill=white,inner sep=1pt] at (\x,0.5) {\(\lab\)};
  }

  \node[below=0.65cm,font=\small] at (4,0) {};
\end{tikzpicture}
\caption{The \(N=m-ku\) unit lengths are placed in one strip. The cuts divide the strip into \(u\) intervals of equal length \(N/u\).}
\label{fig:construction}
\end{figure}


We line up the \(N=m-ku\) unit lengths into one block of total length \(N\). Make cuts at
\[
\frac{N}{u},\frac{2N}{u},\frac{3N}{u},\ldots,\frac{(u-1)N}{u}.
\]

The cut points divide the interval \([0,N]\) into \(u\) consecutive intervals
\[
\mathcal{I}_j=\left[\frac{(j-1)N}{u},\frac{jN}{u}\right],
\qquad
j=1,2,\ldots,u.
\]
Each interval \(\mathcal{I}_j\) has length \(N/u\), so each \(\mathcal{I}_j\) is one of the desired final machine loads.
The unit pieces, however, have boundaries at the integers $0,1,2,\ldots,N$.
Since \(N<u\), each interval \(\mathcal{I}_j\) has length less than \(1\). Therefore an interval \(\mathcal{I}_j\) can contain at most one integer boundary in its interior.

If \(\mathcal{I}_j\) contains no integer boundary, then it remains as one piece of length
\[
\frac{N}{u}.
\]
If \(\mathcal{I}_j\) contains an integer boundary \(i\), then that boundary splits \(\mathcal{I}_j\) into two pieces.

\begin{figure}[ht]
\centering
\begin{tikzpicture}[x=1.35cm,y=1cm]
  \draw[thick] (0,0) -- (5,0);
  \draw[dashed] (0,-0.5) -- (0,0.5);
  \draw[dashed] (5,-0.5) -- (5,0.5);
  \draw[thick] (2.15,-0.35) -- (2.15,0.35);

  \node[below,font=\small,fill=white,inner sep=1pt] at (0,-0.6) {\(\frac{(j-1)N}{u}\)};
  \node[below,font=\small,fill=white,inner sep=1pt] at (5,-0.6) {\(\frac{jN}{u}\)};
  \node[above,font=\small,fill=white,inner sep=1pt] at (2.15,0.4) {\(i\)};

  \draw[decorate,decoration={brace,mirror,amplitude=5pt}] (0,-1.3) -- (2.15,-1.3)
    node[midway,below=7pt,font=\small,fill=white,inner sep=1pt] {\(\frac{\alpha}{u}\)};
  \draw[decorate,decoration={brace,mirror,amplitude=5pt}] (2.15,-1.3) -- (5,-1.3)
    node[midway,below=7pt,font=\small,fill=white,inner sep=1pt] {\(\frac{N-\alpha}{u}\)};
  \draw[decorate,decoration={brace,amplitude=5pt}] (0,0.95) -- (5,0.95)
    node[midway,above=7pt,font=\small,fill=white,inner sep=1pt] {\(\frac{N}{u}\)};
\end{tikzpicture}
\caption{An interval \(\mathcal{I}_j\) has length \(N/u<1\), so it contains at most one unit boundary. If it contains one, it is split into two lengths.}
\label{fig:interval}
\end{figure}

There are \(N-1\) interior integer boundaries:
\[
1,2,\ldots,N-1.
\]
Thus \(N-1\) of the intervals \(\mathcal{I}_j\) are split into two pieces, and the remaining
\[
u-(N-1)=u-N+1
\]
intervals are not split. Therefore there are $u-N+1$ many
pieces of length $\frac{N}{u}$.

It remains to identify the lengths created by the $N-1$ interior integer boundaries. Fix an integer boundary
\[
i\in\{1,2,\ldots,N-1\}.
\]
It lies in a unique interval \(\mathcal{I}_j\), so
\[
\frac{(j-1)N}{u}<i<\frac{jN}{u}.
\]
The two pieces created around \(i\) have lengths
\[
i-\frac{(j-1)N}{u}
\]
and
\[
\frac{jN}{u}-i.
\]
Multiplying each by \(u\), these lengths become
\[
iu-(j-1)N
\]
and
\[
jN-iu.
\]
These are positive integers whose sum is \(N\). Hence the two lengths are
\[
\frac{iu-(j-1)N}{u}
\qquad\text{and}\qquad
\frac{jN-iu}{u}.
\]
Now set
\[
r \coloneqq iu-(j-1)N \in \mathbb{N}.
\]
Then, because
\[
\bigl(iu-(j-1)N\bigr)+\bigl(jN-iu\bigr)=N.
\]
the two lengths can be written as
\[
\frac{r}{u}
\qquad\text{and}\qquad
\frac{N-r}{u}
\]
for some \(r\in\{1,2,\ldots,N-1\}\).

As \(i\) runs through \(1,2,\ldots,N-1\), the residuals
\[
iu \pmod N
\]
run through all nonzero residuals mod \(N\), because \(\gcd(N,u)=1\). 
Therefore the values 
\[
 r = 1,2,\ldots,N-1,
\]
are obtained \textbf{exactly once} each. Each such boundary contributes the pair
\[
\frac{r}{u}
\quad\text{and}\quad
\frac{N-r}{u}.
\]
Consequently, every length
\[
\frac{1}{u},\frac{2}{u},\ldots,\frac{N-1}{u}
\]
appears \textbf{exactly twice}.

Combining this with the unsplit intervals gives the full list of job lengths:
\[
2 \text{ pieces of length } \frac{1}{u},
\quad
2 \text{ pieces of length } \frac{2}{u},
\quad
\ldots \quad ,
\quad
2 \text{ pieces of length } \frac{N-1}{u}
\]
and
\[
u-N+1 \text{ pieces of length } \frac{N}{u}.
\]


For the family of instances introduced above, we do not split unit-length intervals but intervals of length $\frac{u}{m}$.
Scaling all pieces by $\frac{u}{m}$, we get the following list of pieces:

\[
2 \text{ pieces of length } \frac{1}{m},
\quad
2 \text{ pieces of length } \frac{2}{m},
\quad
\ldots,
\quad
2 \text{ pieces of length } \frac{N-1}{m}=\frac{(m-ku)-1}{m}
\]
and

\[
u-(m-ku)+1=((k+1)u-m+1) \text{ pieces of length }  \frac{N}{m}=\frac{m-ku}{m}.
\]


We can now write down $\lambda_{\alpha}(I)$ for the family of instances introduced 
above. Let $m \ge 3$ be the number of machines, and let \(u\) be the number of jobs 
of size \(p\) that are singleton in \(\sigma^{*}(I)\). We only consider
$ 0 < \alpha \leq \frac12$

From now on, we will associate the family of instances with the tuple $(m,u,p)$ as these parameters
uniquely determine its explicit form. 


We want to maximize
\[
\lambda_{\alpha}(m,u,p)=\frac{
(m-u)\left(1-\frac um\right)^2
 +u\left(p+1-\frac um\right)^2
 +(2\alpha-1)B
}{
(m-u)+up^2+(2\alpha-1)B
},
\]
where
\[
\begin{aligned}
B&\coloneqq\sum_{j=1}^{m-ku-1}2\left(\frac{j}{m}\right)^2+\left((k+1)u-m+1\right)\left(\frac{m-ku}{m}\right)^2
+up^2
+(m-u)\left(1-\frac um\right)^2
+(k-1)\frac{u^3}{m^2}\\
&=\frac{(m-ku-1)(m-ku)(2m-2ku-1)}{3m^2}  \\
&+\left((k+1)u-m+1\right)\left(\frac{m-ku}{m}\right)^2
+up^2
+(m-u)\left(1-\frac um\right)^2
+(k-1)\frac{u^3}{m^2}.
\end{aligned}
\]
$B$ is the sum of squared processing times for all jobs of the instance. 

The rest of the proof is algebraic and involves no more constructive ideas. 
It can be found in \hyperlink{proof:LB}{Appendix~\ref{app:proof-LB}}
(Proof of Lemma~\ref*{lem: LB})
\end{proof}

\section{Computational Analysis}
Since the gap between the upper and lower bounds for $\lambda_{\alpha}$ established in the previous 
chapters remains substantial, we aim to investigate computationally the value towards which the 
approximation ratio may converge. To this end, we formulate a mixed-integer nonlinear program (MINLP).
The main challenge is to ensure that every solution produced by the program corresponds to a 
list schedule. The number of inequalities required for a given number of machines and jobs grows rapidly,
as is already apparent from the explicit determination of the approximation ratio for the special 
cases of two and three machines (\hyperlink{proof:mthree}{Appendix~\ref*{app:proof-mthree}
(Proof of Lemma~\ref*{lem: mthree})}). We therefore make the following compromise: we require the last 
job assigned to each machine in the worst-case schedule to have the same starting time.
This condition immediately guarantees that the resulting schedule is a list schedule,
although it restricts the feasible region of the MINLP.
The structure of $\sigma^{\mathrm{WSPT}}$-schedules considered is depicted in Figure \ref{fig:computational-schedule-structure}.

\begin{figure}[ht]
\centering

\begin{tikzpicture}[
    x=0.7cm,
    y=0.7cm,
    job/.style={
        draw,
        thick,
        fill=white,
        font=\small
    },
    machine/.style={font=\small}
]

% ------------------------------------------------
% Template idea:
% row      = displayed vertical row in the picture
% machine  = actual machine label shown on the left
%
% Use \machinerow{row}{machine label}
% Use \dotsrow{row}
% Use \jobblock{row}{start}{length}{job label}
% Use \jobblock[jobhighlight]{row}{start}{length}{job label} for shading
% ------------------------------------------------

% Displayed machine rows
\machinerow{1}{1}
\machinerow{2}{2}
\machinerow{3}{3}
\dotsrow{4}
\machinerow{5}{m-1}
\machinerow{6}{m}

% Example jobs
\jobblock{1}{0}{1}{}
\jobblock{1}{1}{2}{}
\jobblock{1}{3}{1}{}
\jobblock{1}{4}{1}{}
\jobblock{1}{5}{3}{f_1}

\jobblock{2}{0}{3}{}
\jobblock{2}{3}{0.5}{}
\jobblock{2}{3.5}{1.5}{}
\jobblock{2}{5}{5}{f_2}

\jobblock{3}{0}{1}{}
\jobblock{3}{1}{1}{}
\jobblock{3}{2}{1}{}
\jobblock{3}{3}{2}{}
\jobblock{3}{5}{4}{f_3}

\jobblock{5}{0}{2}{}
\jobblock{5}{2}{3}{}
\jobblock{5}{5}{3.5}{f_{m-1}}

\jobblock{6}{0}{4}{}
\jobblock{6}{4}{0.5}{}
\jobblock{6}{4.5}{0.5}{}
\jobblock{6}{5}{3}{f_m}

% Timestamp marker where all f jobs begin
\draw[dashed, thick] (5,0.2) -- (5,-6.0);
\node[above] at (5,0.2) {$1$};

\end{tikzpicture}

\caption{MINLP is restricted to $\mathrm{WSPT}$-schedules of this form. The determination of an optimal schedule $\sigma^*$ is free for the algorithm to choose.}
\label{fig:computational-schedule-structure}
\end{figure}

Formally, fix \(m\in\mathbb{N}\) and \(\alpha \in\Big(0,\frac12\Big]\) and normalize the common starting
time of the last jobs in \(\sigma^W\) to \(1\). Let
\[
  \mathcal{I}\coloneqq\{(i,j): i,j=1,\ldots,m\}.
\]
The set \(\mathcal{I}\) indexes the variables \(p_{ij}\). We define
one combined indexing set
\[
  \mathcal{K}\coloneqq\mathcal{I}\cup\{1,\ldots,m\}.
\]

The continuous decision variables are
\[
  p_{ij}\ge 0\quad ((i,j)\in \mathcal{I}),\qquad
  f_i\ge 0\quad i \in [m],
\]
where \(p_{ij}\) denotes the processing time of the \(j\)-th listed job scheduled on
machine \(i\) and $f_i$ denotes the processing time of the job that is processed last on machine \(i\).
The binary assignment variables are
\[
  x_{kg}\in\{0,1\}\quad k\in \mathcal{K}, g\in [m].
\]
Here \(x_{(i,j),g}=1\) means that the job with processing time \(p_{ij}\) is assigned to machine \(g\),
whereas \(x_{i,g}=1\) means that the job with processing time \(f_i\) is assigned to machine \(g\).
Define the machine loads directly by
\[
  P_g
  =\sum_{(i,j)\in\mathcal{I}}x_{(i,j),g}p_{ij}
  +\sum_{i=1}^m x_{i,g}f_i
  \quad g\in [m].
\]



The model solved by $\verb|SCIP|$ in $\verb|Python 3.11.0|$ is:
\[
\boxed{
\begin{aligned}
\max\quad
  & r \\[0.4em]
\text{s.t.}\quad
  & \sum_{j=1}^m p_{ij}=1
      && i=1,\ldots,m,\\[0.4em]
  & \sum_{g=1}^m x_{kg}=1
      && k\in \mathcal{K},\\[0.4em]
  & P_g=\sum_{(i,j)\in\mathcal{I}}x_{(i,j),g}p_{ij}
      +\sum_{i=1}^m x_{i,g}f_i
      && g=1,\ldots,m,\\[0.4em]
  & P_g\ge \varepsilon
      && g=1,\ldots,m,\\[0.4em]
  & r\left(
      \sum_{g=1}^m P_g^2
      +(2\alpha-1)\left(\sum_{i=1}^m\sum_{j=1}^m p_{ij}^2
      +\sum_{i=1}^m f_i^2\right)
    \right)\\
  &\hspace{5em}\le
    \sum_{i=1}^m(f_i+1)^2
    +(2\alpha-1)\left(\sum_{i=1}^m\sum_{j=1}^m p_{ij}^2
    +\sum_{i=1}^m f_i^2\right),\\[0.4em]
  & \sum_{g=1}^m P_g^2
      +(2\alpha-1)\left(\sum_{i=1}^m\sum_{j=1}^m p_{ij}^2
      +\sum_{i=1}^m f_i^2\right)
      \ge \varepsilon,\\[0.4em]
  & 0\le p_{ij}\le 1
      && i,j=1,\ldots,m,\\[0.4em]
  & 0\le f_i\le 30
      && i=1,\ldots,m,\\[0.4em]
  & 0\le r\le 1+\frac{1}{6\alpha},\\[0.4em]
  & x_{kg}\in\{0,1\}
      && k\in \mathcal{K},\;g=1,\ldots,m.
\end{aligned}
}
\]
We maximize the instance-specific approximation ratio, represented by the
variable \(r\). As shown in Section~\ref{sec:Upper Bound}, the bound
$r \leq 1 + \frac{1}{6\alpha}$
does not exclude any feasible solution. Nevertheless, imposing this upper
bound improves the solver's efficiency.

The variable \(f_i\) denotes the processing time of the job that completes
last on machine \(i\) in the \(\mathrm{WSPT}\) schedule. An upper bound on
\(f_i\) is required to ensure that the solver searches over a bounded domain.
Based on preliminary experiments, we set this upper bound to \(30\). In all
computational results presented below, no
variable \(f_i\) attained its upper bound. We bound $P_g$ away from zero so as to 
avoid any stability issues.


\begin{center}
\begin{tikzpicture}
\begin{axis}[
  width=0.94\textwidth,
  height=0.62\textheight,
  xmode=log,
  log basis x=10,
  xmin=0.001111111111,
  xmax=0.5,
  ymin=1,
  ymax=17,
  xlabel={Parameter \(\alpha\)},
  ylabel={Objective value},
  xtick={
    0.5,
    0.1,
    0.02,
    0.005,
    0.0025,
    0.001111111111
  },
  xticklabels={
    \(\frac12\),
    \(\frac1{10}\),
    \(\frac1{50}\),
    \(\frac1{200}\),
    \(\frac1{400}\),
    \(\frac1{900}\)
  },
  grid=both,
  major grid style={gray!35},
  minor grid style={gray!15},
  legend style={
  at={(0.98,0.90)},
  anchor=north east,
  legend columns=1,
  draw=black,
  fill=white,
  fill opacity=0.9,
  text opacity=1,
  cells={anchor=west},
  inner sep=4pt,
  },
  tick label style={font=\small},
  label style={font=\large},
]

% Breakpoints a_k=1/D_k, where D_k=k^4-2k^3-2k^2+k+3.
\addplot[gray,dashed,forget plot] coordinates {(0.066666666667,1) (0.066666666667,17)};
\node[anchor=south,fill=white,inner sep=1pt] at (axis cs:0.066666666667,16.3)
  {\(\alpha_3\)};

\addplot[gray,dashed,forget plot] coordinates {(0.009708737864,1) (0.009708737864,17)};
\node[anchor=south,fill=white,inner sep=1pt] at (axis cs:0.009708737864,16.3)
  {\(\alpha_4\)};

\addplot[gray,dashed,forget plot] coordinates {(0.003003003003,1) (0.003003003003,17)};
\node[anchor=south,fill=white,inner sep=1pt] at (axis cs:0.003003003003,16.3)
  {\(\alpha_5\)};

\addplot[gray,dashed,forget plot] coordinates {(0.001248439451,1) (0.001248439451,17)};
\node[anchor=south,fill=white,inner sep=1pt] at (axis cs:0.001248439451,16.3)
  {\(\alpha_6\)};

% Piecewise branches of U_k(a).
\addplot[yellow, very thick,domain=0.066666666667:0.5,samples=100]
  {Uk(3,x)};
\addlegendentry{\(B_3(\alpha)\)}

\addplot[blue,very thick,domain=0.009708737864:0.066666666667,samples=100]
  {Uk(4,x)};
\addlegendentry{\(B_4(\alpha)\)}

\addplot[green,very thick,domain=0.003003003003:0.009708737864,samples=100]
  {Uk(5,x)};
\addlegendentry{\(B_5(\alpha)\)}

\addplot[purple,very thick,domain=0.001248439451:0.003003003003,samples=100]
  {Uk(6,x)};
\addlegendentry{\(B_6(\alpha)\)}

\addplot[orange,very thick,domain=0.001111111111:0.001248439451,samples=100]
  {Uk(7,x)};
\addlegendentry{\(B_7(\alpha)\)}

% At each \alpha, this is max over the computed m=3,...,7 objectives.
\addplot[
  only marks,
  mark=*,
  mark size=3.2pt,
  red!85!black,
  fill=red!65,
] coordinates {
  (0.5,          1.333333346656501)
  (0.1,          2.145683320593890)
  (0.02,         4.118401735799101)
  (0.005,        7.694623551818327)
  (0.0025,      10.642889269095093)
  (0.001111111111,15.644793715614036)
};
\addlegendentry{$\max_{\substack{m=3,\ldots,7}} \mathrm{MINLP}(\alpha)$}

% Visible vertical gaps between the curve and computed points.
\addplot[red!70!black,thick,forget plot] coordinates {
  (0.02,4.118401735799101)
  (0.02,4.125000000000000)
};
\addplot[red!70!black,thick,forget plot] coordinates {
  (0.005,7.694623551818327)
  (0.005,7.697804474378321)
};
\addplot[red!70!black,thick,forget plot] coordinates {
  (0.0025,10.642889269095093)
  (0.0025,10.643889586940004)
};
\addplot[red!70!black,thick,forget plot] coordinates {
  (0.001111111111,15.644793715614036)
  (0.001111111111,15.661998372367162)
};
\end{axis}
\end{tikzpicture}
\captionof{figure}{The MINLP is solved for each
\(\alpha \in \left\{\frac{1}{2}, \frac{1}{10}, \frac{1}{50},
\frac{1}{200}, \frac{1}{400}, \frac{1}{900}\right\}\)
and each \(m \in \{3,4,5,6,7\}\). Each run uses the parameter
settings \(0 \leq f_i \leq 30\), a SCIP node limit of \(20{,}000\),
and \(\varepsilon = 10^{-5}\). For each value of \(\alpha\), the
plotted value is the maximum objective value obtained across the
different numbers of machines. The resulting objective values are
then compared with the piecewise-defined analytical bounds \(B_i\).}
\label{fig:objective-comparison}
\end{center}


\begin{center}
\begin{tikzpicture}
\begin{groupplot}[
  group style={group size=2 by 2,horizontal sep=2.2cm,vertical sep=2.2cm},
  width=0.42\textwidth,
  height=0.22\textheight,
  grid=major,
  major grid style={gray!20},
  xlabel={\(\alpha\)},
  ylabel={Objective},
  tick label style={font=\scriptsize},
  label style={font=\small},
]
\nextgroupplot[
  title={Near \(\alpha=\frac1{50}\)},
  xmin=0.0195,xmax=0.0205,
  ymin=4.116,ymax=4.127,
]
\addplot[blue,very thick,domain=0.0195:0.0205,samples=50] {Uk(4,x)};
\addplot[only marks,mark=*,red,mark size=3pt] coordinates {(0.02,4.118401735799101)};

\nextgroupplot[
  title={Near \(\alpha=\frac1{200}\)},
  xmin=0.00485,xmax=0.00515,
  ymin=7.692,ymax=7.700,
]
\addplot[green,very thick,domain=0.00485:0.00515,samples=50] {Uk(5,x)};
\addplot[only marks,mark=*,red,mark size=3pt] coordinates {(0.005,7.694623551818327)};

\nextgroupplot[
  title={Near \(\alpha=\frac1{400}\)},
  xmin=0.00243,xmax=0.00257,
  ymin=10.6415,ymax=10.6452,
]
\addplot[purple,very thick,domain=0.00243:0.00257,samples=50] {Uk(6,x)};
\addplot[only marks,mark=*,red,mark size=3pt] coordinates {(0.0025,10.642889269095093)};

\nextgroupplot[
  title={Near \(\alpha=\frac1{900}\)},
  xmin=0.00109,xmax=0.001135,
  ymin=15.64,ymax=15.67,
]
\addplot[orange,very thick,domain=0.00109:0.001135,samples=50] {Uk(7,x)};
\addplot[only marks,mark=*,red,mark size=3pt] coordinates {(0.001111111111,15.644793715614036)};
\end{groupplot}
\end{tikzpicture}
\captionof{figure}{Enlarged views of the differences between the computed objective
values and the analytical lower bounds. As in
Figure~\ref{fig:objective-comparison}, the markers represent the results of
the MINLP evaluations, while the lines represent the analytical lower bounds
\(B_i\).}
\label{fig:objective-zoom}
\end{center}


Figure \ref{fig:objective-comparison} compares the piecewise defined lower bound we determined in 
Section \ref{sec:lower-bound} with sample values determined by the MINLP we just introduced. 
At the six sampled values of $\alpha$, the best computed objectives over
$m=3,\ldots,7$ practically agree with the analytical lower-bound curve.
The enlarged plot Figure \ref{fig:objective-zoom} shows how close the numerical results come to this curve for the four
smallest sampled values of $\alpha$, showing that the numerical search closely
approximates the ratios attained by the explicit instance family.

This agreement supports the family as a candidate for the worst-case ratio,
although the restriction to common starting times, the bound $f_i\leq30$, and
the finite $\verb|SCIP|$ node budget mean that these computations do not establish
global optimality.

\chapter{Outlook}\label{chap:outlook}
We set out to determine the approximation ratio \(\lambda_{\alpha}\) of the
\(\mathrm{WSPT}\)-algorithm for minimizing the total weighted
\(\alpha\)-completion time on identical parallel machines.
The cases $\alpha=0$ and $\alpha\in\Big[\frac12,1\Big]$ are solved by the
unboundedness result and the exact analysis presented in this thesis,
respectively. The open problem is therefore to determine the
approximation ratio for $\alpha\in\Big(0,\frac12\Big)$ by closing the remaining
gap between the upper and lower bounds. Resolving this gap is
also relevant to the stochastic motivation of this thesis: through the
transfer results of Jäger and Skutella, sharper deterministic bounds for
the $\alpha$-completion-time problem may yield improved performance
guarantees for $\mathrm{WSEPT}$.\medskip


Our computational analysis shows that the approximation ratio for $\alpha\in\Big(0,\frac12\Big)$ is likely to be close to the lower bound family of instances introduced in this thesis.
It is this observation that motivates the following conjecture.
\begin{conjecture}
\[
\lambda_{\alpha} =
\begin{cases}
B_3(\alpha),
& \displaystyle \alpha_3\leq\alpha\leq\frac12,\\[6pt]
B_{m+1}(\alpha),
& \alpha_{m+1}\leq\alpha<\alpha_{m},\quad m\geq 3,
\end{cases}
\qquad
\alpha_m=\frac{1}{m^4-2m^3-2m^2+m+3}.
\]
\end{conjecture}
The instances that attain the ratio of $B_{m}(\alpha)$ are of the following form:

\begin{figure}[ht]
\centering
\begin{subfigure}[t]{0.47\textwidth}
\centering
\begin{tikzpicture}[
    x=0.7cm,
    y=0.7cm,
    job/.style={
        draw,
        thick,
        fill=white,
        font=\small
    },
    machine/.style={font=\small}
]

% ------------------------------------------------
% Template idea:
% row      = displayed vertical row in the picture
% machine  = actual machine label shown on the left
%
% Use \machinerow{row}{machine label}
% Use \dotsrow{row}
% Use \jobblock{row}{start}{length}{job label}
% Use \jobblock[jobhighlight]{row}{start}{length}{job label} for shading
% ------------------------------------------------

% Displayed machine rows
\machinerow{1}{1}
\dotsrow{2}
\machinerow{3}{m-1}
\machinerow{4}{m}


% Example jobs
\jobblock{1}{0}{3}{1-\frac{1}{m}}

\jobblock{3}{0}{3}{1-\frac{1}{m}}

\jobblock{4}{0}{1}{\frac{1}{m}}
\jobblock{4}{1}{1}{\dots}
\jobblock{4}{2}{1}{\frac{1}{m}}
\jobblock{4}{3}{8}{p_{\alpha}}
\end{tikzpicture}

\caption{Schedule $\sigma^{*}(I)$.}
\end{subfigure}
\hfill
\begin{subfigure}[t]{0.47\textwidth}
\centering
\begin{tikzpicture}[
    x=0.7cm,
    y=0.7cm,
    job/.style={
        draw,
        thick,
        fill=white,
        font=\small
    },
    machine/.style={font=\small}
]

% ------------------------------------------------
% Template idea:
% row      = displayed vertical row in the picture
% machine  = actual machine label shown on the left
%
% Use \machinerow{row}{machine label}
% Use \dotsrow{row}
% Use \jobblock{row}{start}{length}{job label}
% Use \jobblock[jobhighlight]{row}{start}{length}{job label} for shading
% ------------------------------------------------

% Displayed machine rows
\machinerow{1}{1}
\dotsrow{2}
\machinerow{3}{m-1}
\machinerow{4}{m}


% Example jobs
\jobblock{1}{0}{2}{\frac{1}{m}}
\jobblock{1}{2}{3}{1-\frac{1}{m}}

\jobblock{3}{0}{2}{\frac{1}{m}}
\jobblock{3}{2}{3}{1-\frac{1}{m}}

\jobblock{4}{0}{8}{p_{\alpha}}



\end{tikzpicture}

\caption{Schedule $\sigma^{\mathrm{WSPT}}(I)$.}
\end{subfigure}
\caption{Schedules for $\alpha=0$.}
\label{fig:insjnkntance-43-ratio}
\end{figure}
One possible approach to proving this conjecture is to establish further
structural reductions. In particular, it would be useful to identify a
\(\lambda_{\alpha}\)-dominant class in which every instance has at most
\[
    n\leq 2(m-1)+1=2m-1
\]
jobs. This bound coincides with the number of jobs in the instances attaining
\(B_m(\alpha)\) and would be an improvement over the current best bound of \(n\leq 3m-2\) jobs. 

A further useful reduction would be to show that there is a dominant class in which 
every instance admits a cost-maximizing list schedule
\(\sigma^{W}(I)\) in which \(m-1\) machines are singleton machines, meaning
that each of these machines processes exactly one job.

The two upper bounds established in this thesis use only a limited subset of
the known structural properties of \(\lambda_{\alpha}\)-dominant instances.
Consequently, the existing upper-bound analysis may leave relevant structural
information unused. Incorporating these may therefore yield sharper upper bounds.
\appendix

\chapter{Detailed Calculations}
\label{app:detailed-calculations}
We provide the proofs and calculations that were omitted in the main text.
\phantomsection
\label{app:proof-queyranne}
\begin{proof}[Proof of Lemma~\ref{corr:queyranne}]
Since we deal only with a single machine, the maximal machine load is the sum of all processing times, i.e., $\machLoadMax=\sum_{j=1}^n p_j$.
For the following calculation, we assume without loss of generality that the jobs are scheduled
in order of increasing index.
\[
\begin{aligned}
        \sum_{j=1}^n p_jC_j^{\alpha}(\sigma(I))
        &=
        \sum_{j=1}^n p_j\left(\sum_{k=1}^{j-1} p_k +\alpha p_j \right)                       \\
        &=
        \sum_{j=1}^n p_jC_{j-1}(\sigma(I))+\alpha\left(\sum_{j=1}^n p_j^2  \right)                   \\
        &=
        \sum_{j=1}^n p_j(C_j(\sigma(I))-p_j)+\alpha\left(\sum_{j=1}^n p_j^2  \right)                 \\
        &=
        \sum_{j=1}^n p_jC_j(\sigma(I))+(\alpha-1)\left(\sum_{j=1}^n p_j^2  \right)                      \\
        &=
        \frac12\left((\machLoadMax)^2+\sum_{j=1}^n p_j^2\right)+(\alpha-1)\left(\sum_{j=1}^n p_j^2  \right)\\
        &=
        \frac12\left((\machLoadMax)^2+(2\alpha-1)\left(\sum_{j=1}^n p_j^2  \right) \right)
\end{aligned}
\]
Where, as in \cite{Schwiegelshohn:2011}, we make use of a result from \cite{queyranne1993structure} for the fifth equality of this calculation.
\end{proof}


\phantomsection
\label{app:proof-cases}
\begin{proof}[Proof of Result~\ref{Res:Case-Comp}]

First, observe that
\begin{align*}
\left(2\sqrt{\alpha}+\sqrt{2}\right)^2=
2\left(1+\sqrt{2\alpha}\right)^2.
\end{align*}

Therefore,
\begin{align*}
\lambda_{\alpha}(I_{2})
&=
\frac{1}{
1-\dfrac{2}{\left(2\sqrt{\alpha}+\sqrt{2}\right)^2}
}
=
\frac{1}{
1-\dfrac{2}{2\left(1+\sqrt{2\alpha}\right)^2}
}
=
\frac{\left(1+\sqrt{2\alpha}\right)^2}
{\left(1+\sqrt{2\alpha}\right)^2-1}.
\end{align*}
We compute the difference:
\begin{align*}
\lambda_{\alpha}(I_{2})-\lambda_{\alpha}(I_{1})
&=
\frac{\left(1+\sqrt{2\alpha}\right)^2}
{\left(1+\sqrt{2\alpha}\right)^2-1}
-
\left(
1+\frac{1}{\left(1+\sqrt{2\alpha}\right)^2}
\right)
=
\frac{\left(1+\sqrt{2\alpha}\right)^2}
{\left(1+\sqrt{2\alpha}\right)^2-1}
-
\frac{\left(1+\sqrt{2\alpha}\right)^2+1}
{\left(1+\sqrt{2\alpha}\right)^2}\\
&=
\frac{1}{
\left(1+\sqrt{2\alpha}\right)^2
\left(
\left(1+\sqrt{2\alpha}\right)^2-1
\right)
}.
\end{align*}
For every \(\alpha\in\Big[\frac12,1\Big]\), the denominator is strictly positive.
Hence, \(\lambda_{\alpha}(I_2)>\lambda_{\alpha}(I_1)\) holds as claimed.
\end{proof}

\label{app:proofmtwo}
\begin{proof}[\hypertarget{proof:mtwo}{Proof of Lemma~\ref{lem:mtwo}}]
Let
\[
Q := p_1^2+p_3^2+p_2^2,
\]
and let \(\alpha\) be a parameter with \(0 \le \alpha \le \tfrac12\).
We want to determine
\[
\lambda^{2}_{\alpha}=\max_{p_1,p_2,p_3}\frac{{(p_1+p_2)}^2+{(p_3)}^2+(2\alpha-1)Q}{{(p_1+p_3)}^2+{(p_2)}^2+(2\alpha-1)Q}
\]
where \(p_1,p_2,p_3 \ge 0\), subject to the following constraints:
\begin{align}
&p_3 \ge p_1\\
&p_3 \le p_2\\
&p_1 \le p_2.
\end{align}
These must hold so that the schedule actually is a list schedule.
We know from Lemma \ref{corr:unification} and Theorem \ref{corr: singletonjob} that it is sufficient to consider 
instances with only three jobs for $m=2$ to determine $\lambda_{\alpha}$ or, in this case, $\lambda^2_{\alpha}$.
\begin{figure}[htbp]
\centering
\setlength{\fboxsep}{6pt}
\fbox{%
\begin{minipage}{0.94\textwidth}
    \begin{minipage}[t]{0.47\linewidth}
        \centering
        \scalebox{0.8}{\schedulemtwoone}
        \par\smallskip
        \small(a) \(\sigma^{*}(I)\)
    \end{minipage}
    \hfill
    \begin{minipage}[t]{0.47\linewidth}
        \centering
        \scalebox{0.8}{\schedulemtwotwo}
        \par\smallskip
        \small(b) \(\sigma^{W}(I)\)
    \end{minipage}
\end{minipage}%
}
\caption{General schedules for $m=2$.}
\label{figure-scawdne}
\end{figure}

We can rewrite the numerator
\begin{align*}
&{(p_1+p_2)}^2+{(p_3)}^2+(2\alpha-1)Q\\
&= (p_1^2+p_2^2+2p_1p_2)+(p_3^2)+(2\alpha-1)Q\\
&= 2p_1p_2 + 2\alpha Q.
\end{align*}
Similarly, the denominator becomes
\[
{(p_3+p_1)}^2+{(p_2)}^2+(2\alpha-1)Q = 2p_1p_3+2\alpha Q.
\]
Therefore
\[
\lambda_{\alpha}(I)
=
\frac{\alpha Q+p_1p_2}{\alpha Q+p_1p_3}=1+\frac{p_1(p_2-p_3)}{\alpha Q+p_1p_3}.
\]

\(\lambda_{\alpha}(I) > 1\) only if \(p_1(p_2-p_3)>0\).
So we may restrict our search to instances for which \(p_2>p_3\).

The restrictions read: \(p_1 \le p_3 < p_2\)

\(\lambda_{\alpha}(I)\) is monotonically decreasing in \(p_3\). 
We thus set \(p_3=p_1\) as this is the smallest feasible value of $p_3$. 

$p_3=p_1$ simplifies our objective: 
\[
\max_{p_1,p_2\ge 0,\;p_2\ge p_1}
1+\frac{p_1(p_2-p_1)}{\alpha(2p_1^2+p_2^2)+p_1^2},
\qquad 0<\alpha\le \frac12.
\]

Set
\[
x=\frac{p_2}{p_1}.
\]
For \(p_1>0\), the restriction \(p_2\ge p_1\) becomes \(x\ge 1\), and
\[
1+\frac{p_1(p_2-p_1)}{\alpha(2p_1^2+p_2^2)+p_1^2}
=
1+\frac{\frac{1}{p_1}(p_2-p_1)}{\frac{1}{p_1^2}(\alpha(2p_1^2+p_2^2)+p_1^2)}
=
1+\frac{\frac{p_2}{p_1}-1}{\alpha(2+(\frac{p_2}{p_1})^2)+1}
=
1+\frac{x-1}{\alpha(2+x^2)+1}.
\]
Thus we maximize
\[
f(x)=1+\frac{x-1}{\alpha x^2+2\alpha+1},
\qquad x\ge 1.
\]

The derivative is
\[
f'(x)
=
\frac{\alpha x^2+2\alpha+1-2\alpha x(x-1)}
{{(\alpha x^2+2\alpha+1)}^2}
=
\frac{-\alpha x^2+2\alpha x+2\alpha+1}
{{(\alpha x^2+2\alpha+1)}^2}.
\]

The denominator is positive for \(0<\alpha\le \frac12\), so the critical points are determined by
\[
-\alpha x^2+2\alpha x+2\alpha+1=0.
\]

Which has the solutions
\[
x_{\alpha}
=
1\pm \sqrt{3+\frac1\alpha}.
\]
Since \(x\ge 1\), only
\[
x_{\alpha}^*=1+\sqrt{3+\frac1\alpha}
\]
is relevant.
This is a maximum, since  
\[
f'(x)=-\alpha x^2+2\alpha x+2\alpha+1.
\]
is a downward-opened parabola.


Thus, the maximum is attained for all \(p_1>0\) with
\[
\boxed{
p_2=p_1\left(1+\sqrt{3+\frac1\alpha}\right) \quad \text{and} \quad p_3=p_1
}
\]
and the maximum value is
\[
\boxed{
f(x_{\alpha}^*)
=
1+\frac{1}{2\alpha\left(1+\sqrt{3+\frac1\alpha}\right)}
}.
\]
\begin{figure}[htbp]
\centering
\setlength{\fboxsep}{6pt}
\fbox{%
\begin{minipage}{0.94\textwidth}
    \begin{minipage}[t]{0.47\linewidth}
        \centering
        \scalebox{0.8}{\scheduleappendixmm}
        \par\smallskip
        \small(a) \(\sigma^{*}(I_1)\)
    \end{minipage}
    \hfill
    \begin{minipage}[t]{0.47\linewidth}
        \centering
        \scalebox{0.8}{\scheduleappendixm}
        \par\smallskip
        \small(b) \(\sigma^{W}(I)\)
    \end{minipage}
\end{minipage}%
}
\caption{$\lambda^2_{\alpha}$-maximizing schedules.}
\end{figure}

If \(\alpha=0\), then
\[
\lambda_{\alpha}(I)=\frac{p_1(p_2-p_1)}{\alpha Q +p_1^2}=
\frac{p_1(p_2-p_1)}{p_1^2}
=
\frac{p_2}{p_1}-1,
\]
which is unbounded for \(p_2/p_1\to\infty\). Thus for \(\alpha=0\), no finite maximum exists.

Consequently, the approximation ratio is
\[
\boxed{
\lambda^2_{\alpha}
=
\begin{cases}
1+\dfrac{1}{2\alpha\left(1+\sqrt{3+\frac1\alpha}\right)},
    & 0<\alpha\le \dfrac12,\\[1.2ex]
+\infty,
    & \alpha=0.
\end{cases}
}.
\]
\end{proof}



\label{app:proof-mthree}
\begin{proof}[\hypertarget{proof:mthree}{Proof of Lemma~\ref{lem: mthree}}]
We know from Theorems \ref{thm:job-bound} and \ref{corr: singletonjob}, that we may only consider instances
with $n\le 3m-2=7$ jobs. Furthermore, we know that there is a singleton in $\sigma^*(I)$ and that 
there are at most $m=3$ jobs scheduled on any other machine in $\sigma^*(I)$.
The remaining class of instances can be expressed as follows: 

Let
\[
0\le \alpha\le \frac12 \quad\text{and}\quad
p_1,p_2,p_3,p_4,p_5,p_6,p_7>0.
\]
Set
\[
S:=p_1^2+p_2^2+p_3^2+p_4^2+p_5^2+p_6^2+p_7^2 \qquad\text{and}\qquad \mathcal P
=
\{(p_1,\ldots,p_7):p_i\geq0\text{ for all }i\in[7]\},
\]

Then,

\[
\lambda^3_{\alpha}
=
\max_{\mathcal{P}}
\frac{
(p_1+p_4)^2+(p_2+p_5)^2+(p_3+p_6+p_7)^2+(2\alpha-1)S
}{
(p_1+p_2+p_3)^2+(p_4+p_5+p_6)^2+p_7^2+(2\alpha-1)S
}
\]
subject to
\begin{align}
p_1+p_2+p_3&\ge p_4+p_5,
    \label{con:m3-load-123}\\
p_4+p_5+p_6&\ge p_1+p_2,
    \label{con:m3-load-456}\\
p_7&\ge p_1+p_2,
    \label{con:m3-p7-12}\\
p_7&\ge p_4+p_5,
    \label{con:m3-p7-45}\\
p_1+p_4&\ge p_3+p_6,
    \label{con:m3-p14-36}\\
p_2+p_5&\ge p_3+p_6,
    \label{con:m3-p25-36}\\
p_1+p_4&\ge p_2,
    \label{con:m3-p14-2}\\
p_2+p_5&\ge p_1,
    \label{con:m3-p25-1}\\
p_3+p_6+p_7&\ge p_1,
    \label{con:m3-p367-1}\\
p_3+p_6+p_7&\ge p_2.
    \label{con:m3-p367-2}
\end{align}

Where equations \eqref{con:m3-load-123} - \eqref{con:m3-p7-45} are sufficient conditions for 
$\sigma^{*}(I)$ to be a list schedule and \eqref{con:m3-p14-36} - \eqref{con:m3-p367-2} are sufficient 
for $\sigma^{W}(I)$ to be a list schedule.


\begin{figure}[htbp]
\centering
\setlength{\fboxsep}{6pt}
\fbox{%
\begin{minipage}{0.94\textwidth}
    \begin{minipage}[t]{0.47\linewidth}
        \centering
        \scalebox{0.8}{\schedulethreeone}
        \par\smallskip
        \small(a) \(\sigma^{*}(I)\)
    \end{minipage}
    \hfill
    \begin{minipage}[t]{0.47\linewidth}
        \centering
        \scalebox{0.8}{\schedulethreetwo}
        \par\smallskip
        \small(b) \(\sigma^{W}(I)\)
    \end{minipage}
\end{minipage}%
}
\caption{General schedules for $m=3$.}
\label{figure-scawd}
\end{figure}


We will normalize the problem by setting
$p_7=1$
but keep the notation $p_7$ for the moment as this makes
our calculations easier to follow.

Constraints \eqref{con:m3-p367-1} and \eqref{con:m3-p367-2} are redundant. Indeed, since
\[
p_7\ge p_1+p_2
\]
and \(p_1,p_2>0\), we have \(p_7\ge p_1\) and \(p_7\ge p_2\). Hence,
\[
p_3+p_6+p_7\ge p_1\quad \text{and} \quad p_3+p_6+p_7\ge p_2
\]
hold automatically.



Define
\[
L:=p_1p_2+p_1p_3+p_2p_3+p_4p_5+p_4p_6+p_5p_6
\quad \text{and}\quad
Q:=p_1p_4+p_2p_5+p_3p_6+p_7(p_3+p_6).
\]
Expanding the denominator of $\lambda_{\alpha}(I)$ gives
\[
\begin{aligned}
& (p_1+p_2+p_3)^2+(p_4+p_5+p_6)^2+p_7^2+(2\alpha-1)S\\
&=
2\alpha S+2(p_1p_2+p_1p_3+p_2p_3+p_4p_5+p_4p_6+p_5p_6)\\
&=
2(\alpha S+L).
\end{aligned}
\]
Similarly for the numerator,
\[
\begin{aligned}
& (p_1+p_4)^2+(p_2+p_5)^2+(p_3+p_6+p_7)^2+(2\alpha-1)S\\
&=
2\alpha S+2(p_1p_4+p_2p_5+p_3p_6+p_7(p_3+p_6))\\
&=
2(\alpha S+Q).
\end{aligned}
\]
Therefore the problem is equivalent to maximizing
\[
\frac{\alpha S+Q}{\alpha S+L}=1+
\frac{Q-L}{\alpha S+L}.
\]

Only one term in \(Q\) involves \(p_7\), namely
\[
p_7(p_3+p_6).
\]
The denominator contains the compensating term
\[
\alpha p_7^2.
\]
Thus, increasing \(p_7\) cannot make the ratio
arbitrarily large.

Let
\[
x\coloneqq p_3+p_6
\]
and rewrite the constraints \ref{con:m3-p14-36} and \ref{con:m3-p25-36} as
\[
p_1+p_4\ge x,\qquad p_2+p_5\ge x.
\]
We will show next that, \(p_1+p_4=p_2+p_5=x\) and \(p_3=p_6=x/2\) must hold for the maximum.

We now record the algebraic reduction leading to the active boundary.

We begin by introducing the sums
\[
u \coloneqq p_1+p_4,\qquad v \coloneqq p_2+p_5,\qquad x\coloneqq p_3+p_6.
\]
The constraints include
\[
u\ge x,\qquad v\ge x.
\]
Also define
\[
h\coloneqq p_1+p_2,\qquad d \coloneqq p_4+p_5.
\]
The constraints
\[
p_7\ge p_1+p_2,\qquad p_7\ge p_4+p_5
\]
become
\[
p_7\ge h,\qquad p_7\ge d.
\]
Finally,
\[
p_1+p_4\ge p_2,\qquad p_2+p_5\ge p_1
\]
become
\[
u\ge p_2,\qquad v\ge p_1.
\]

For fixed \(p_7\) and \(x=p_3+p_6\), expanding \(Q-L\) in the variables
\(u,v,h,p_3\) gives
\[
\begin{aligned}
Q-L
=p_7x-p_3^2+p_3(u+v+x-2h)-h^2
+h(u+v+x)-uv-x(u+v).
\end{aligned}
\]
We now set
\[
p_1+p_4=x,\qquad p_2+p_5=x.
\]
We will later see in a perturbation type argument in the variables $r$ and $s$, that choosing no slack in 
the constraints is indeed optimal. 
With this choice, we can analyze the remaining constraint \(p_2+p_5\ge p_1\). Choose the
maximizing "split"
\[
p_3=p_6=\frac{x}{2},
\]
which is the maximal feasible choice for fixed \(x\), since it maximizes
\(p_3p_6\) (which appears in the numerator of $\frac{\alpha S + Q}{\alpha S+L}$) and minimizes \(p_3^2+p_6^2\) (which appears in the denominator of $\frac{\alpha S + Q}{\alpha S+L}$).

Put
\[
y=p_1+p_2.
\]
Then, inserting gives
\[
Q-L
=
p_7x-\frac74x^2+2xy-y^2.
\]
For fixed \(x\) and \(p_7\), this quadratic in \(y\) is maximized at
\begin{equation}\label{Wolfram7}
y=x. 
\end{equation}

Thus, the maximizing point satisfies
\[
p_1+p_2=x.
\]
Here, the denominator is
\[
\alpha S+L
=
\alpha p_7^2
+\left(1+\frac52\alpha\right)x^2
+(2-4\alpha)p_1p_2.
\]
Since \(0<\alpha\le \frac12\), the coefficient \(2-4\alpha\) is
nonnegative. Hence the denominator is minimized by
\[
p_1p_2=0.
\]
Consequently, either
\[
(p_1,p_2)=(x,0)
\qquad\text{or}\qquad
(p_1,p_2)=(0,x).
\]
The two cases are symmetric. We choose the first one, which implies
\[
p_5=x,\qquad p_2+p_5=x=p_1.
\]
Thus, the constraint
\(p_2+p_5\ge p_1\) is tight.

Finally, we need to check that this does not benefit from slack in
\(u\ge x\) or \(v\ge x\). We write
\[
u=x+r,\qquad v=x+s,\qquad r,s\ge0.
\]
This gives
\[
(p_1,p_2,p_3,p_4,p_5,p_6,p_7)
=
\left(x,s,\frac{x}{2},r,x,\frac{x}{2},p_7\right),
\]
so that
\[
p_1+p_4=x+r,\qquad
p_2+p_5=x+s,\qquad
p_3+p_6=x.
\]
For this family the excess numerator is
\[
Q-L
=
p_7x-\frac34x^2-\frac{x}{2}(r+s),
\]
whereas the denominator is
\[
\alpha S+L
=
\alpha p_7^2
+\left(1+\frac52\alpha\right)x^2
+\frac{3x}{2}(r+s)
+\alpha(r^2+s^2).
\]
Thus every positive slack \(r\) or \(s\) subtracts from the numerator
and adds to the denominator. Consequently, for fixed \(x\) and \(p_7\),
the largest value in this slack family is attained at
\[
r=s=0.
\]
This proves that setting $p_1+p_4=x$ and $p_2+p_5=x$ is indeed justified.

Substituting the active equalities gives
\[
(p_1,p_2,p_3,p_4,p_5,p_6,p_7)
=
\left(x,0,\frac{x}{2},0,x,\frac{x}{2},p_7\right).
\]
We normalize \(p_7\) to get
\[
\left(t,0,\frac{t}{2},0,t,\frac{t}{2},1\right).
\]
On this boundary,
\[
Q-L=t-\frac34t^2
\]
and
\[
\alpha S+L
=
\alpha+\left(1+\frac52\alpha\right)t^2.
\]
Therefore every feasible point is bounded above by the one-variable
problem
\[
\sup_{t>0}
\frac{t-\frac34t^2}
{\alpha+\left(1+\frac52\alpha\right)t^2}.
\]

We continue by maximizing the one-variable function $f_{\alpha}$
Define
\[
f_\alpha(t)
\coloneqq
\frac{t-\frac34t^2}{\alpha+\left(1+\frac52\alpha\right)t^2},
\qquad t>0.
\]
Set
\[
\beta=1+\frac52\alpha,
\]
which leads to
\[
f_\alpha(t)=\frac{t-\frac34t^2}{\alpha+\beta t^2}.
\]
Differentiate:
\begin{align}
f_\alpha'(t)
=
\frac{
\left(1-\frac32t\right)(\alpha+\beta t^2)
-
\left(t-\frac34t^2\right)(2\beta t)
}{
(\alpha+\beta t^2)^2
} \label{Wolfram3}.
\end{align}
The numerator simplifies to
\[
\alpha-\frac32\alpha t-\beta t^2.
\]
Thus the critical point satisfies
\[
\beta t^2+\frac32\alpha t-\alpha=0.
\]
The positive root is
\begin{align}
\label{Wolfram4}
t_\alpha
=
\frac{
-\frac32\alpha+\sqrt{\frac94\alpha^2+4\alpha\beta}
}{
2\beta
}.
\end{align}
Since
\[
\beta=1+\frac52\alpha,
\]
we have
\[
\frac94\alpha^2+4\alpha\beta
=
4\alpha+\frac{49}{4}\alpha^2
=
\frac14\alpha(16+49\alpha).
\]
Therefore
\[
t_\alpha
=
\frac{\sqrt{\alpha(16+49\alpha)}-3\alpha}{4+10\alpha}.
\]

So the supremum over the reduced one-variable problem is

\[
\lambda_\alpha(t_\alpha)
= 1+
\frac{2}{\sqrt{\alpha(16+49\alpha)}+3\alpha}.
\]


The maximizing boundary point is only infeasible since in this instance we enforce 
there to be seven nonzero jobs. 
Therefore the expression above is only a supremum which can be approached by the family

\[
\left(
t_\alpha,\varepsilon,\frac{t_\alpha}{2},
\varepsilon,t_\alpha,\frac{t_\alpha}{2},1
\right)
\]
and $\varepsilon \rightarrow 0$.


Taking the limit and deleting all jobs with processing time zero leaves us with a 
feasible instance that achieves the already established result for $\lambda^3_{\alpha}$ (see Section \ref{sec:FixedmachineNumber}).
\[
\left(
t_\alpha,\frac{t_\alpha}{2},t_\alpha,\frac{t_\alpha}{2},1
\right).
\]
The worst and optimal schedules corresponding to the job set we just determined,
are depicted in Figure \ref{figure-sometjing}.

\begin{figure}[htbp]
\centering
\setlength{\fboxsep}{6pt}
\fbox{%
\begin{minipage}{0.94\textwidth}
    \begin{minipage}[t]{0.47\linewidth}
        \centering
        \scalebox{0.8}{\schedulethreemaxone}
        \par\smallskip
        \small(a) \(\sigma^{*}(I)\)
    \end{minipage}
    \hfill
    \begin{minipage}[t]{0.47\linewidth}
        \centering
        \scalebox{0.8}{\schedulethreemaxtwo}
        \par\smallskip
        \small(b) \(\sigma^{W}(I)\)
    \end{minipage}
\end{minipage}%
}
\caption{$\lambda_{\alpha}^3$-maximizing schedules.}
\label{figure-sometjing}
\end{figure}
\end{proof}



\label{app:proof-JobNumber}
\begin{proof}[\hypertarget{proof:JobNumber}{Proof of Theorem~\ref{thm:job-bound}}]
\begin{equation} \label{app:ratio-calculation}
\begin{aligned}
&
\frac{\gamma_{\alpha}(\sigma^W(I(\varepsilon)))}{\gamma_{\alpha}(\sigma^{*}(I(\varepsilon)))}-
\frac{\gamma_{\alpha}(\sigma^W(I))}{\gamma_{\alpha}(\sigma^{*}(I))}
=
\frac{A^W+(2\alpha-1)Q(I(\varepsilon))}
     {A^*+(2\alpha-1)Q(I(\varepsilon))}
-
\frac{A^W+(2\alpha-1)Q}
     {A^*+(2\alpha-1)Q}
\\[1mm]
&=
\frac{
\bigl(A^W+(2\alpha-1)Q(I(\varepsilon))\bigr)
\bigl(A^*+(2\alpha-1)Q\bigr)}
{\bigl(A^*+(2\alpha-1)Q(I(\varepsilon))\bigr)
 \bigl(A^*+(2\alpha-1)Q\bigr)}
-
\frac{
\bigl(A^W+(2\alpha-1)Q\bigr)
\bigl(A^*+(2\alpha-1)Q(I(\varepsilon))\bigr)}
{\bigl(A^*+(2\alpha-1)Q(I(\varepsilon))\bigr)
 \bigl(A^*+(2\alpha-1)Q\bigr)}
\\[1mm]
&=
\frac{
(2\alpha-1)A^WQ
+(2\alpha-1)A^* Q(I(\varepsilon))
-(2\alpha-1)A^W Q(I(\varepsilon))
-(2\alpha-1)A^*Q}
{\bigl(A^* +(2\alpha-1)Q(I(\varepsilon))\bigr)
 \bigl(A^* +(2\alpha-1)Q\bigr)}
\\[1mm]
&=
\frac{
(2\alpha-1)
\bigl(Q(I(\varepsilon))-Q\bigr)
\bigl(A^* - A^W\bigr)}
{\bigl(A^* +(2\alpha-1)Q(I(\varepsilon))\bigr)
 \bigl(A^* +(2\alpha-1)Q\bigr)}
\\[1mm]
&=
\frac{
(1-2\alpha)
\bigl(Q(I(\varepsilon))-Q\bigr)
\bigl(A^W-A^*\bigr)}
{\bigl(A^*+(2\alpha-1)Q(I(\varepsilon))\bigr)
 \bigl(A^*+(2\alpha-1)Q\bigr)} > 0.
\end{aligned}
\end{equation}

The strict inequality holds because
\[
1-2\alpha>0,\qquad
Q(I(\varepsilon))-Q>0,\qquad
A^W-A^*>0,
\]
and
\[
A^*+(2\alpha-1)Q(I(\varepsilon))>0,
\qquad
A^W+(2\alpha-1)Q(I(\varepsilon))>0.
\]
\end{proof}

\label{app:proof-LB}
\begin{proof}[\hypertarget{proof:LB}{Proof of Lemma~\ref{lem: LB}}]

Recall that we want to maximize 
\[
\lambda_{\alpha}(m,u,p)=\frac{
(m-u)\left(1-\frac um\right)^2
 +u\left(p+1-\frac um\right)^2
 +(2\alpha-1)B
}{
(m-u)+up^2+(2\alpha-1)B
},
\]
where
\[
\begin{aligned}
B&=\sum_{j=1}^{m-ku-1}2\left(\frac{j}{m}\right)^2+\left((k+1)u-m+1\right)\left(\frac{m-ku}{m}\right)^2
+up^2
+(m-u)\left(1-\frac um\right)^2
+(k-1)\frac{u^3}{m^2}\\
&=\frac{(m-ku-1)(m-ku)(2m-2ku-1)}{3m^2}  \\
&+\left((k+1)u-m+1\right)\left(\frac{m-ku}{m}\right)^2
+up^2
+(m-u)\left(1-\frac um\right)^2
+(k-1)\frac{u^3}{m^2}.
\end{aligned}
\]
Set
\[
        x=\frac{u}{m}, \qquad r=1-kx=\frac{m-ku}{m}.
\]
The constraint on $\frac{u}{m}$ becomes
\[
        \frac{1}{k+1}\leq x\leq \frac1k,
        \qquad 0\leq r = 1-kx \le 1-\frac{k}{k+1} = \frac{1}{k+1}.
\]
Also set
\[
        b=1-2\alpha.
\]
Then \(0\leq b\leq 1\), and
\[
        2\alpha-1=-b.
\]

We first separate the \(p^2\)-term in B. Write
\[
        B=mxp^2+R,
\]
where
\[
\begin{aligned}
R={}&
\frac{(m-ku-1)(m-ku)(2m-2ku-1)}{3m^2}  \\
&+\left((k+1)u-m+1\right)\left(\frac{m-ku}{m}\right)^2
+(m-u)\left(1-\frac um\right)^2
+(k-1)\frac{u^3}{m^2}.
\end{aligned}
\]

Since \(m-ku=mr\), the first summand satisfies
\[
\begin{aligned}
\frac{(m-ku-1)(m-ku)(2m-2ku-1)}{3m^2}
&=
\frac{(mr-1)(mr)(2mr-1)}{3m^2}.
\end{aligned}
\]
Dividing by \(m\),
\[
\begin{aligned}
\frac{(mr-1)(mr)(2mr-1)}{3m^3}
&=
\frac{2r^3}{3}-\frac{r^2}{m}+\frac{r}{3m^2}.
\end{aligned}
\]
The second summand gives
\[
\begin{aligned}
\frac{1}{m}
\left((k+1)u-m+1\right)\left(\frac{m-ku}{m}\right)^2
&=
\left((k+1)x-1+\frac1m\right)r^2.
\end{aligned}
\]
Using \(1=kx+r\), we have
\[
       \left((k+1)x-1+\frac1m\right)r^2= \left(x-r+\frac1m\right)r^2.
\]
If we divide the third summand by $m$, it satisfies
\[
\frac{m-u}{m}\left(1-\frac{u}{m}\right)^2= \left(1-\frac{u}{m}\right)^3.
\]
Combining the observations for each of the summands
\[
\begin{aligned}
\frac{R}{m}
&=
\left(\frac{2r^3}{3}-\frac{r^2}{m}+\frac{r}{3m^2}\right)
+
\left(x-r+\frac1m\right)r^2
+(1-x)^3
+(k-1)x^3 \\
&=
\frac{2r^3}{3}-\frac{r^2}{m}+\frac{r}{3m^2}
+xr^2-r^3+\frac{r^2}{m}
+(1-x)^3
+(k-1)x^3 \\
&=
xr^2-\frac{r^3}{3}
+(1-x)^3
+(k-1)x^3
+\frac{r}{3m^2}.
\end{aligned}
\]
Denote this quantity by
\[
        \rho(x,m)
        =
        \frac{R}{m}
        =
        xr^2-\frac{r^3}{3}
        +(1-x)^3
        +(k-1)x^3
        +\frac{r}{3m^2}.
\]

Subtracting the numerator (N) and denominator (D) of the competitive ratio $\lambda_{\alpha}(u,m,p)$ gives

\[
\begin{aligned}
N-D
={}&
(m-u)\left(1-\frac um\right)^2
+u\left(p+1-\frac um\right)^2
-(m-u)-up^2.
\end{aligned}
\]
Substituting \(x=\frac{u}{m}\), this becomes
\[
\begin{aligned}
N-D
&=
m(1-x)^3
 +mx\bigl(p+1-x\bigr)^2
 -m(1-x)-mxp^2 \\
&=
mx(1-x)(2p-1).
\end{aligned}
\]
Thus,
\[
\lambda_{\alpha}(x,m,p)=\frac{D+(N-D)}{D}=1+\frac{mx(1-x)(2p-1)}{D}.
\]

Where we associate the class of instances we are maximizing over with the three-tuple $(x,m,p)$
and thus replace it in the argument of $\lambda$

The denominator is
\[
\begin{aligned}
D
&=
m(1-x)+mxp^2-b(mx p^2+R) \\
&=
m\left(2\alpha xp^2+(1-x)-b\rho(x,m)\right).
\end{aligned}
\]
Define
\[
        c(x,m,\alpha)=(1-x)-(1-2\alpha)\rho(x,m).
\]
Then, for \(\alpha>0\),
\[
        \lambda_{\alpha}(x,m,p)
        =
        1+
        \frac{x(1-x)(2p-1)}{2\alpha xp^2+c(x,m,\alpha)}.
\]

\subsection*{\texorpdfstring{Maximizing in \(p\)}{Maximizing in p}}

We want to determine the optimal length of long jobs.
For fixed \(x,m,\alpha\), abbreviate
\[
        c=c(x,m,\alpha).
\]
Then
\[
        \lambda_{\alpha}(x,m,p)
        =
        1+
        \frac{x(1-x)(2p-1)}{2\alpha xp^2+c}.
\]
Equivalently,
\[
        \lambda_{\alpha}(x,m,p)
        =
        1+
        \frac{1-x}{2\alpha}
        \frac{2p-1}{p^2+q},
        \qquad
        q=\frac{c}{2\alpha x}.
\]
Differentiate the \(p\)-dependent factor:
\[
\begin{aligned}
\frac{d}{dp}\left(\frac{2p-1}{p^2+q}\right)
&=
\frac{2(p^2+q)-(2p-1)(2p)}{(p^2+q)^2} \\
&=
\frac{-2p^2+2p+2q}{(p^2+q)^2}.
\end{aligned}
\]
The critical point is therefore determined by
\[
        -p^2+p+q=0.
\]
Since \(p>0\), the maximizing critical point is
\[
        p^*
        =
        \frac{1+\sqrt{1+4q}}{2}
        =
        \frac{1+\sqrt{1+\frac{2c}{\alpha x}}}{2}.
\]
At this point,
\[
        q=(p^*)^2-p^*,
\]
so
\[
\begin{aligned}
\frac{2p^*-1}{(p^*)^2+q}
&=
\frac{2p^*-1}{(p^*)^2+(p^*)^2-p^*}  \\
&=
\frac{2p^*-1}{p^*(2p^*-1)} \\
&=
\frac{1}{p^*}.
\end{aligned}
\]
Hence the optimized value at fixed \(x,m, \alpha\) is
\[
\boxed{
        \lambda^*_{\alpha}(x,m)
        =
        1+
        \frac{1-x}{2\alpha p^*}
        =
        1+
        \frac{1-x}{\alpha\left(1+\sqrt{1+\frac{2c(x,m,\alpha)}{\alpha x}}\right)}
        }.
\]

\subsection*{\texorpdfstring{Maximizing in \(x\) and \(m\)}{Maximizing in x and m}}


We first argue that after optimizing over \(p\), the best value is
obtained at one of the two endpoints of our $\frac{u}{m}$ interval
\[
        x=\frac{1}{k+1}
        \qquad\text{or}\qquad
        x=\frac1k.
\]

We set
\[
        d=m-ku.
\]
The constraint
\[
        \frac{1}{k+1}\leq \frac{u}{m}\leq \frac1k
\]
is equivalent to
\[
        0\leq d\leq u.
\]
So that we can write:
\[
        t=\frac{d}{u},
        \qquad 0\leq t\leq 1.
\]
Since
\[
        m=ku+d=u(k+t),
\]
we get
\[
        x=\frac{u}{m}=\frac{1}{k+t},
        \qquad
        r=1-kx=\frac{t}{k+t}.
\]
Thus, the upper endpoint \(x=1/k\) corresponds to \(t=0\), and the lower
endpoint \(x=1/(k+1)\) corresponds to \(t=1\). Moreover, we will be able to determine the maximum
with respect to $t$, giving us insight into which $\frac{u}{m}$ ratio e.g what ratio between
singletons and machines, is optimal.

Substituting $x=\frac{u}{m}$ and $r=\frac{t}{k+t}$ into \(\rho(x,m)\), we obtain 
\[
\begin{aligned}
\rho(t,u)
&=
\frac{1}{(k+t)^3}
\left[
t^2-\frac{t^3}{3}
 +(k+t-1)^3
 +(k-1)
 +\frac{t}{3u^2}
\right].
\end{aligned}
\]
Define
\[
        s\coloneqq k+t,
\]
and
\[
        A_u(t)
        \coloneqq
        t^2-\frac{t^3}{3}
        +(k+t-1)^3
        +(k-1)
        +\frac{t}{3u^2}.
\]
Then,
\[
        \rho(t,u)=\frac{A_u(t)}{s^3}.
\]
Also
\[
        1-x=\frac{s-1}{s} \quad \text{and}
        \quad
        \frac{c(x,m,\alpha)}{x}
        =
        s-1-(1-2\alpha)\frac{A_u(t)}{s^2}.
\]
Therefore, the optimized value can be written as
\[
        \lambda_{\alpha}^*(t,u)
        \coloneqq
        1+
        \frac{s-1}{
        \alpha s\left(
        1+\sqrt{
        1+\frac{2}{\alpha}
        \left(
        s-1-(1-2\alpha)\frac{A_u(t)}{s^2}
        \right)
        }
        \right)
        }.
\]

For fixed \(t\), the only dependence on \(u\) is through the term
\[
        \frac{t}{3u^2}
\]
inside \(A_u(t)\). Since \(1-2\alpha\geq 0\), increasing \(A_u(t)\) decreases
\[
        s-1-(1-2\alpha)\frac{A_u(t)}{s^2},
\]
and therefore decreases the square root in the denominator of \(\lambda^*(t,u)\).
Thus, \(\lambda^*(t,u)\) is increasing in \(A_u(t)\). Hence, for fixed \(t\),
the largest possible value is obtained by taking the smallest admissible
\(u\) (which will always be one and corresponds to exactly one singleton).

In particular, for every admissible integer pair \((m,u)\),
\[
        \lambda_{\alpha}^*(t,u)\leq \overline \lambda_{\alpha}(t) \coloneqq \lambda_{\alpha}(t,1),
\]
where \(s=k+t\) and
\[
        \overline \lambda_{\alpha}(t)
        =
        1+
        \frac{s-1}{
        \alpha s\left(
        1+\sqrt{
        1+\frac{2}{\alpha}
        \left(
        s-1-(1-2\alpha)\frac{\overline A(t)}{s^2}
        \right)
        }
        \right)
        },
\]
and
\[
        \overline A(t) \coloneqq A_1(t)
        =
        t^2-\frac{t^3}{3}
        +(k+t-1)^3
        +(k-1)
        +\frac{t}{3}.
\]
This upper bound is obtained by taking \(u=1\) (as this is the smallest feasible choice for $u$).

We claim that \(\overline \lambda\) has no interior maximum on \([0,1]\) which we show in the next lemma.

\begin{lemma}
Assume
\[
        k\geq 2,\qquad 0<\alpha\leq \frac12.
\]
Every critical point of \(\overline \lambda_{\alpha}\) in \((0,1)\) is a strict local
minimum.
\end{lemma}

\begin{proof}
First,
\[
        \overline A'(t)
        =
        2t-t^2+3(k+t-1)^2+\frac13.
\]
Let \(t\in(0,1)\) be a critical point of \(\overline \lambda_{\alpha}\). Differentiating twice then substituting
\(\overline \lambda'_{\alpha}=0\) gives
\begin{align}
        \overline \lambda''_{\alpha}(t)
        =
        \frac{
        (1-2\alpha)Q_{k,\alpha}(t)
        }{
        \alpha^2s^3(s-1)
        \left(
        1+\frac{2}{\alpha}
        \left(
        s-1-(1-2\alpha)\frac{\overline A(t)}{s^2}
        \right)
        \right)^{3/2}
        },
    \label{Wolfram5}
\end{align}
where
\begin{align}
\label{Wolfram6}
Q_{k,\alpha}(t)
={}&
\frac{\alpha^2}{3s^2}
\left[
3(k-1)^2+6(k-1)t+2t^2+1
\right]^2                                     \\
&+
\frac{2\alpha}{s}
\left[
3(k-1)^2+6(k-1)t+2t^2+1
\right]  \notag                                    \\
&+
\frac{2\alpha}{s}
(1-t)\left(3(k-1)+2t\right) \notag.
\end{align}
For
\[
        k\geq 2,\qquad 0<\alpha\leq \frac12,\qquad 0<t<1,
\]
the denominator is positive. Indeed, \(\alpha>0\), \(s=k+t>1\), and the
argument in the square root is positive because \(\overline \lambda_{\alpha}(t)\) is
real-valued on our interval. Also \(1-2\alpha\geq 0\), and
every term in \(Q_{k,\alpha}(t)\) is nonnegative. Moreover,
\[
        3(k-1)^2+6(k-1)t+2t^2+1>0,
\]
so \(Q_{k,\alpha}(t)>0\). Hence
\[
        \overline \lambda_{\alpha}''(t)>0.
\]
Thus every interior critical point is a strict local minimum.
\end{proof}

Therefore \(\overline \lambda_{\alpha}\) has no interior local maximum. Since
\(\overline \lambda_{\alpha}\) is continuous on \([0,1]\), its maximum is attained at an
endpoint:
\[
        \overline \lambda_{\alpha}(t)
        \leq
        \max\{\overline \lambda_{\alpha}(0),\overline \lambda_{\alpha}(1)\}
        \qquad (0\leq t\leq 1).
\]
Finally, because \(\lambda_{\alpha}^*(t,u)\leq \overline \lambda_{\alpha}(t)\), the original integer
problem is also maximized at one of the two endpoints,
\[
        t=0
        \qquad\text{or}\qquad
        t=1.
\]

\subsection*{Maximum at the endpoints}

The remaining maximization is over integer pairs \((m,u)\) satisfying
\[
        \frac{1}{k+1}\leq \frac um\leq \frac1k.
\]
The maximum is attained at one of the two boundary ratios:
\[
        \frac um=\frac{1}{k+1}
        \qquad\text{or}\qquad
        \frac um=\frac1k.
\]
The corresponding integer maximizers are
\[
        (m,u)=(k+1,1)
        \qquad\text{and}\qquad
        (m,u)=(k,1).
\]

\subsection*{Lower boundary point}

Take
\[
        m=k+1,\qquad u=1,\qquad x=\frac{1}{k+1}.
\]
And evaluate the formulas at this point. We have
\[
        r=1-kx=\frac{1}{k+1}.
\]
Using the formula for \(\rho\),
\[
\begin{aligned}
\rho_1
&=
xr^2-\frac{r^3}{3}
 +(1-x)^3
 +(k-1)x^3
 +\frac{r}{3m^2} \\
&=
\frac{k^3+k}{(k+1)^3} \\
&=
\frac{k(k^2+1)}{(k+1)^3}.
\end{aligned}
\]
Here, $\rho_1$ denotes the value of \(\rho\) at the lower endpoint \(x=1/(k+1)\).
Thus,
\[
\begin{aligned}
c_1
&=
1-\frac{1}{k+1}
-(1-2\alpha)\frac{k(k^2+1)}{(k+1)^3} \\
&=
\frac{2k\left(k+\alpha(k^2+1)\right)}{(k+1)^3}.
\end{aligned}
\]
Therefore
\[
\frac{2c_1}{\alpha x}
=
\frac{2c_1}{\alpha/(k+1)}
=
\frac{4k\left(k+\alpha(k^2+1)\right)}{\alpha(k+1)^2}.
\]
So the optimizing \(p\) on the lower branch $B^{(1)}_k$ is
\[
\boxed{
p_1
=
\frac{
1+\sqrt{
1+\frac{4k\left(k+\alpha(k^2+1)\right)}{\alpha(k+1)^2}
}
}
{2}.
}
\]
The corresponding value is
\[
\boxed{
B^{(1)}_k(\alpha)
\coloneqq
1+
\frac{k}{
\alpha(k+1)\left(
1+\sqrt{
1+\frac{4k\left(k+\alpha(k^2+1)\right)}{\alpha(k+1)^2}
}
\right)}
.}
\]

\subsection*{Upper boundary point}

Take
\[
        m=k,\qquad u=1,\qquad x=\frac1k.
\]
Then
\[
        r=1-kx=0.
\]
As in the previous case, we can evaluate \(\rho\) at this point:
\[
\begin{aligned}
\rho_2
&=
(1-x)^3+(k-1)x^3 \\
&=
\frac{(k-1)\left((k-1)^2+1\right)}{k^3}.
\end{aligned}
\]
Thus
\[
\begin{aligned}
c_2
&=
1-\frac1k
-(1-2\alpha)
\frac{(k-1)\left((k-1)^2+1\right)}{k^3} \\
&=
\frac{
2(k-1)^2
 +2\alpha(k-1)\left((k-1)^2+1\right)
}{k^3} \\
&=
\frac{
2(k-1)\left((k-1)+\alpha\left((k-1)^2+1\right)\right)
}{k^3}.
\end{aligned}
\]
By observing that,
\[
        (k-1)^2+1=k^2-2k+2, 
\]
We can write
\[
\frac{2c_2}{\alpha x}
=
\frac{2c_2}{\alpha/k}
=
\frac{
4(k-1)\left((k-1)+\alpha(k^2-2k+2)\right)
}{\alpha k^2}.
\]
The optimizing \(p\) on the upper branch is
\[
\boxed{
p_2
=
\frac{
1+\sqrt{
1+
\frac{
4(k-1)\left((k-1)+\alpha(k^2-2k+2)\right)
}{\alpha k^2}
}
}{2}
.}
\]
The corresponding value is
\[
\boxed{
B^{(2)}_k(\alpha)
\coloneqq
1+
\frac{k-1}{
\alpha k\left(
1+\sqrt{
1+
\frac{
4(k-1)\left((k-1)+\alpha(k^2-2k+2)\right)
}{\alpha k^2}
}
\right)}
.}
\]

\subsection*{Comparison of the Cases}

For \(0<\alpha\leq 1/2\), the maximum is
\[
        \max\{B^{(1)}_k(\alpha),B^{(2)}_k(\alpha)\}.
\]
The switching point is obtained by solving
\[
        B^{(1)}_k(\alpha)=B^{(2)}_k(\alpha).
\]
First, by comparing $B^{(1)}_k$ and $B^{(2)}_k$, one realizes that (unsurprisingly) $B^{(2)}_{k+1}(\alpha)=B^{(1)}_k(\alpha)$
so the maximum is
\[
        \max\{B^{(2)}_{k+1}(\alpha),B^{(2)}_{k}(\alpha)\}.
\]
and the switching point can be obtained by solving
\[
        B^{(2)}_{k+1}(\alpha)=B^{(2)}_{k}(\alpha).
\]
This equality reduces to the linear equation
\[
        1-\alpha\left(k^4-2k^3-2k^2+k+3\right)=0.
\]
Since we only work with $B^{(2)}_k$, we define $B_k \coloneqq B^{(2)}_k$.
Therefore, the "switching point" $\alpha_k$, (at which $B_k(\alpha_k) =B_{k+1}(\alpha_k)$) is
\[
        \alpha_k \coloneqq
        \frac{1}{k^4-2k^3-2k^2+k+3}.
\]

As in the determination of $B_k$, $k$ was equal to $m$, we can use the machine number in our formulas.
giving the switching point as

\[
\boxed{
        \alpha_m=
        \frac{1}{m^4-2m^3-2m^2+m+3}
        .}
\]

For \(m=2\), the denominator is
\[
        2^4-2\cdot 2^3-2\cdot 2^2+2+3
        =
        -3,
\]
so there is no positive switching point. This gives insight
into the fact that the set of instances restricted to $m=2$ machines
is non-dominant for every $\alpha \in \Big(0,\frac12\Big)$.
In this case \(B_3(\alpha)>B_2(\alpha)\) for every
\[
        0<\alpha\leq \frac12.
\]

For \(m\geq 3\), the denominator of $\alpha_m$ is positive and gives a genuine switching point. Thus
\[
        B_{m+1}(\alpha)>B_m(\alpha)
        \quad\Longleftrightarrow\quad
        0<\alpha<\alpha_m,
\]
and
\[
        B_m(\alpha)>B_{m+1}(\alpha)
        \quad\Longleftrightarrow\quad
        \alpha_m<\alpha\leq \frac12.
\]
At \(\alpha=\alpha_m\), both branches give the same value.

\subsection*{Maximum}\label{subsec: max Lbound}
Summarizing,

\[
\boxed{
\lambda_{\alpha} \ge
\begin{cases}
B_3(\alpha),
& \displaystyle \alpha_3\leq\alpha\leq\frac12,\\[6pt]
B_{m+1}(\alpha),
& \alpha_{m+1}\leq\alpha \le \alpha_{m},\quad m\geq 3.
\end{cases}
}
\]
\end{proof}
\backmatter{}
\listoffigures
\printbibliography[heading=bibintoc]
\end{document}

%%% Local Variables:
%%% mode: latex
%%% TeX-engine: default
%%% TeX-command-extra-options: "-shell-escape"
%%% ispell-local-dictionary: "american"
%%% eval: (setenv "TEXINPUTS" ".//:")
%%% TeX-master: t
%%% End:
